% Sprawozdanie: ticketbench, porównanie PostgreSQL, MariaDB, MongoDB i Valkey.
% Kompilacja: tectonic -X compile sprawozdanie.tex  (albo latexmk -xelatex)
% Tabele i dane wykresów w generated/ tworzy `tbench report`, nie edytować ręcznie.
\documentclass[11pt,a4paper]{article}

\usepackage{fontspec}
\usepackage[polish]{babel}
\usepackage[a4paper,margin=2.3cm]{geometry}
\usepackage{microtype}
\usepackage{amssymb}
\usepackage{booktabs}
\usepackage{graphicx}
\usepackage[table]{xcolor}
\usepackage{array}
\usepackage{float}
\usepackage{caption}
\usepackage{subcaption}
\usepackage{enumitem}
\usepackage{fancyvrb}
\usepackage{fvextra}
\usepackage{listings}
\usepackage{siunitx}
\usepackage{xfp}
\usepackage{pgfplots}
\usepgfplotslibrary{groupplots}
\usetikzlibrary{arrows.meta,positioning,fit,backgrounds}
\usepackage[hidelinks]{hyperref}

\pgfplotsset{compat=1.18, /pgf/number format/.cd, use comma, 1000 sep={\,}}
\sisetup{output-decimal-marker={,}, group-separator={\,}, group-minimum-digits=4}
\setlist{itemsep=2pt,topsep=4pt}
\setlength{\emergencystretch}{3em}
% Long inline code may break at spaces and after underscores.
\let\origtexttt\texttt
\renewcommand{\texttt}[1]{\origtexttt{\def\_{\textunderscore\allowbreak}#1}}
\lstset{basicstyle=\ttfamily\footnotesize,breaklines=true,columns=fullflexible,frame=single,framesep=3pt}
\captionsetup{font=small,labelfont=bf}

\definecolor{cpg}{HTML}{2F6DB5}
\definecolor{cmaria}{HTML}{C0392B}
\definecolor{cmongo}{HTML}{2E8B57}
\definecolor{cvalkey}{HTML}{8E44AD}

\pgfplotsset{
  tb/.style={
    width=\linewidth, height=6.2cm, grid=major, grid style={gray!25},
    legend style={font=\footnotesize, at={(0.5,1.02)}, anchor=south, legend columns=4, draw=none},
    tick label style={font=\footnotesize}, label style={font=\footnotesize},
    unbounded coords=jump, every axis plot/.append style={mark options={scale=0.8}},
  },
  pg/.style={cpg, thick, mark=*},
  maria/.style={cmaria, thick, mark=square*},
  mongo/.style={cmongo, thick, mark=triangle*},
  valkey/.style={cvalkey, thick, mark=diamond*},
}

% generated by tbench report
\newcommand{\HOneCases}{12}
\newcommand{\HOneGrows}{12}
\newcommand{\HOneVerdict}{potwierdzona}
\newcommand{\HTwoCases}{54}
\newcommand{\HTwoGrows}{36}
\newcommand{\HTwoVerdict}{potwierdzona częściowo}
\newcommand{\SizesList}{500 tys., 1 mln, 10 mln}
\newcommand{\ScenarioRuns}{5424}
\newcommand{\ContentionRuns}{468}
\newcommand{\ContentionInconsistent}{0}

% generated by tbench report
\makeatletter
\@namedef{tb@s@mariadb@1000000@off@C1}{0.3031}
\@namedef{tb@s@mariadb@1000000@off@C2}{0.3116}
\@namedef{tb@s@mariadb@1000000@off@C3}{28.225}
\@namedef{tb@s@mariadb@1000000@off@C4}{0.156}
\@namedef{tb@s@mariadb@1000000@off@C5}{0.2962}
\@namedef{tb@s@mariadb@1000000@off@C6}{0.2055}
\@namedef{tb@s@mariadb@1000000@off@D1}{180.2841}
\@namedef{tb@s@mariadb@1000000@off@D2}{178.4935}
\@namedef{tb@s@mariadb@1000000@off@D3}{28.9419}
\@namedef{tb@s@mariadb@1000000@off@D4}{320.6309}
\@namedef{tb@s@mariadb@1000000@off@D5}{162.8766}
\@namedef{tb@s@mariadb@1000000@off@D6}{293.6292}
\@namedef{tb@s@mariadb@1000000@off@J1}{0.3471}
\@namedef{tb@s@mariadb@1000000@off@J2}{22.3513}
\@namedef{tb@s@mariadb@1000000@off@J3}{0.2242}
\@namedef{tb@s@mariadb@1000000@off@R1}{88.2863}
\@namedef{tb@s@mariadb@1000000@off@R2}{84.4595}
\@namedef{tb@s@mariadb@1000000@off@R3}{0.1793}
\@namedef{tb@s@mariadb@1000000@off@R4}{74.9006}
\@namedef{tb@s@mariadb@1000000@off@R5}{342.9176}
\@namedef{tb@s@mariadb@1000000@off@R6}{90.154}
\@namedef{tb@s@mariadb@1000000@off@U1}{0.1491}
\@namedef{tb@s@mariadb@1000000@off@U2}{0.1642}
\@namedef{tb@s@mariadb@1000000@off@U3}{167.279}
\@namedef{tb@s@mariadb@1000000@off@U4}{0.2085}
\@namedef{tb@s@mariadb@1000000@off@U5}{0.179}
\@namedef{tb@s@mariadb@1000000@off@U6}{260.4769}
\@namedef{tb@s@mariadb@1000000@on@C1}{0.3416}
\@namedef{tb@s@mariadb@1000000@on@C2}{0.3367}
\@namedef{tb@s@mariadb@1000000@on@C3}{41.8109}
\@namedef{tb@s@mariadb@1000000@on@C4}{0.2503}
\@namedef{tb@s@mariadb@1000000@on@C5}{0.4899}
\@namedef{tb@s@mariadb@1000000@on@C6}{0.2758}
\@namedef{tb@s@mariadb@1000000@on@D1}{0.4127}
\@namedef{tb@s@mariadb@1000000@on@D2}{0.4223}
\@namedef{tb@s@mariadb@1000000@on@D3}{2.7065}
\@namedef{tb@s@mariadb@1000000@on@D4}{0.6222}
\@namedef{tb@s@mariadb@1000000@on@D5}{5.2555}
\@namedef{tb@s@mariadb@1000000@on@D6}{9.7439}
\@namedef{tb@s@mariadb@1000000@on@J1}{0.0913}
\@namedef{tb@s@mariadb@1000000@on@J2}{13.5085}
\@namedef{tb@s@mariadb@1000000@on@J3}{0.2005}
\@namedef{tb@s@mariadb@1000000@on@R1}{0.1231}
\@namedef{tb@s@mariadb@1000000@on@R2}{1.3235}
\@namedef{tb@s@mariadb@1000000@on@R3}{0.1336}
\@namedef{tb@s@mariadb@1000000@on@R4}{0.0938}
\@namedef{tb@s@mariadb@1000000@on@R5}{18.7137}
\@namedef{tb@s@mariadb@1000000@on@R6}{74.6524}
\@namedef{tb@s@mariadb@1000000@on@U1}{0.5166}
\@namedef{tb@s@mariadb@1000000@on@U2}{0.1608}
\@namedef{tb@s@mariadb@1000000@on@U3}{0.9025}
\@namedef{tb@s@mariadb@1000000@on@U4}{0.1744}
\@namedef{tb@s@mariadb@1000000@on@U5}{0.1675}
\@namedef{tb@s@mariadb@1000000@on@U6}{7.9297}
\@namedef{tb@s@mongo@1000000@off@C1}{1.0271}
\@namedef{tb@s@mongo@1000000@off@C2}{0.7619}
\@namedef{tb@s@mongo@1000000@off@C3}{43.9187}
\@namedef{tb@s@mongo@1000000@off@C4}{0.5981}
\@namedef{tb@s@mongo@1000000@off@C5}{0.4659}
\@namedef{tb@s@mongo@1000000@off@C6}{1.1158}
\@namedef{tb@s@mongo@1000000@off@D1}{0.7955}
\@namedef{tb@s@mongo@1000000@off@D2}{0.9491}
\@namedef{tb@s@mongo@1000000@off@D3}{43.919}
\@namedef{tb@s@mongo@1000000@off@D4}{84.2524}
\@namedef{tb@s@mongo@1000000@off@D5}{240.4191}
\@namedef{tb@s@mongo@1000000@off@D6}{297.5274}
\@namedef{tb@s@mongo@1000000@off@J1}{0.2739}
\@namedef{tb@s@mongo@1000000@off@J2}{30.6218}
\@namedef{tb@s@mongo@1000000@off@J3}{0.8391}
\@namedef{tb@s@mongo@1000000@off@R1}{0.1138}
\@namedef{tb@s@mongo@1000000@off@R2}{282.7316}
\@namedef{tb@s@mongo@1000000@off@R3}{0.6306}
\@namedef{tb@s@mongo@1000000@off@R4}{42.304}
\@namedef{tb@s@mongo@1000000@off@R5}{202.2441}
\@namedef{tb@s@mongo@1000000@off@R6}{204.4975}
\@namedef{tb@s@mongo@1000000@off@U1}{1.5906}
\@namedef{tb@s@mongo@1000000@off@U2}{1.3784}
\@namedef{tb@s@mongo@1000000@off@U3}{252.3331}
\@namedef{tb@s@mongo@1000000@off@U4}{0.8658}
\@namedef{tb@s@mongo@1000000@off@U5}{0.5308}
\@namedef{tb@s@mongo@1000000@off@U6}{55.4355}
\@namedef{tb@s@mongo@1000000@on@C1}{1.5727}
\@namedef{tb@s@mongo@1000000@on@C2}{0.8448}
\@namedef{tb@s@mongo@1000000@on@C3}{56.9408}
\@namedef{tb@s@mongo@1000000@on@C4}{0.4711}
\@namedef{tb@s@mongo@1000000@on@C5}{0.3155}
\@namedef{tb@s@mongo@1000000@on@C6}{0.9281}
\@namedef{tb@s@mongo@1000000@on@D1}{1.312}
\@namedef{tb@s@mongo@1000000@on@D2}{0.9543}
\@namedef{tb@s@mongo@1000000@on@D3}{6.1309}
\@namedef{tb@s@mongo@1000000@on@D4}{1.9892}
\@namedef{tb@s@mongo@1000000@on@D5}{11.1225}
\@namedef{tb@s@mongo@1000000@on@D6}{15.7368}
\@namedef{tb@s@mongo@1000000@on@J1}{0.134}
\@namedef{tb@s@mongo@1000000@on@J2}{10.0286}
\@namedef{tb@s@mongo@1000000@on@J3}{0.9607}
\@namedef{tb@s@mongo@1000000@on@R1}{0.0988}
\@namedef{tb@s@mongo@1000000@on@R2}{1.0662}
\@namedef{tb@s@mongo@1000000@on@R3}{0.1238}
\@namedef{tb@s@mongo@1000000@on@R4}{0.1285}
\@namedef{tb@s@mongo@1000000@on@R5}{10.0956}
\@namedef{tb@s@mongo@1000000@on@R6}{213.5855}
\@namedef{tb@s@mongo@1000000@on@U1}{0.5462}
\@namedef{tb@s@mongo@1000000@on@U2}{0.7075}
\@namedef{tb@s@mongo@1000000@on@U3}{2.3558}
\@namedef{tb@s@mongo@1000000@on@U4}{0.4691}
\@namedef{tb@s@mongo@1000000@on@U5}{0.8814}
\@namedef{tb@s@mongo@1000000@on@U6}{19.4717}
\@namedef{tb@s@postgres@1000000@off@C1}{0.37}
\@namedef{tb@s@postgres@1000000@off@C2}{0.3928}
\@namedef{tb@s@postgres@1000000@off@C3}{9.5662}
\@namedef{tb@s@postgres@1000000@off@C4}{0.2045}
\@namedef{tb@s@postgres@1000000@off@C5}{0.3913}
\@namedef{tb@s@postgres@1000000@off@C6}{0.3164}
\@namedef{tb@s@postgres@1000000@off@D1}{36.2509}
\@namedef{tb@s@postgres@1000000@off@D2}{32.5838}
\@namedef{tb@s@postgres@1000000@off@D3}{4.2557}
\@namedef{tb@s@postgres@1000000@off@D4}{51.0204}
\@namedef{tb@s@postgres@1000000@off@D5}{23.6169}
\@namedef{tb@s@postgres@1000000@off@D6}{52.0704}
\@namedef{tb@s@postgres@1000000@off@J1}{0.145}
\@namedef{tb@s@postgres@1000000@off@J2}{4.0175}
\@namedef{tb@s@postgres@1000000@off@J3}{0.2249}
\@namedef{tb@s@postgres@1000000@off@R1}{15.2427}
\@namedef{tb@s@postgres@1000000@off@R2}{14.8504}
\@namedef{tb@s@postgres@1000000@off@R3}{0.1398}
\@namedef{tb@s@postgres@1000000@off@R4}{12.7874}
\@namedef{tb@s@postgres@1000000@off@R5}{42.1449}
\@namedef{tb@s@postgres@1000000@off@R6}{22.5364}
\@namedef{tb@s@postgres@1000000@off@U1}{0.2841}
\@namedef{tb@s@postgres@1000000@off@U2}{0.3426}
\@namedef{tb@s@postgres@1000000@off@U3}{28.3668}
\@namedef{tb@s@postgres@1000000@off@U4}{0.2176}
\@namedef{tb@s@postgres@1000000@off@U5}{0.202}
\@namedef{tb@s@postgres@1000000@off@U6}{48.1667}
\@namedef{tb@s@postgres@1000000@on@C1}{0.3304}
\@namedef{tb@s@postgres@1000000@on@C2}{0.3007}
\@namedef{tb@s@postgres@1000000@on@C3}{18.1528}
\@namedef{tb@s@postgres@1000000@on@C4}{0.1656}
\@namedef{tb@s@postgres@1000000@on@C5}{0.3163}
\@namedef{tb@s@postgres@1000000@on@C6}{0.229}
\@namedef{tb@s@postgres@1000000@on@D1}{0.3368}
\@namedef{tb@s@postgres@1000000@on@D2}{0.2931}
\@namedef{tb@s@postgres@1000000@on@D3}{0.3558}
\@namedef{tb@s@postgres@1000000@on@D4}{0.5953}
\@namedef{tb@s@postgres@1000000@on@D5}{0.658}
\@namedef{tb@s@postgres@1000000@on@D6}{1.7704}
\@namedef{tb@s@postgres@1000000@on@J1}{0.0631}
\@namedef{tb@s@postgres@1000000@on@J2}{1.9005}
\@namedef{tb@s@postgres@1000000@on@J3}{0.1719}
\@namedef{tb@s@postgres@1000000@on@R1}{0.0687}
\@namedef{tb@s@postgres@1000000@on@R2}{0.711}
\@namedef{tb@s@postgres@1000000@on@R3}{0.0721}
\@namedef{tb@s@postgres@1000000@on@R4}{0.0619}
\@namedef{tb@s@postgres@1000000@on@R5}{2.239}
\@namedef{tb@s@postgres@1000000@on@R6}{11.5752}
\@namedef{tb@s@postgres@1000000@on@U1}{0.1902}
\@namedef{tb@s@postgres@1000000@on@U2}{0.1671}
\@namedef{tb@s@postgres@1000000@on@U3}{0.8806}
\@namedef{tb@s@postgres@1000000@on@U4}{0.2035}
\@namedef{tb@s@postgres@1000000@on@U5}{0.1757}
\@namedef{tb@s@postgres@1000000@on@U6}{9.8058}
\@namedef{tb@s@valkey@1000000@off@C1}{0.2325}
\@namedef{tb@s@valkey@1000000@off@C2}{0.2432}
\@namedef{tb@s@valkey@1000000@off@C3}{11.8325}
\@namedef{tb@s@valkey@1000000@off@C4}{0.3186}
\@namedef{tb@s@valkey@1000000@off@C5}{0.3472}
\@namedef{tb@s@valkey@1000000@off@C6}{0.2452}
\@namedef{tb@s@valkey@1000000@off@D1}{0.7524}
\@namedef{tb@s@valkey@1000000@off@D2}{0.4285}
\@namedef{tb@s@valkey@1000000@off@D3}{1123.4272}
\@namedef{tb@s@valkey@1000000@off@D4}{961.2958}
\@namedef{tb@s@valkey@1000000@off@D5}{0.1117}
\@namedef{tb@s@valkey@1000000@off@D6}{1162.8407}
\@namedef{tb@s@valkey@1000000@off@J1}{0.6143}
\@namedef{tb@s@valkey@1000000@off@J2}{1105.8321}
\@namedef{tb@s@valkey@1000000@off@J3}{0.1078}
\@namedef{tb@s@valkey@1000000@off@R1}{0.161}
\@namedef{tb@s@valkey@1000000@off@R2}{1.7669}
\@namedef{tb@s@valkey@1000000@off@R3}{0.5931}
\@namedef{tb@s@valkey@1000000@off@R4}{938.2485}
\@namedef{tb@s@valkey@1000000@off@R5}{0.6217}
\@namedef{tb@s@valkey@1000000@off@R6}{0.1426}
\@namedef{tb@s@valkey@1000000@off@U1}{0.1592}
\@namedef{tb@s@valkey@1000000@off@U2}{0.1052}
\@namedef{tb@s@valkey@1000000@off@U3}{2.1099}
\@namedef{tb@s@valkey@1000000@off@U4}{0.188}
\@namedef{tb@s@valkey@1000000@off@U5}{0.1332}
\@namedef{tb@s@valkey@1000000@off@U6}{967.015}
\@namedef{tb@s@valkey@1000000@on@C1}{0.121}
\@namedef{tb@s@valkey@1000000@on@C2}{0.1261}
\@namedef{tb@s@valkey@1000000@on@C3}{8.3585}
\@namedef{tb@s@valkey@1000000@on@C4}{0.1133}
\@namedef{tb@s@valkey@1000000@on@C5}{0.2065}
\@namedef{tb@s@valkey@1000000@on@C6}{0.1012}
\@namedef{tb@s@valkey@1000000@on@D1}{0.185}
\@namedef{tb@s@valkey@1000000@on@D2}{0.1477}
\@namedef{tb@s@valkey@1000000@on@D3}{6.3495}
\@namedef{tb@s@valkey@1000000@on@D4}{0.3045}
\@namedef{tb@s@valkey@1000000@on@D5}{0.1032}
\@namedef{tb@s@valkey@1000000@on@D6}{20.1048}
\@namedef{tb@s@valkey@1000000@on@J1}{0.1908}
\@namedef{tb@s@valkey@1000000@on@J2}{44.0838}
\@namedef{tb@s@valkey@1000000@on@J3}{0.0927}
\@namedef{tb@s@valkey@1000000@on@R1}{0.1019}
\@namedef{tb@s@valkey@1000000@on@R2}{1.5657}
\@namedef{tb@s@valkey@1000000@on@R3}{0.1585}
\@namedef{tb@s@valkey@1000000@on@R4}{0.1258}
\@namedef{tb@s@valkey@1000000@on@R5}{0.2313}
\@namedef{tb@s@valkey@1000000@on@R6}{0.1197}
\@namedef{tb@s@valkey@1000000@on@U1}{0.104}
\@namedef{tb@s@valkey@1000000@on@U2}{0.1142}
\@namedef{tb@s@valkey@1000000@on@U3}{2.0101}
\@namedef{tb@s@valkey@1000000@on@U4}{0.1359}
\@namedef{tb@s@valkey@1000000@on@U5}{0.1191}
\@namedef{tb@s@valkey@1000000@on@U6}{17.8309}
\@namedef{tb@s@mariadb@10000000@off@C1}{0.3664}
\@namedef{tb@s@mariadb@10000000@off@C2}{0.4208}
\@namedef{tb@s@mariadb@10000000@off@C3}{29.4855}
\@namedef{tb@s@mariadb@10000000@off@C4}{0.1496}
\@namedef{tb@s@mariadb@10000000@off@C5}{0.2553}
\@namedef{tb@s@mariadb@10000000@off@C6}{0.1946}
\@namedef{tb@s@mariadb@10000000@off@D1}{3305.3433}
\@namedef{tb@s@mariadb@10000000@off@D2}{3302.2437}
\@namedef{tb@s@mariadb@10000000@off@D3}{252.6478}
\@namedef{tb@s@mariadb@10000000@off@D4}{4789.83}
\@namedef{tb@s@mariadb@10000000@off@D5}{3101.511}
\@namedef{tb@s@mariadb@10000000@off@D6}{4703.2798}
\@namedef{tb@s@mariadb@10000000@off@J1}{2.2304}
\@namedef{tb@s@mariadb@10000000@off@J2}{221.3637}
\@namedef{tb@s@mariadb@10000000@off@J3}{0.2412}
\@namedef{tb@s@mariadb@10000000@off@R1}{2196.9663}
\@namedef{tb@s@mariadb@10000000@off@R2}{2237.3024}
\@namedef{tb@s@mariadb@10000000@off@R3}{0.6326}
\@namedef{tb@s@mariadb@10000000@off@R4}{1465.1823}
\@namedef{tb@s@mariadb@10000000@off@R5}{4214.5175}
\@namedef{tb@s@mariadb@10000000@off@R6}{2333.1199}
\@namedef{tb@s@mariadb@10000000@off@U1}{0.1195}
\@namedef{tb@s@mariadb@10000000@off@U2}{0.1982}
\@namedef{tb@s@mariadb@10000000@off@U3}{3104.5163}
\@namedef{tb@s@mariadb@10000000@off@U4}{0.2954}
\@namedef{tb@s@mariadb@10000000@off@U5}{0.1644}
\@namedef{tb@s@mariadb@10000000@off@U6}{4060.9173}
\@namedef{tb@s@mariadb@10000000@on@C1}{1.4003}
\@namedef{tb@s@mariadb@10000000@on@C2}{0.7544}
\@namedef{tb@s@mariadb@10000000@on@C3}{48.3965}
\@namedef{tb@s@mariadb@10000000@on@C4}{0.2201}
\@namedef{tb@s@mariadb@10000000@on@C5}{0.3888}
\@namedef{tb@s@mariadb@10000000@on@C6}{0.3101}
\@namedef{tb@s@mariadb@10000000@on@D1}{1.1298}
\@namedef{tb@s@mariadb@10000000@on@D2}{0.5515}
\@namedef{tb@s@mariadb@10000000@on@D3}{4.0194}
\@namedef{tb@s@mariadb@10000000@on@D4}{1.5883}
\@namedef{tb@s@mariadb@10000000@on@D5}{7.0249}
\@namedef{tb@s@mariadb@10000000@on@D6}{12.5376}
\@namedef{tb@s@mariadb@10000000@on@J1}{0.3998}
\@namedef{tb@s@mariadb@10000000@on@J2}{159.7937}
\@namedef{tb@s@mariadb@10000000@on@J3}{0.2504}
\@namedef{tb@s@mariadb@10000000@on@R1}{1.0988}
\@namedef{tb@s@mariadb@10000000@on@R2}{2.5563}
\@namedef{tb@s@mariadb@10000000@on@R3}{0.1971}
\@namedef{tb@s@mariadb@10000000@on@R4}{0.1039}
\@namedef{tb@s@mariadb@10000000@on@R5}{58.7821}
\@namedef{tb@s@mariadb@10000000@on@R6}{699.898}
\@namedef{tb@s@mariadb@10000000@on@U1}{0.2619}
\@namedef{tb@s@mariadb@10000000@on@U2}{0.3735}
\@namedef{tb@s@mariadb@10000000@on@U3}{2.9048}
\@namedef{tb@s@mariadb@10000000@on@U4}{0.3578}
\@namedef{tb@s@mariadb@10000000@on@U5}{0.28}
\@namedef{tb@s@mariadb@10000000@on@U6}{10.7561}
\@namedef{tb@s@mongo@10000000@off@C1}{1.2068}
\@namedef{tb@s@mongo@10000000@off@C2}{1.3502}
\@namedef{tb@s@mongo@10000000@off@C3}{43.0979}
\@namedef{tb@s@mongo@10000000@off@C4}{0.584}
\@namedef{tb@s@mongo@10000000@off@C5}{0.3523}
\@namedef{tb@s@mongo@10000000@off@C6}{0.9411}
\@namedef{tb@s@mongo@10000000@off@D1}{0.9566}
\@namedef{tb@s@mongo@10000000@off@D2}{1.1295}
\@namedef{tb@s@mongo@10000000@off@D3}{533.5324}
\@namedef{tb@s@mongo@10000000@off@D4}{1036.7572}
\@namedef{tb@s@mongo@10000000@off@D5}{2695.7061}
\@namedef{tb@s@mongo@10000000@off@D6}{3306.3932}
\@namedef{tb@s@mongo@10000000@off@J1}{1.2181}
\@namedef{tb@s@mongo@10000000@off@J2}{432.5127}
\@namedef{tb@s@mongo@10000000@off@J3}{0.935}
\@namedef{tb@s@mongo@10000000@off@R1}{0.1268}
\@namedef{tb@s@mongo@10000000@off@R2}{3066.9223}
\@namedef{tb@s@mongo@10000000@off@R3}{1.3072}
\@namedef{tb@s@mongo@10000000@off@R4}{510.0516}
\@namedef{tb@s@mongo@10000000@off@R5}{2401.6889}
\@namedef{tb@s@mongo@10000000@off@R6}{2375.3541}
\@namedef{tb@s@mongo@10000000@off@U1}{0.1934}
\@namedef{tb@s@mongo@10000000@off@U2}{1.2939}
\@namedef{tb@s@mongo@10000000@off@U3}{2632.936}
\@namedef{tb@s@mongo@10000000@off@U4}{0.8764}
\@namedef{tb@s@mongo@10000000@off@U5}{0.7355}
\@namedef{tb@s@mongo@10000000@off@U6}{484.6253}
\@namedef{tb@s@mongo@10000000@on@C1}{1.4903}
\@namedef{tb@s@mongo@10000000@on@C2}{0.9189}
\@namedef{tb@s@mongo@10000000@on@C3}{60.9042}
\@namedef{tb@s@mongo@10000000@on@C4}{0.4043}
\@namedef{tb@s@mongo@10000000@on@C5}{0.6075}
\@namedef{tb@s@mongo@10000000@on@C6}{1.129}
\@namedef{tb@s@mongo@10000000@on@D1}{1.7197}
\@namedef{tb@s@mongo@10000000@on@D2}{1.3598}
\@namedef{tb@s@mongo@10000000@on@D3}{6.6613}
\@namedef{tb@s@mongo@10000000@on@D4}{1.9931}
\@namedef{tb@s@mongo@10000000@on@D5}{11.7421}
\@namedef{tb@s@mongo@10000000@on@D6}{19.0091}
\@namedef{tb@s@mongo@10000000@on@J1}{0.3289}
\@namedef{tb@s@mongo@10000000@on@J2}{446.6642}
\@namedef{tb@s@mongo@10000000@on@J3}{1.1546}
\@namedef{tb@s@mongo@10000000@on@R1}{0.1363}
\@namedef{tb@s@mongo@10000000@on@R2}{1.413}
\@namedef{tb@s@mongo@10000000@on@R3}{0.2534}
\@namedef{tb@s@mongo@10000000@on@R4}{0.1285}
\@namedef{tb@s@mongo@10000000@on@R5}{16.2019}
\@namedef{tb@s@mongo@10000000@on@R6}{2370.1445}
\@namedef{tb@s@mongo@10000000@on@U1}{1.3326}
\@namedef{tb@s@mongo@10000000@on@U2}{0.5809}
\@namedef{tb@s@mongo@10000000@on@U3}{2.4303}
\@namedef{tb@s@mongo@10000000@on@U4}{0.6425}
\@namedef{tb@s@mongo@10000000@on@U5}{0.7159}
\@namedef{tb@s@mongo@10000000@on@U6}{19.5667}
\@namedef{tb@s@postgres@10000000@off@C1}{0.3175}
\@namedef{tb@s@postgres@10000000@off@C2}{0.5974}
\@namedef{tb@s@postgres@10000000@off@C3}{12.6475}
\@namedef{tb@s@postgres@10000000@off@C4}{0.1891}
\@namedef{tb@s@postgres@10000000@off@C5}{0.3318}
\@namedef{tb@s@postgres@10000000@off@C6}{0.2682}
\@namedef{tb@s@postgres@10000000@off@D1}{383.7764}
\@namedef{tb@s@postgres@10000000@off@D2}{496.0865}
\@namedef{tb@s@postgres@10000000@off@D3}{51.2733}
\@namedef{tb@s@postgres@10000000@off@D4}{481.1605}
\@namedef{tb@s@postgres@10000000@off@D5}{363.1775}
\@namedef{tb@s@postgres@10000000@off@D6}{571.8316}
\@namedef{tb@s@postgres@10000000@off@J1}{0.4927}
\@namedef{tb@s@postgres@10000000@off@J2}{23.405}
\@namedef{tb@s@postgres@10000000@off@J3}{0.3053}
\@namedef{tb@s@postgres@10000000@off@R1}{138.36}
\@namedef{tb@s@postgres@10000000@off@R2}{133.6604}
\@namedef{tb@s@postgres@10000000@off@R3}{0.3016}
\@namedef{tb@s@postgres@10000000@off@R4}{118.1356}
\@namedef{tb@s@postgres@10000000@off@R5}{469.3586}
\@namedef{tb@s@postgres@10000000@off@R6}{234.361}
\@namedef{tb@s@postgres@10000000@off@U1}{0.576}
\@namedef{tb@s@postgres@10000000@off@U2}{0.2413}
\@namedef{tb@s@postgres@10000000@off@U3}{350.0178}
\@namedef{tb@s@postgres@10000000@off@U4}{0.5339}
\@namedef{tb@s@postgres@10000000@off@U5}{0.1901}
\@namedef{tb@s@postgres@10000000@off@U6}{565.1649}
\@namedef{tb@s@postgres@10000000@on@C1}{0.3427}
\@namedef{tb@s@postgres@10000000@on@C2}{0.3219}
\@namedef{tb@s@postgres@10000000@on@C3}{22.1191}
\@namedef{tb@s@postgres@10000000@on@C4}{0.1778}
\@namedef{tb@s@postgres@10000000@on@C5}{0.3011}
\@namedef{tb@s@postgres@10000000@on@C6}{0.2003}
\@namedef{tb@s@postgres@10000000@on@D1}{0.5788}
\@namedef{tb@s@postgres@10000000@on@D2}{0.5945}
\@namedef{tb@s@postgres@10000000@on@D3}{0.3792}
\@namedef{tb@s@postgres@10000000@on@D4}{26.8282}
\@namedef{tb@s@postgres@10000000@on@D5}{0.7382}
\@namedef{tb@s@postgres@10000000@on@D6}{1.9772}
\@namedef{tb@s@postgres@10000000@on@J1}{0.0805}
\@namedef{tb@s@postgres@10000000@on@J2}{15.7561}
\@namedef{tb@s@postgres@10000000@on@J3}{0.2562}
\@namedef{tb@s@postgres@10000000@on@R1}{0.357}
\@namedef{tb@s@postgres@10000000@on@R2}{1.4066}
\@namedef{tb@s@postgres@10000000@on@R3}{0.0974}
\@namedef{tb@s@postgres@10000000@on@R4}{0.0803}
\@namedef{tb@s@postgres@10000000@on@R5}{6.6622}
\@namedef{tb@s@postgres@10000000@on@R6}{100.4583}
\@namedef{tb@s@postgres@10000000@on@U1}{0.5749}
\@namedef{tb@s@postgres@10000000@on@U2}{0.1716}
\@namedef{tb@s@postgres@10000000@on@U3}{1.6892}
\@namedef{tb@s@postgres@10000000@on@U4}{0.1898}
\@namedef{tb@s@postgres@10000000@on@U5}{0.4216}
\@namedef{tb@s@postgres@10000000@on@U6}{6.8411}
\@namedef{tb@s@valkey@10000000@off@C1}{0.1317}
\@namedef{tb@s@valkey@10000000@off@C2}{0.1904}
\@namedef{tb@s@valkey@10000000@off@C3}{7.1761}
\@namedef{tb@s@valkey@10000000@off@C4}{0.1401}
\@namedef{tb@s@valkey@10000000@off@C5}{0.2235}
\@namedef{tb@s@valkey@10000000@off@C6}{0.1173}
\@namedef{tb@s@valkey@10000000@off@D1}{0.2018}
\@namedef{tb@s@valkey@10000000@off@D2}{0.1969}
\@namedef{tb@s@valkey@10000000@off@D3}{18353.7742}
\@namedef{tb@s@valkey@10000000@off@D4}{9946.718}
\@namedef{tb@s@valkey@10000000@off@D5}{0.1172}
\@namedef{tb@s@valkey@10000000@off@D6}{9749.0215}
\@namedef{tb@s@valkey@10000000@off@J1}{3.9978}
\@namedef{tb@s@valkey@10000000@off@J2}{9761.1176}
\@namedef{tb@s@valkey@10000000@off@J3}{0.1141}
\@namedef{tb@s@valkey@10000000@off@R1}{0.0902}
\@namedef{tb@s@valkey@10000000@off@R2}{1.5699}
\@namedef{tb@s@valkey@10000000@off@R3}{4.6918}
\@namedef{tb@s@valkey@10000000@off@R4}{9858.1627}
\@namedef{tb@s@valkey@10000000@off@R5}{4.0665}
\@namedef{tb@s@valkey@10000000@off@R6}{0.1557}
\@namedef{tb@s@valkey@10000000@off@U1}{0.1303}
\@namedef{tb@s@valkey@10000000@off@U2}{0.1154}
\@namedef{tb@s@valkey@10000000@off@U3}{2.9033}
\@namedef{tb@s@valkey@10000000@off@U4}{0.1698}
\@namedef{tb@s@valkey@10000000@off@U5}{0.1339}
\@namedef{tb@s@valkey@10000000@off@U6}{9882.99}
\@namedef{tb@s@valkey@10000000@on@C1}{0.1197}
\@namedef{tb@s@valkey@10000000@on@C2}{0.1343}
\@namedef{tb@s@valkey@10000000@on@C3}{6.7639}
\@namedef{tb@s@valkey@10000000@on@C4}{0.0913}
\@namedef{tb@s@valkey@10000000@on@C5}{0.177}
\@namedef{tb@s@valkey@10000000@on@C6}{0.0941}
\@namedef{tb@s@valkey@10000000@on@D1}{0.1761}
\@namedef{tb@s@valkey@10000000@on@D2}{0.1444}
\@namedef{tb@s@valkey@10000000@on@D3}{6.6309}
\@namedef{tb@s@valkey@10000000@on@D4}{0.4446}
\@namedef{tb@s@valkey@10000000@on@D5}{0.106}
\@namedef{tb@s@valkey@10000000@on@D6}{17.067}
\@namedef{tb@s@valkey@10000000@on@J1}{0.5634}
\@namedef{tb@s@valkey@10000000@on@J2}{444.1919}
\@namedef{tb@s@valkey@10000000@on@J3}{0.1157}
\@namedef{tb@s@valkey@10000000@on@R1}{0.0918}
\@namedef{tb@s@valkey@10000000@on@R2}{1.5408}
\@namedef{tb@s@valkey@10000000@on@R3}{0.2121}
\@namedef{tb@s@valkey@10000000@on@R4}{0.1127}
\@namedef{tb@s@valkey@10000000@on@R5}{0.2278}
\@namedef{tb@s@valkey@10000000@on@R6}{0.1345}
\@namedef{tb@s@valkey@10000000@on@U1}{0.1074}
\@namedef{tb@s@valkey@10000000@on@U2}{0.109}
\@namedef{tb@s@valkey@10000000@on@U3}{2.9887}
\@namedef{tb@s@valkey@10000000@on@U4}{0.1383}
\@namedef{tb@s@valkey@10000000@on@U5}{0.1383}
\@namedef{tb@s@valkey@10000000@on@U6}{17.7589}
\@namedef{tb@s@mariadb@500000@off@C1}{0.4849}
\@namedef{tb@s@mariadb@500000@off@C2}{0.5286}
\@namedef{tb@s@mariadb@500000@off@C3}{33.3248}
\@namedef{tb@s@mariadb@500000@off@C4}{0.1419}
\@namedef{tb@s@mariadb@500000@off@C5}{0.3614}
\@namedef{tb@s@mariadb@500000@off@C6}{0.1771}
\@namedef{tb@s@mariadb@500000@off@D1}{85.4856}
\@namedef{tb@s@mariadb@500000@off@D2}{85.5047}
\@namedef{tb@s@mariadb@500000@off@D3}{13.1932}
\@namedef{tb@s@mariadb@500000@off@D4}{153.7643}
\@namedef{tb@s@mariadb@500000@off@D5}{82.2589}
\@namedef{tb@s@mariadb@500000@off@D6}{141.5569}
\@namedef{tb@s@mariadb@500000@off@J1}{0.2406}
\@namedef{tb@s@mariadb@500000@off@J2}{10.5816}
\@namedef{tb@s@mariadb@500000@off@J3}{0.2704}
\@namedef{tb@s@mariadb@500000@off@R1}{42.5687}
\@namedef{tb@s@mariadb@500000@off@R2}{43.4462}
\@namedef{tb@s@mariadb@500000@off@R3}{0.1355}
\@namedef{tb@s@mariadb@500000@off@R4}{19.2783}
\@namedef{tb@s@mariadb@500000@off@R5}{225.6834}
\@namedef{tb@s@mariadb@500000@off@R6}{43.9811}
\@namedef{tb@s@mariadb@500000@off@U1}{0.1762}
\@namedef{tb@s@mariadb@500000@off@U2}{0.173}
\@namedef{tb@s@mariadb@500000@off@U3}{87.8889}
\@namedef{tb@s@mariadb@500000@off@U4}{0.2938}
\@namedef{tb@s@mariadb@500000@off@U5}{0.2405}
\@namedef{tb@s@mariadb@500000@off@U6}{130.5517}
\@namedef{tb@s@mariadb@500000@on@C1}{0.615}
\@namedef{tb@s@mariadb@500000@on@C2}{0.6984}
\@namedef{tb@s@mariadb@500000@on@C3}{44.6764}
\@namedef{tb@s@mariadb@500000@on@C4}{0.1673}
\@namedef{tb@s@mariadb@500000@on@C5}{0.3631}
\@namedef{tb@s@mariadb@500000@on@C6}{0.1926}
\@namedef{tb@s@mariadb@500000@on@D1}{0.7808}
\@namedef{tb@s@mariadb@500000@on@D2}{0.2868}
\@namedef{tb@s@mariadb@500000@on@D3}{2.816}
\@namedef{tb@s@mariadb@500000@on@D4}{0.7393}
\@namedef{tb@s@mariadb@500000@on@D5}{5.5067}
\@namedef{tb@s@mariadb@500000@on@D6}{9.6529}
\@namedef{tb@s@mariadb@500000@on@J1}{0.0977}
\@namedef{tb@s@mariadb@500000@on@J2}{6.6386}
\@namedef{tb@s@mariadb@500000@on@J3}{0.2456}
\@namedef{tb@s@mariadb@500000@on@R1}{0.1013}
\@namedef{tb@s@mariadb@500000@on@R2}{1.464}
\@namedef{tb@s@mariadb@500000@on@R3}{0.1349}
\@namedef{tb@s@mariadb@500000@on@R4}{0.0961}
\@namedef{tb@s@mariadb@500000@on@R5}{20.6573}
\@namedef{tb@s@mariadb@500000@on@R6}{38.9872}
\@namedef{tb@s@mariadb@500000@on@U1}{0.5652}
\@namedef{tb@s@mariadb@500000@on@U2}{0.4333}
\@namedef{tb@s@mariadb@500000@on@U3}{1.1001}
\@namedef{tb@s@mariadb@500000@on@U4}{0.2061}
\@namedef{tb@s@mariadb@500000@on@U5}{0.1454}
\@namedef{tb@s@mariadb@500000@on@U6}{8.2001}
\@namedef{tb@s@mongo@500000@off@C1}{0.5233}
\@namedef{tb@s@mongo@500000@off@C2}{0.9653}
\@namedef{tb@s@mongo@500000@off@C3}{46.7728}
\@namedef{tb@s@mongo@500000@off@C4}{0.2993}
\@namedef{tb@s@mongo@500000@off@C5}{0.3667}
\@namedef{tb@s@mongo@500000@off@C6}{0.9666}
\@namedef{tb@s@mongo@500000@off@D1}{0.9559}
\@namedef{tb@s@mongo@500000@off@D2}{0.8417}
\@namedef{tb@s@mongo@500000@off@D3}{23.5572}
\@namedef{tb@s@mongo@500000@off@D4}{43.961}
\@namedef{tb@s@mongo@500000@off@D5}{133.9634}
\@namedef{tb@s@mongo@500000@off@D6}{156.6257}
\@namedef{tb@s@mongo@500000@off@J1}{0.1985}
\@namedef{tb@s@mongo@500000@off@J2}{16.5435}
\@namedef{tb@s@mongo@500000@off@J3}{0.6063}
\@namedef{tb@s@mongo@500000@off@R1}{0.1073}
\@namedef{tb@s@mongo@500000@off@R2}{147.1645}
\@namedef{tb@s@mongo@500000@off@R3}{0.2366}
\@namedef{tb@s@mongo@500000@off@R4}{22.6944}
\@namedef{tb@s@mongo@500000@off@R5}{109.6199}
\@namedef{tb@s@mongo@500000@off@R6}{103.7005}
\@namedef{tb@s@mongo@500000@off@U1}{0.8361}
\@namedef{tb@s@mongo@500000@off@U2}{0.8206}
\@namedef{tb@s@mongo@500000@off@U3}{134.6072}
\@namedef{tb@s@mongo@500000@off@U4}{0.6669}
\@namedef{tb@s@mongo@500000@off@U5}{0.382}
\@namedef{tb@s@mongo@500000@off@U6}{33.5904}
\@namedef{tb@s@mongo@500000@on@C1}{1.8582}
\@namedef{tb@s@mongo@500000@on@C2}{0.8175}
\@namedef{tb@s@mongo@500000@on@C3}{61.7405}
\@namedef{tb@s@mongo@500000@on@C4}{0.3683}
\@namedef{tb@s@mongo@500000@on@C5}{0.6285}
\@namedef{tb@s@mongo@500000@on@C6}{1.1761}
\@namedef{tb@s@mongo@500000@on@D1}{0.9613}
\@namedef{tb@s@mongo@500000@on@D2}{0.8942}
\@namedef{tb@s@mongo@500000@on@D3}{7.8803}
\@namedef{tb@s@mongo@500000@on@D4}{1.4284}
\@namedef{tb@s@mongo@500000@on@D5}{13.1677}
\@namedef{tb@s@mongo@500000@on@D6}{15.6966}
\@namedef{tb@s@mongo@500000@on@J1}{0.1357}
\@namedef{tb@s@mongo@500000@on@J2}{3.0157}
\@namedef{tb@s@mongo@500000@on@J3}{1.0736}
\@namedef{tb@s@mongo@500000@on@R1}{0.1081}
\@namedef{tb@s@mongo@500000@on@R2}{1.3323}
\@namedef{tb@s@mongo@500000@on@R3}{0.1388}
\@namedef{tb@s@mongo@500000@on@R4}{0.1503}
\@namedef{tb@s@mongo@500000@on@R5}{12.0605}
\@namedef{tb@s@mongo@500000@on@R6}{120.4588}
\@namedef{tb@s@mongo@500000@on@U1}{0.2076}
\@namedef{tb@s@mongo@500000@on@U2}{0.8219}
\@namedef{tb@s@mongo@500000@on@U3}{1.8443}
\@namedef{tb@s@mongo@500000@on@U4}{0.5589}
\@namedef{tb@s@mongo@500000@on@U5}{0.4837}
\@namedef{tb@s@mongo@500000@on@U6}{19.4289}
\@namedef{tb@s@postgres@500000@off@C1}{0.3125}
\@namedef{tb@s@postgres@500000@off@C2}{0.3125}
\@namedef{tb@s@postgres@500000@off@C3}{8.4267}
\@namedef{tb@s@postgres@500000@off@C4}{0.1871}
\@namedef{tb@s@postgres@500000@off@C5}{0.3182}
\@namedef{tb@s@postgres@500000@off@C6}{0.2653}
\@namedef{tb@s@postgres@500000@off@D1}{15.0783}
\@namedef{tb@s@postgres@500000@off@D2}{15.4417}
\@namedef{tb@s@postgres@500000@off@D3}{2.0997}
\@namedef{tb@s@postgres@500000@off@D4}{21.4311}
\@namedef{tb@s@postgres@500000@off@D5}{11.3706}
\@namedef{tb@s@postgres@500000@off@D6}{25.6565}
\@namedef{tb@s@postgres@500000@off@J1}{0.1166}
\@namedef{tb@s@postgres@500000@off@J2}{2.1167}
\@namedef{tb@s@postgres@500000@off@J3}{0.1716}
\@namedef{tb@s@postgres@500000@off@R1}{7.8983}
\@namedef{tb@s@postgres@500000@off@R2}{7.9811}
\@namedef{tb@s@postgres@500000@off@R3}{0.1095}
\@namedef{tb@s@postgres@500000@off@R4}{3.6782}
\@namedef{tb@s@postgres@500000@off@R5}{20.3837}
\@namedef{tb@s@postgres@500000@off@R6}{11.1218}
\@namedef{tb@s@postgres@500000@off@U1}{0.1879}
\@namedef{tb@s@postgres@500000@off@U2}{0.1851}
\@namedef{tb@s@postgres@500000@off@U3}{13.1378}
\@namedef{tb@s@postgres@500000@off@U4}{0.1807}
\@namedef{tb@s@postgres@500000@off@U5}{0.1895}
\@namedef{tb@s@postgres@500000@off@U6}{28.7238}
\@namedef{tb@s@postgres@500000@on@C1}{0.3336}
\@namedef{tb@s@postgres@500000@on@C2}{0.3071}
\@namedef{tb@s@postgres@500000@on@C3}{18.794}
\@namedef{tb@s@postgres@500000@on@C4}{0.1734}
\@namedef{tb@s@postgres@500000@on@C5}{0.3203}
\@namedef{tb@s@postgres@500000@on@C6}{0.2351}
\@namedef{tb@s@postgres@500000@on@D1}{0.3297}
\@namedef{tb@s@postgres@500000@on@D2}{0.2925}
\@namedef{tb@s@postgres@500000@on@D3}{0.3421}
\@namedef{tb@s@postgres@500000@on@D4}{0.3786}
\@namedef{tb@s@postgres@500000@on@D5}{0.6843}
\@namedef{tb@s@postgres@500000@on@D6}{1.641}
\@namedef{tb@s@postgres@500000@on@J1}{0.063}
\@namedef{tb@s@postgres@500000@on@J2}{1.2263}
\@namedef{tb@s@postgres@500000@on@J3}{0.1705}
\@namedef{tb@s@postgres@500000@on@R1}{0.0744}
\@namedef{tb@s@postgres@500000@on@R2}{1.4221}
\@namedef{tb@s@postgres@500000@on@R3}{0.0605}
\@namedef{tb@s@postgres@500000@on@R4}{0.0784}
\@namedef{tb@s@postgres@500000@on@R5}{2.4283}
\@namedef{tb@s@postgres@500000@on@R6}{5.9887}
\@namedef{tb@s@postgres@500000@on@U1}{0.1696}
\@namedef{tb@s@postgres@500000@on@U2}{0.1641}
\@namedef{tb@s@postgres@500000@on@U3}{0.7564}
\@namedef{tb@s@postgres@500000@on@U4}{0.1648}
\@namedef{tb@s@postgres@500000@on@U5}{0.1702}
\@namedef{tb@s@postgres@500000@on@U6}{6.3966}
\@namedef{tb@s@valkey@500000@off@C1}{0.1071}
\@namedef{tb@s@valkey@500000@off@C2}{0.1309}
\@namedef{tb@s@valkey@500000@off@C3}{8.8694}
\@namedef{tb@s@valkey@500000@off@C4}{0.1005}
\@namedef{tb@s@valkey@500000@off@C5}{0.1198}
\@namedef{tb@s@valkey@500000@off@C6}{0.1083}
\@namedef{tb@s@valkey@500000@off@D1}{0.0996}
\@namedef{tb@s@valkey@500000@off@D2}{0.1108}
\@namedef{tb@s@valkey@500000@off@D3}{505.5686}
\@namedef{tb@s@valkey@500000@off@D4}{499.4295}
\@namedef{tb@s@valkey@500000@off@D5}{0.1145}
\@namedef{tb@s@valkey@500000@off@D6}{518.2361}
\@namedef{tb@s@valkey@500000@off@J1}{0.4386}
\@namedef{tb@s@valkey@500000@off@J2}{498.7966}
\@namedef{tb@s@valkey@500000@off@J3}{0.1122}
\@namedef{tb@s@valkey@500000@off@R1}{0.0951}
\@namedef{tb@s@valkey@500000@off@R2}{1.8059}
\@namedef{tb@s@valkey@500000@off@R3}{0.3898}
\@namedef{tb@s@valkey@500000@off@R4}{506.091}
\@namedef{tb@s@valkey@500000@off@R5}{0.4409}
\@namedef{tb@s@valkey@500000@off@R6}{0.1255}
\@namedef{tb@s@valkey@500000@off@U1}{0.1135}
\@namedef{tb@s@valkey@500000@off@U2}{0.1137}
\@namedef{tb@s@valkey@500000@off@U3}{2.0011}
\@namedef{tb@s@valkey@500000@off@U4}{0.129}
\@namedef{tb@s@valkey@500000@off@U5}{0.1155}
\@namedef{tb@s@valkey@500000@off@U6}{479.7103}
\@namedef{tb@s@valkey@500000@on@C1}{0.3559}
\@namedef{tb@s@valkey@500000@on@C2}{0.1325}
\@namedef{tb@s@valkey@500000@on@C3}{6.9666}
\@namedef{tb@s@valkey@500000@on@C4}{0.0805}
\@namedef{tb@s@valkey@500000@on@C5}{0.17}
\@namedef{tb@s@valkey@500000@on@C6}{0.0885}
\@namedef{tb@s@valkey@500000@on@D1}{0.1229}
\@namedef{tb@s@valkey@500000@on@D2}{0.1103}
\@namedef{tb@s@valkey@500000@on@D3}{5.9784}
\@namedef{tb@s@valkey@500000@on@D4}{0.2829}
\@namedef{tb@s@valkey@500000@on@D5}{0.1394}
\@namedef{tb@s@valkey@500000@on@D6}{17.9773}
\@namedef{tb@s@valkey@500000@on@J1}{0.3474}
\@namedef{tb@s@valkey@500000@on@J2}{19.8033}
\@namedef{tb@s@valkey@500000@on@J3}{0.0976}
\@namedef{tb@s@valkey@500000@on@R1}{0.0946}
\@namedef{tb@s@valkey@500000@on@R2}{2.138}
\@namedef{tb@s@valkey@500000@on@R3}{0.1525}
\@namedef{tb@s@valkey@500000@on@R4}{0.1186}
\@namedef{tb@s@valkey@500000@on@R5}{0.2576}
\@namedef{tb@s@valkey@500000@on@R6}{0.1304}
\@namedef{tb@s@valkey@500000@on@U1}{0.1053}
\@namedef{tb@s@valkey@500000@on@U2}{0.1025}
\@namedef{tb@s@valkey@500000@on@U3}{2.4093}
\@namedef{tb@s@valkey@500000@on@U4}{0.1619}
\@namedef{tb@s@valkey@500000@on@U5}{0.139}
\@namedef{tb@s@valkey@500000@on@U6}{17.4677}
\@namedef{tb@c@mariadb@1000000@cond@1000}{18.41}
\@namedef{tb@p@mariadb@1000000@cond@1000}{13650.9425}
\@namedef{tb@k@mariadb@1000000@cond@1000}{8587.6667}
\@namedef{tb@c@mariadb@1000000@cond@100}{6155.8067}
\@namedef{tb@p@mariadb@1000000@cond@100}{6.3272}
\@namedef{tb@k@mariadb@1000000@cond@100}{3406.6667}
\@namedef{tb@c@mariadb@1000000@cond@10}{6209.95}
\@namedef{tb@p@mariadb@1000000@cond@10}{0.7074}
\@namedef{tb@k@mariadb@1000000@cond@10}{1966.3333}
\@namedef{tb@c@mariadb@1000000@cond@1}{5208.0833}
\@namedef{tb@p@mariadb@1000000@cond@1}{0.2535}
\@namedef{tb@k@mariadb@1000000@cond@1}{0}
\@namedef{tb@c@mariadb@1000000@forupdate@1000}{14.9533}
\@namedef{tb@p@mariadb@1000000@forupdate@1000}{15199.4809}
\@namedef{tb@k@mariadb@1000000@forupdate@1000}{8820}
\@namedef{tb@c@mariadb@1000000@forupdate@100}{6970.6533}
\@namedef{tb@p@mariadb@1000000@forupdate@100}{5.1933}
\@namedef{tb@k@mariadb@1000000@forupdate@100}{3634}
\@namedef{tb@c@mariadb@1000000@forupdate@10}{5943.85}
\@namedef{tb@p@mariadb@1000000@forupdate@10}{0.7945}
\@namedef{tb@k@mariadb@1000000@forupdate@10}{1876}
\@namedef{tb@c@mariadb@1000000@forupdate@1}{4349.4467}
\@namedef{tb@p@mariadb@1000000@forupdate@1}{0.2994}
\@namedef{tb@k@mariadb@1000000@forupdate@1}{0}
\@namedef{tb@c@mariadb@1000000@nowait@1000}{327.8767}
\@namedef{tb@p@mariadb@1000000@nowait@1000}{639.6985}
\@namedef{tb@k@mariadb@1000000@nowait@1000}{12077}
\@namedef{tb@c@mariadb@1000000@nowait@100}{8232.3467}
\@namedef{tb@p@mariadb@1000000@nowait@100}{4.1912}
\@namedef{tb@k@mariadb@1000000@nowait@100}{3441}
\@namedef{tb@c@mariadb@1000000@nowait@10}{6037.4567}
\@namedef{tb@p@mariadb@1000000@nowait@10}{0.7757}
\@namedef{tb@k@mariadb@1000000@nowait@10}{1987}
\@namedef{tb@c@mariadb@1000000@nowait@1}{4564.6867}
\@namedef{tb@p@mariadb@1000000@nowait@1}{0.2997}
\@namedef{tb@k@mariadb@1000000@nowait@1}{0}
\@namedef{tb@c@mariadb@1000000@serializable@1000}{18.66}
\@namedef{tb@p@mariadb@1000000@serializable@1000}{6836.9523}
\@namedef{tb@k@mariadb@1000000@serializable@1000}{13917.3333}
\@namedef{tb@c@mariadb@1000000@serializable@100}{5179.8667}
\@namedef{tb@p@mariadb@1000000@serializable@100}{8.3905}
\@namedef{tb@k@mariadb@1000000@serializable@100}{3529.3333}
\@namedef{tb@c@mariadb@1000000@serializable@10}{5197.0967}
\@namedef{tb@p@mariadb@1000000@serializable@10}{0.8576}
\@namedef{tb@k@mariadb@1000000@serializable@10}{1925.6667}
\@namedef{tb@c@mariadb@1000000@serializable@1}{3895.3733}
\@namedef{tb@p@mariadb@1000000@serializable@1}{0.3614}
\@namedef{tb@k@mariadb@1000000@serializable@1}{0}
\@namedef{tb@c@mongo@1000000@atomic@1000}{2163.7567}
\@namedef{tb@p@mongo@1000000@atomic@1000}{79.4654}
\@namedef{tb@k@mongo@1000000@atomic@1000}{6111.3333}
\@namedef{tb@c@mongo@1000000@atomic@100}{3417.32}
\@namedef{tb@p@mongo@1000000@atomic@100}{9.6189}
\@namedef{tb@k@mongo@1000000@atomic@100}{3247.6667}
\@namedef{tb@c@mongo@1000000@atomic@10}{2356.2967}
\@namedef{tb@p@mongo@1000000@atomic@10}{2.3998}
\@namedef{tb@k@mongo@1000000@atomic@10}{1929}
\@namedef{tb@c@mongo@1000000@atomic@1}{703.7}
\@namedef{tb@p@mongo@1000000@atomic@1}{2.3422}
\@namedef{tb@k@mongo@1000000@atomic@1}{0}
\@namedef{tb@c@mongo@1000000@txn@1000}{1582.23}
\@namedef{tb@p@mongo@1000000@txn@1000}{108.5421}
\@namedef{tb@k@mongo@1000000@txn@1000}{5400}
\@namedef{tb@c@mongo@1000000@txn@100}{2637.83}
\@namedef{tb@p@mongo@1000000@txn@100}{8.8608}
\@namedef{tb@k@mongo@1000000@txn@100}{3133.6667}
\@namedef{tb@c@mongo@1000000@txn@10}{1698.0533}
\@namedef{tb@p@mongo@1000000@txn@10}{2.2273}
\@namedef{tb@k@mongo@1000000@txn@10}{1920.6667}
\@namedef{tb@c@mongo@1000000@txn@1}{962.87}
\@namedef{tb@p@mongo@1000000@txn@1}{1.6143}
\@namedef{tb@k@mongo@1000000@txn@1}{0}
\@namedef{tb@c@postgres@1000000@cond@1000}{2602.7133}
\@namedef{tb@p@postgres@1000000@cond@1000}{141.8514}
\@namedef{tb@k@postgres@1000000@cond@1000}{7495}
\@namedef{tb@c@postgres@1000000@cond@100}{6979.2533}
\@namedef{tb@p@postgres@1000000@cond@100}{5.0356}
\@namedef{tb@k@postgres@1000000@cond@100}{3964}
\@namedef{tb@c@postgres@1000000@cond@10}{6261.7733}
\@namedef{tb@p@postgres@1000000@cond@10}{0.7263}
\@namedef{tb@k@postgres@1000000@cond@10}{1941.3333}
\@namedef{tb@c@postgres@1000000@cond@1}{5905.04}
\@namedef{tb@p@postgres@1000000@cond@1}{0.2146}
\@namedef{tb@k@postgres@1000000@cond@1}{0}
\@namedef{tb@c@postgres@1000000@forupdate@1000}{3731.9567}
\@namedef{tb@p@postgres@1000000@forupdate@1000}{79.3961}
\@namedef{tb@k@postgres@1000000@forupdate@1000}{3685}
\@namedef{tb@c@postgres@1000000@forupdate@100}{6112.9567}
\@namedef{tb@p@postgres@1000000@forupdate@100}{5.1388}
\@namedef{tb@k@postgres@1000000@forupdate@100}{3107}
\@namedef{tb@c@postgres@1000000@forupdate@10}{3559.4467}
\@namedef{tb@p@postgres@1000000@forupdate@10}{0.9454}
\@namedef{tb@k@postgres@1000000@forupdate@10}{1902}
\@namedef{tb@c@postgres@1000000@forupdate@1}{4325.8367}
\@namedef{tb@p@postgres@1000000@forupdate@1}{0.2896}
\@namedef{tb@k@postgres@1000000@forupdate@1}{0}
\@namedef{tb@c@postgres@1000000@nowait@1000}{1089.8267}
\@namedef{tb@p@postgres@1000000@nowait@1000}{91.8871}
\@namedef{tb@k@postgres@1000000@nowait@1000}{17419}
\@namedef{tb@c@postgres@1000000@nowait@100}{4427.3633}
\@namedef{tb@p@postgres@1000000@nowait@100}{5.5765}
\@namedef{tb@k@postgres@1000000@nowait@100}{3879}
\@namedef{tb@c@postgres@1000000@nowait@10}{3877.44}
\@namedef{tb@p@postgres@1000000@nowait@10}{0.7776}
\@namedef{tb@k@postgres@1000000@nowait@10}{1958.3333}
\@namedef{tb@c@postgres@1000000@nowait@1}{4617.38}
\@namedef{tb@p@postgres@1000000@nowait@1}{0.3268}
\@namedef{tb@k@postgres@1000000@nowait@1}{0}
\@namedef{tb@c@postgres@1000000@serializable@1000}{1232.3967}
\@namedef{tb@p@postgres@1000000@serializable@1000}{138.3564}
\@namedef{tb@k@postgres@1000000@serializable@1000}{8905.6667}
\@namedef{tb@c@postgres@1000000@serializable@100}{3063}
\@namedef{tb@p@postgres@1000000@serializable@100}{10.0538}
\@namedef{tb@k@postgres@1000000@serializable@100}{4504.6667}
\@namedef{tb@c@postgres@1000000@serializable@10}{4735.7}
\@namedef{tb@p@postgres@1000000@serializable@10}{0.9092}
\@namedef{tb@k@postgres@1000000@serializable@10}{2045}
\@namedef{tb@c@postgres@1000000@serializable@1}{4333.66}
\@namedef{tb@p@postgres@1000000@serializable@1}{0.2826}
\@namedef{tb@k@postgres@1000000@serializable@1}{0}
\@namedef{tb@c@valkey@1000000@lua@1000}{35009.4333}
\@namedef{tb@p@valkey@1000000@lua@1000}{4.2551}
\@namedef{tb@k@valkey@1000000@lua@1000}{4088.6667}
\@namedef{tb@c@valkey@1000000@lua@100}{35824.8}
\@namedef{tb@p@valkey@1000000@lua@100}{1.0347}
\@namedef{tb@k@valkey@1000000@lua@100}{2957}
\@namedef{tb@c@valkey@1000000@lua@10}{28600.17}
\@namedef{tb@p@valkey@1000000@lua@10}{0.1215}
\@namedef{tb@k@valkey@1000000@lua@10}{1900.3333}
\@namedef{tb@c@valkey@1000000@lua@1}{16917.84}
\@namedef{tb@p@valkey@1000000@lua@1}{0.0693}
\@namedef{tb@k@valkey@1000000@lua@1}{0}
\@namedef{tb@c@valkey@1000000@setnx@1000}{27634.77}
\@namedef{tb@p@valkey@1000000@setnx@1000}{6.8129}
\@namedef{tb@k@valkey@1000000@setnx@1000}{6674.6667}
\@namedef{tb@c@valkey@1000000@setnx@100}{51762.6867}
\@namedef{tb@p@valkey@1000000@setnx@100}{0.7118}
\@namedef{tb@k@valkey@1000000@setnx@100}{3151.3333}
\@namedef{tb@c@valkey@1000000@setnx@10}{22142.8567}
\@namedef{tb@p@valkey@1000000@setnx@10}{0.2433}
\@namedef{tb@k@valkey@1000000@setnx@10}{1924}
\@namedef{tb@c@valkey@1000000@setnx@1}{7896.9267}
\@namedef{tb@p@valkey@1000000@setnx@1}{0.1506}
\@namedef{tb@k@valkey@1000000@setnx@1}{0}
\@namedef{tb@c@valkey@1000000@watch@1000}{90.1833}
\@namedef{tb@p@valkey@1000000@watch@1000}{26.2472}
\@namedef{tb@k@valkey@1000000@watch@1000}{271904}
\@namedef{tb@c@valkey@1000000@watch@100}{1223.3167}
\@namedef{tb@p@valkey@1000000@watch@100}{2.0023}
\@namedef{tb@k@valkey@1000000@watch@100}{26630.6667}
\@namedef{tb@c@valkey@1000000@watch@10}{2856.23}
\@namedef{tb@p@valkey@1000000@watch@10}{0.5956}
\@namedef{tb@k@valkey@1000000@watch@10}{4058.6667}
\@namedef{tb@c@valkey@1000000@watch@1}{4593.0433}
\@namedef{tb@p@valkey@1000000@watch@1}{0.2932}
\@namedef{tb@k@valkey@1000000@watch@1}{0}
\@namedef{tb@c@mariadb@10000000@cond@1000}{24.5333}
\@namedef{tb@p@mariadb@10000000@cond@1000}{15453.7684}
\@namedef{tb@k@mariadb@10000000@cond@1000}{7399.6667}
\@namedef{tb@c@mariadb@10000000@cond@100}{2117.2033}
\@namedef{tb@p@mariadb@10000000@cond@100}{60.5437}
\@namedef{tb@k@mariadb@10000000@cond@100}{3693.3333}
\@namedef{tb@c@mariadb@10000000@cond@10}{4690.0933}
\@namedef{tb@p@mariadb@10000000@cond@10}{0.9838}
\@namedef{tb@k@mariadb@10000000@cond@10}{1939.3333}
\@namedef{tb@c@mariadb@10000000@cond@1}{4059.5133}
\@namedef{tb@p@mariadb@10000000@cond@1}{0.4238}
\@namedef{tb@k@mariadb@10000000@cond@1}{0}
\@namedef{tb@c@mariadb@10000000@forupdate@1000}{56.7}
\@namedef{tb@p@mariadb@10000000@forupdate@1000}{9095.479}
\@namedef{tb@k@mariadb@10000000@forupdate@1000}{7895.6667}
\@namedef{tb@c@mariadb@10000000@forupdate@100}{7395.7967}
\@namedef{tb@p@mariadb@10000000@forupdate@100}{5.0858}
\@namedef{tb@k@mariadb@10000000@forupdate@100}{3443.6667}
\@namedef{tb@c@mariadb@10000000@forupdate@10}{5834.2267}
\@namedef{tb@p@mariadb@10000000@forupdate@10}{0.8848}
\@namedef{tb@k@mariadb@10000000@forupdate@10}{1945.3333}
\@namedef{tb@c@mariadb@10000000@forupdate@1}{3962.21}
\@namedef{tb@p@mariadb@10000000@forupdate@1}{0.3654}
\@namedef{tb@k@mariadb@10000000@forupdate@1}{0}
\@namedef{tb@c@mariadb@10000000@nowait@1000}{585.2067}
\@namedef{tb@p@mariadb@10000000@nowait@1000}{355.1633}
\@namedef{tb@k@mariadb@10000000@nowait@1000}{11168}
\@namedef{tb@c@mariadb@10000000@nowait@100}{8379.87}
\@namedef{tb@p@mariadb@10000000@nowait@100}{4.1405}
\@namedef{tb@k@mariadb@10000000@nowait@100}{3535.3333}
\@namedef{tb@c@mariadb@10000000@nowait@10}{5525.65}
\@namedef{tb@p@mariadb@10000000@nowait@10}{0.9913}
\@namedef{tb@k@mariadb@10000000@nowait@10}{1980}
\@namedef{tb@c@mariadb@10000000@nowait@1}{3716.54}
\@namedef{tb@p@mariadb@10000000@nowait@1}{0.4925}
\@namedef{tb@k@mariadb@10000000@nowait@1}{0}
\@namedef{tb@c@mariadb@10000000@serializable@1000}{58.4533}
\@namedef{tb@p@mariadb@10000000@serializable@1000}{2240.8186}
\@namedef{tb@k@mariadb@10000000@serializable@1000}{18612}
\@namedef{tb@c@mariadb@10000000@serializable@100}{5805.23}
\@namedef{tb@p@mariadb@10000000@serializable@100}{5.7729}
\@namedef{tb@k@mariadb@10000000@serializable@100}{3523}
\@namedef{tb@c@mariadb@10000000@serializable@10}{4919.73}
\@namedef{tb@p@mariadb@10000000@serializable@10}{0.9825}
\@namedef{tb@k@mariadb@10000000@serializable@10}{1938.6667}
\@namedef{tb@c@mariadb@10000000@serializable@1}{3690.2167}
\@namedef{tb@p@mariadb@10000000@serializable@1}{0.3718}
\@namedef{tb@k@mariadb@10000000@serializable@1}{0}
\@namedef{tb@c@mongo@10000000@atomic@1000}{2047.1067}
\@namedef{tb@p@mongo@10000000@atomic@1000}{85.4749}
\@namedef{tb@k@mongo@10000000@atomic@1000}{5214.3333}
\@namedef{tb@c@mongo@10000000@atomic@100}{3473.27}
\@namedef{tb@p@mongo@10000000@atomic@100}{8.8178}
\@namedef{tb@k@mongo@10000000@atomic@100}{3293.3333}
\@namedef{tb@c@mongo@10000000@atomic@10}{2379.08}
\@namedef{tb@p@mongo@10000000@atomic@10}{2.4152}
\@namedef{tb@k@mongo@10000000@atomic@10}{1932.6667}
\@namedef{tb@c@mongo@10000000@atomic@1}{731.3}
\@namedef{tb@p@mongo@10000000@atomic@1}{2.3081}
\@namedef{tb@k@mongo@10000000@atomic@1}{0}
\@namedef{tb@c@mongo@10000000@txn@1000}{1613.8767}
\@namedef{tb@p@mongo@10000000@txn@1000}{137.596}
\@namedef{tb@k@mongo@10000000@txn@1000}{4025}
\@namedef{tb@c@mongo@10000000@txn@100}{1925.5833}
\@namedef{tb@p@mongo@10000000@txn@100}{12.9007}
\@namedef{tb@k@mongo@10000000@txn@100}{3219}
\@namedef{tb@c@mongo@10000000@txn@10}{1244.23}
\@namedef{tb@p@mongo@10000000@txn@10}{3.2637}
\@namedef{tb@k@mongo@10000000@txn@10}{1941}
\@namedef{tb@c@mongo@10000000@txn@1}{926.1833}
\@namedef{tb@p@mongo@10000000@txn@1}{1.8203}
\@namedef{tb@k@mongo@10000000@txn@1}{0}
\@namedef{tb@c@postgres@10000000@cond@1000}{1975.43}
\@namedef{tb@p@postgres@10000000@cond@1000}{208.7821}
\@namedef{tb@k@postgres@10000000@cond@1000}{9444}
\@namedef{tb@c@postgres@10000000@cond@100}{6923.3067}
\@namedef{tb@p@postgres@10000000@cond@100}{6.8298}
\@namedef{tb@k@postgres@10000000@cond@100}{3410.3333}
\@namedef{tb@c@postgres@10000000@cond@10}{5808.3833}
\@namedef{tb@p@postgres@10000000@cond@10}{0.7477}
\@namedef{tb@k@postgres@10000000@cond@10}{1982.6667}
\@namedef{tb@c@postgres@10000000@cond@1}{5548.4267}
\@namedef{tb@p@postgres@10000000@cond@1}{0.2263}
\@namedef{tb@k@postgres@10000000@cond@1}{0}
\@namedef{tb@c@postgres@10000000@forupdate@1000}{3488.9733}
\@namedef{tb@p@postgres@10000000@forupdate@1000}{86.7856}
\@namedef{tb@k@postgres@10000000@forupdate@1000}{3502.6667}
\@namedef{tb@c@postgres@10000000@forupdate@100}{6498.2933}
\@namedef{tb@p@postgres@10000000@forupdate@100}{4.5417}
\@namedef{tb@k@postgres@10000000@forupdate@100}{3170.6667}
\@namedef{tb@c@postgres@10000000@forupdate@10}{3838.9967}
\@namedef{tb@p@postgres@10000000@forupdate@10}{0.7867}
\@namedef{tb@k@postgres@10000000@forupdate@10}{1933.3333}
\@namedef{tb@c@postgres@10000000@forupdate@1}{4413.8833}
\@namedef{tb@p@postgres@10000000@forupdate@1}{0.2797}
\@namedef{tb@k@postgres@10000000@forupdate@1}{0}
\@namedef{tb@c@postgres@10000000@nowait@1000}{1066.3933}
\@namedef{tb@p@postgres@10000000@nowait@1000}{91.66}
\@namedef{tb@k@postgres@10000000@nowait@1000}{15663}
\@namedef{tb@c@postgres@10000000@nowait@100}{4892.1033}
\@namedef{tb@p@postgres@10000000@nowait@100}{4.8752}
\@namedef{tb@k@postgres@10000000@nowait@100}{3957.6667}
\@namedef{tb@c@postgres@10000000@nowait@10}{3879.7067}
\@namedef{tb@p@postgres@10000000@nowait@10}{0.7844}
\@namedef{tb@k@postgres@10000000@nowait@10}{1934.6667}
\@namedef{tb@c@postgres@10000000@nowait@1}{4399.73}
\@namedef{tb@p@postgres@10000000@nowait@1}{0.2887}
\@namedef{tb@k@postgres@10000000@nowait@1}{0}
\@namedef{tb@c@postgres@10000000@serializable@1000}{912.06}
\@namedef{tb@p@postgres@10000000@serializable@1000}{158.8357}
\@namedef{tb@k@postgres@10000000@serializable@1000}{11289.6667}
\@namedef{tb@c@postgres@10000000@serializable@100}{2830.1133}
\@namedef{tb@p@postgres@10000000@serializable@100}{10.7494}
\@namedef{tb@k@postgres@10000000@serializable@100}{4133.3333}
\@namedef{tb@c@postgres@10000000@serializable@10}{4528.5033}
\@namedef{tb@p@postgres@10000000@serializable@10}{0.8927}
\@namedef{tb@k@postgres@10000000@serializable@10}{2207}
\@namedef{tb@c@postgres@10000000@serializable@1}{4338.3067}
\@namedef{tb@p@postgres@10000000@serializable@1}{0.2815}
\@namedef{tb@k@postgres@10000000@serializable@1}{0}
\@namedef{tb@c@valkey@10000000@lua@1000}{46619.4967}
\@namedef{tb@p@valkey@10000000@lua@1000}{3.2236}
\@namedef{tb@k@valkey@10000000@lua@1000}{3369}
\@namedef{tb@c@valkey@10000000@lua@100}{36871.4767}
\@namedef{tb@p@valkey@10000000@lua@100}{0.8479}
\@namedef{tb@k@valkey@10000000@lua@100}{2988}
\@namedef{tb@c@valkey@10000000@lua@10}{24973.7833}
\@namedef{tb@p@valkey@10000000@lua@10}{0.165}
\@namedef{tb@k@valkey@10000000@lua@10}{1904.3333}
\@namedef{tb@c@valkey@10000000@lua@1}{16835.01}
\@namedef{tb@p@valkey@10000000@lua@1}{0.0695}
\@namedef{tb@k@valkey@10000000@lua@1}{0}
\@namedef{tb@c@valkey@10000000@setnx@1000}{22170.51}
\@namedef{tb@p@valkey@10000000@setnx@1000}{5.6186}
\@namedef{tb@k@valkey@10000000@setnx@1000}{7537.6667}
\@namedef{tb@c@valkey@10000000@setnx@100}{45901.24}
\@namedef{tb@p@valkey@10000000@setnx@100}{0.8077}
\@namedef{tb@k@valkey@10000000@setnx@100}{3201.3333}
\@namedef{tb@c@valkey@10000000@setnx@10}{21146.7967}
\@namedef{tb@p@valkey@10000000@setnx@10}{0.2554}
\@namedef{tb@k@valkey@10000000@setnx@10}{1958}
\@namedef{tb@c@valkey@10000000@setnx@1}{7609.75}
\@namedef{tb@p@valkey@10000000@setnx@1}{0.1498}
\@namedef{tb@k@valkey@10000000@setnx@1}{0}
\@namedef{tb@c@valkey@10000000@watch@1000}{76.07}
\@namedef{tb@p@valkey@10000000@watch@1000}{32.5353}
\@namedef{tb@k@valkey@10000000@watch@1000}{291179.6667}
\@namedef{tb@c@valkey@10000000@watch@100}{1175.82}
\@namedef{tb@p@valkey@10000000@watch@100}{1.8837}
\@namedef{tb@k@valkey@10000000@watch@100}{27607.6667}
\@namedef{tb@c@valkey@10000000@watch@10}{3282.6133}
\@namedef{tb@p@valkey@10000000@watch@10}{0.4893}
\@namedef{tb@k@valkey@10000000@watch@10}{4140.6667}
\@namedef{tb@c@valkey@10000000@watch@1}{4654.4867}
\@namedef{tb@p@valkey@10000000@watch@1}{0.2628}
\@namedef{tb@k@valkey@10000000@watch@1}{0}
\@namedef{tb@c@mariadb@500000@cond@1000}{36.42}
\@namedef{tb@p@mariadb@500000@cond@1000}{8040.7325}
\@namedef{tb@k@mariadb@500000@cond@1000}{7180.3333}
\@namedef{tb@c@mariadb@500000@cond@100}{3995.0333}
\@namedef{tb@p@mariadb@500000@cond@100}{22.9426}
\@namedef{tb@k@mariadb@500000@cond@100}{2852}
\@namedef{tb@c@mariadb@500000@cond@10}{4507.6767}
\@namedef{tb@p@mariadb@500000@cond@10}{1.6514}
\@namedef{tb@k@mariadb@500000@cond@10}{1851.3333}
\@namedef{tb@c@mariadb@500000@cond@1}{3093.4967}
\@namedef{tb@p@mariadb@500000@cond@1}{1.0418}
\@namedef{tb@k@mariadb@500000@cond@1}{0}
\@namedef{tb@c@mariadb@500000@forupdate@1000}{24.1867}
\@namedef{tb@p@mariadb@500000@forupdate@1000}{8684.8847}
\@namedef{tb@k@mariadb@500000@forupdate@1000}{7493}
\@namedef{tb@c@mariadb@500000@forupdate@100}{6096.3667}
\@namedef{tb@p@mariadb@500000@forupdate@100}{7.4368}
\@namedef{tb@k@mariadb@500000@forupdate@100}{3522.3333}
\@namedef{tb@c@mariadb@500000@forupdate@10}{4671.0333}
\@namedef{tb@p@mariadb@500000@forupdate@10}{1.5283}
\@namedef{tb@k@mariadb@500000@forupdate@10}{1912.3333}
\@namedef{tb@c@mariadb@500000@forupdate@1}{2893.56}
\@namedef{tb@p@mariadb@500000@forupdate@1}{1.0232}
\@namedef{tb@k@mariadb@500000@forupdate@1}{0}
\@namedef{tb@c@mariadb@500000@nowait@1000}{187.82}
\@namedef{tb@p@mariadb@500000@nowait@1000}{1465.5827}
\@namedef{tb@k@mariadb@500000@nowait@1000}{12216.3333}
\@namedef{tb@c@mariadb@500000@nowait@100}{6948.9333}
\@namedef{tb@p@mariadb@500000@nowait@100}{5.8085}
\@namedef{tb@k@mariadb@500000@nowait@100}{3508.6667}
\@namedef{tb@c@mariadb@500000@nowait@10}{4226.26}
\@namedef{tb@p@mariadb@500000@nowait@10}{1.7658}
\@namedef{tb@k@mariadb@500000@nowait@10}{1981}
\@namedef{tb@c@mariadb@500000@nowait@1}{3081.7433}
\@namedef{tb@p@mariadb@500000@nowait@1}{0.7398}
\@namedef{tb@k@mariadb@500000@nowait@1}{0}
\@namedef{tb@c@mariadb@500000@serializable@1000}{24.3033}
\@namedef{tb@p@mariadb@500000@serializable@1000}{5824.5518}
\@namedef{tb@k@mariadb@500000@serializable@1000}{14939.3333}
\@namedef{tb@c@mariadb@500000@serializable@100}{6570.03}
\@namedef{tb@p@mariadb@500000@serializable@100}{6.0316}
\@namedef{tb@k@mariadb@500000@serializable@100}{3257.6667}
\@namedef{tb@c@mariadb@500000@serializable@10}{4123.6833}
\@namedef{tb@p@mariadb@500000@serializable@10}{1.6074}
\@namedef{tb@k@mariadb@500000@serializable@10}{1936}
\@namedef{tb@c@mariadb@500000@serializable@1}{2538.92}
\@namedef{tb@p@mariadb@500000@serializable@1}{0.963}
\@namedef{tb@k@mariadb@500000@serializable@1}{0}
\@namedef{tb@c@mongo@500000@atomic@1000}{1796.9533}
\@namedef{tb@p@mongo@500000@atomic@1000}{83.6313}
\@namedef{tb@k@mongo@500000@atomic@1000}{6422}
\@namedef{tb@c@mongo@500000@atomic@100}{3329.4667}
\@namedef{tb@p@mongo@500000@atomic@100}{9.8339}
\@namedef{tb@k@mongo@500000@atomic@100}{3046.6667}
\@namedef{tb@c@mongo@500000@atomic@10}{2476.4933}
\@namedef{tb@p@mongo@500000@atomic@10}{2.3068}
\@namedef{tb@k@mongo@500000@atomic@10}{1916.3333}
\@namedef{tb@c@mongo@500000@atomic@1}{724.91}
\@namedef{tb@p@mongo@500000@atomic@1}{2.3665}
\@namedef{tb@k@mongo@500000@atomic@1}{0}
\@namedef{tb@c@mongo@500000@txn@1000}{1806.93}
\@namedef{tb@p@mongo@500000@txn@1000}{111.5107}
\@namedef{tb@k@mongo@500000@txn@1000}{4303}
\@namedef{tb@c@mongo@500000@txn@100}{2618.24}
\@namedef{tb@p@mongo@500000@txn@100}{9.3056}
\@namedef{tb@k@mongo@500000@txn@100}{3091.3333}
\@namedef{tb@c@mongo@500000@txn@10}{1733.8167}
\@namedef{tb@p@mongo@500000@txn@10}{2.1677}
\@namedef{tb@k@mongo@500000@txn@10}{1930}
\@namedef{tb@c@mongo@500000@txn@1}{958.9033}
\@namedef{tb@p@mongo@500000@txn@1}{1.5965}
\@namedef{tb@k@mongo@500000@txn@1}{0}
\@namedef{tb@c@postgres@500000@cond@1000}{1820.3867}
\@namedef{tb@p@postgres@500000@cond@1000}{206.6391}
\@namedef{tb@k@postgres@500000@cond@1000}{6993.3333}
\@namedef{tb@c@postgres@500000@cond@100}{7427.12}
\@namedef{tb@p@postgres@500000@cond@100}{8.6861}
\@namedef{tb@k@postgres@500000@cond@100}{3739.6667}
\@namedef{tb@c@postgres@500000@cond@10}{6974.53}
\@namedef{tb@p@postgres@500000@cond@10}{0.669}
\@namedef{tb@k@postgres@500000@cond@10}{1934.6667}
\@namedef{tb@c@postgres@500000@cond@1}{5244.13}
\@namedef{tb@p@postgres@500000@cond@1}{0.2283}
\@namedef{tb@k@postgres@500000@cond@1}{0}
\@namedef{tb@c@postgres@500000@forupdate@1000}{4055.0967}
\@namedef{tb@p@postgres@500000@forupdate@1000}{77.6079}
\@namedef{tb@k@postgres@500000@forupdate@1000}{3303.3333}
\@namedef{tb@c@postgres@500000@forupdate@100}{7106.1967}
\@namedef{tb@p@postgres@500000@forupdate@100}{4.121}
\@namedef{tb@k@postgres@500000@forupdate@100}{3120}
\@namedef{tb@c@postgres@500000@forupdate@10}{3958.26}
\@namedef{tb@p@postgres@500000@forupdate@10}{0.7707}
\@namedef{tb@k@postgres@500000@forupdate@10}{1918}
\@namedef{tb@c@postgres@500000@forupdate@1}{4341.16}
\@namedef{tb@p@postgres@500000@forupdate@1}{0.2848}
\@namedef{tb@k@postgres@500000@forupdate@1}{0}
\@namedef{tb@c@postgres@500000@nowait@1000}{1082.5133}
\@namedef{tb@p@postgres@500000@nowait@1000}{93.4461}
\@namedef{tb@k@postgres@500000@nowait@1000}{15633.3333}
\@namedef{tb@c@postgres@500000@nowait@100}{5923.3}
\@namedef{tb@p@postgres@500000@nowait@100}{4.0891}
\@namedef{tb@k@postgres@500000@nowait@100}{3694.6667}
\@namedef{tb@c@postgres@500000@nowait@10}{3906.61}
\@namedef{tb@p@postgres@500000@nowait@10}{0.7652}
\@namedef{tb@k@postgres@500000@nowait@10}{1967.6667}
\@namedef{tb@c@postgres@500000@nowait@1}{4213.5267}
\@namedef{tb@p@postgres@500000@nowait@1}{0.2818}
\@namedef{tb@k@postgres@500000@nowait@1}{0}
\@namedef{tb@c@postgres@500000@serializable@1000}{1229.58}
\@namedef{tb@p@postgres@500000@serializable@1000}{127.997}
\@namedef{tb@k@postgres@500000@serializable@1000}{9922.3333}
\@namedef{tb@c@postgres@500000@serializable@100}{3689.4033}
\@namedef{tb@p@postgres@500000@serializable@100}{9.1986}
\@namedef{tb@k@postgres@500000@serializable@100}{3741}
\@namedef{tb@c@postgres@500000@serializable@10}{4837.1067}
\@namedef{tb@p@postgres@500000@serializable@10}{0.886}
\@namedef{tb@k@postgres@500000@serializable@10}{2106.3333}
\@namedef{tb@c@postgres@500000@serializable@1}{4460.9067}
\@namedef{tb@p@postgres@500000@serializable@1}{0.2763}
\@namedef{tb@k@postgres@500000@serializable@1}{0}
\@namedef{tb@c@valkey@500000@lua@1000}{35253.1533}
\@namedef{tb@p@valkey@500000@lua@1000}{4.182}
\@namedef{tb@k@valkey@500000@lua@1000}{3938.3333}
\@namedef{tb@c@valkey@500000@lua@100}{33954.8933}
\@namedef{tb@p@valkey@500000@lua@100}{0.832}
\@namedef{tb@k@valkey@500000@lua@100}{2936.3333}
\@namedef{tb@c@valkey@500000@lua@10}{27934.9833}
\@namedef{tb@p@valkey@500000@lua@10}{0.1217}
\@namedef{tb@k@valkey@500000@lua@10}{1922.3333}
\@namedef{tb@c@valkey@500000@lua@1}{16880.74}
\@namedef{tb@p@valkey@500000@lua@1}{0.0691}
\@namedef{tb@k@valkey@500000@lua@1}{0}
\@namedef{tb@c@valkey@500000@setnx@1000}{28303.14}
\@namedef{tb@p@valkey@500000@setnx@1000}{6.0902}
\@namedef{tb@k@valkey@500000@setnx@1000}{6056.3333}
\@namedef{tb@c@valkey@500000@setnx@100}{46741.6567}
\@namedef{tb@p@valkey@500000@setnx@100}{0.759}
\@namedef{tb@k@valkey@500000@setnx@100}{3199.3333}
\@namedef{tb@c@valkey@500000@setnx@10}{23070.0333}
\@namedef{tb@p@valkey@500000@setnx@10}{0.2283}
\@namedef{tb@k@valkey@500000@setnx@10}{1942.6667}
\@namedef{tb@c@valkey@500000@setnx@1}{7471.5733}
\@namedef{tb@p@valkey@500000@setnx@1}{0.1571}
\@namedef{tb@k@valkey@500000@setnx@1}{0}
\@namedef{tb@c@valkey@500000@watch@1000}{83.4333}
\@namedef{tb@p@valkey@500000@watch@1000}{29.7743}
\@namedef{tb@k@valkey@500000@watch@1000}{283117}
\@namedef{tb@c@valkey@500000@watch@100}{1240.9167}
\@namedef{tb@p@valkey@500000@watch@100}{1.752}
\@namedef{tb@k@valkey@500000@watch@100}{26993.6667}
\@namedef{tb@c@valkey@500000@watch@10}{3067.4767}
\@namedef{tb@p@valkey@500000@watch@10}{0.5735}
\@namedef{tb@k@valkey@500000@watch@10}{3964}
\@namedef{tb@c@valkey@500000@watch@1}{4680.55}
\@namedef{tb@p@valkey@500000@watch@1}{0.2462}
\@namedef{tb@k@valkey@500000@watch@1}{0}
\@namedef{tb@l@mariadb@1000000@indexes_on}{1.1932}
\@namedef{tb@l@mariadb@1000000@load}{2.3994}
\@namedef{tb@l@mariadb@1000000@reset}{0.0343}
\@namedef{tb@l@mongo@1000000@indexes_on}{2.985}
\@namedef{tb@l@mongo@1000000@load}{3.3123}
\@namedef{tb@l@mongo@1000000@reset}{0.1266}
\@namedef{tb@l@postgres@1000000@indexes_on}{0.7993}
\@namedef{tb@l@postgres@1000000@load}{2.0113}
\@namedef{tb@l@postgres@1000000@reset}{0.0293}
\@namedef{tb@l@valkey@1000000@indexes_on}{1.46}
\@namedef{tb@l@valkey@1000000@load}{0.7974}
\@namedef{tb@l@valkey@1000000@reset}{0.0005}
\@namedef{tb@l@mariadb@10000000@indexes_on}{17.0267}
\@namedef{tb@l@mariadb@10000000@load}{19.3588}
\@namedef{tb@l@mariadb@10000000@reset}{0.1015}
\@namedef{tb@l@mongo@10000000@indexes_on}{20.4731}
\@namedef{tb@l@mongo@10000000@load}{29.3856}
\@namedef{tb@l@mongo@10000000@reset}{0.1278}
\@namedef{tb@l@postgres@10000000@indexes_on}{4.2791}
\@namedef{tb@l@postgres@10000000@load}{12.3786}
\@namedef{tb@l@postgres@10000000@reset}{0.063}
\@namedef{tb@l@valkey@10000000@indexes_on}{15.4408}
\@namedef{tb@l@valkey@10000000@load}{7.5306}
\@namedef{tb@l@valkey@10000000@reset}{0.0028}
\@namedef{tb@l@mariadb@500000@indexes_on}{0.5973}
\@namedef{tb@l@mariadb@500000@load}{1.1366}
\@namedef{tb@l@mariadb@500000@reset}{0.0269}
\@namedef{tb@l@mongo@500000@indexes_on}{2.0365}
\@namedef{tb@l@mongo@500000@load}{1.4595}
\@namedef{tb@l@mongo@500000@reset}{0.1082}
\@namedef{tb@l@postgres@500000@indexes_on}{0.5139}
\@namedef{tb@l@postgres@500000@load}{0.9567}
\@namedef{tb@l@postgres@500000@reset}{0.0161}
\@namedef{tb@l@valkey@500000@indexes_on}{0.7012}
\@namedef{tb@l@valkey@500000@load}{0.3166}
\@namedef{tb@l@valkey@500000@reset}{0.0003}
\@namedef{tb@m@postgres@500000@off}{135.2661}
\@namedef{tb@m@postgres@500000@on}{174.8286}
\@namedef{tb@m@mariadb@500000@off}{103}
\@namedef{tb@m@mariadb@500000@on}{156.7031}
\@namedef{tb@m@mongo@500000@off}{58.2852}
\@namedef{tb@m@mongo@500000@on}{82.5664}
\@namedef{tb@m@valkey@500000@off}{165.9224}
\@namedef{tb@m@valkey@500000@on}{194.7192}
\@namedef{tb@m@postgres@1000000@off}{261.5552}
\@namedef{tb@m@postgres@1000000@on}{316.2271}
\@namedef{tb@m@mariadb@1000000@off}{197.3281}
\@namedef{tb@m@mariadb@1000000@on}{287.2969}
\@namedef{tb@m@mongo@1000000@off}{110.1758}
\@namedef{tb@m@mongo@1000000@on}{146.3008}
\@namedef{tb@m@valkey@1000000@off}{333.4226}
\@namedef{tb@m@valkey@1000000@on}{372.6064}
\@namedef{tb@m@postgres@10000000@off}{2507.9274}
\@namedef{tb@m@postgres@10000000@on}{2832.2114}
\@namedef{tb@m@mariadb@10000000@off}{1875.4688}
\@namedef{tb@m@mariadb@10000000@on}{2548.1719}
\@namedef{tb@m@mongo@10000000@off}{1187.0117}
\@namedef{tb@m@mongo@10000000@on}{1392.6484}
\@namedef{tb@m@valkey@10000000@off}{3427.1587}
\@namedef{tb@m@valkey@10000000@on}{3662.298}
\makeatother
\newcommand{\tbget}[1]{\ifcsname tb@#1\endcsname\csname tb@#1\endcsname\else nan\fi}
\newcommand{\tbnum}[1]{\ifcsname tb@#1\endcsname\num[round-mode=figures,round-precision=3]{\csname tb@#1\endcsname}\else --\fi}
\newcommand{\val}[4]{\tbnum{s@#1@#2@#3@#4}}
\newcommand{\valp}[4]{\tbget{s@#1@#2@#3@#4}}
\newcommand{\cval}[4]{\tbnum{c@#1@#2@#3@#4}}
\newcommand{\cpn}[4]{\tbnum{p@#1@#2@#3@#4}}
\newcommand{\ckn}[4]{\tbnum{k@#1@#2@#3@#4}}
\newcommand{\lval}[3]{\tbnum{l@#1@#2@#3}}
\newcommand{\sval}[3]{\tbnum{m@#1@#2@#3}}
\newcommand{\ratio}[5]{\ifcsname tb@s@#1@#3@#4@#5\endcsname\ifcsname tb@s@#2@#3@#4@#5\endcsname\num[round-mode=figures,round-precision=2]{\fpeval{\csname tb@s@#1@#3@#4@#5\endcsname/\csname tb@s@#2@#3@#4@#5\endcsname}}\else --\fi\else --\fi}
\newcommand{\cratio}[6]{\ifcsname tb@c@#1@#3@#4@#6\endcsname\ifcsname tb@c@#2@#3@#5@#6\endcsname\num[round-mode=figures,round-precision=2]{\fpeval{\csname tb@c@#1@#3@#4@#6\endcsname/\csname tb@c@#2@#3@#5@#6\endcsname}}\else --\fi\else --\fi}

% generated by tbench report
\makeatletter
\@namedef{tb@f@postgres@500000@off}{12}
\@namedef{tb@f@postgres@500000@on}{2}
\@namedef{tb@f@mariadb@500000@off}{11}
\@namedef{tb@f@mariadb@500000@on}{2}
\@namedef{tb@f@mongo@500000@off}{11}
\@namedef{tb@f@mongo@500000@on}{1}
\@namedef{tb@f@valkey@500000@off}{7}
\@namedef{tb@f@valkey@500000@on}{0}
\@namedef{tb@f@postgres@1000000@off}{12}
\@namedef{tb@f@postgres@1000000@on}{2}
\@namedef{tb@f@mariadb@1000000@off}{11}
\@namedef{tb@f@mariadb@1000000@on}{1}
\@namedef{tb@f@mongo@1000000@off}{11}
\@namedef{tb@f@mongo@1000000@on}{1}
\@namedef{tb@f@valkey@1000000@off}{7}
\@namedef{tb@f@valkey@1000000@on}{0}
\@namedef{tb@f@postgres@10000000@off}{12}
\@namedef{tb@f@postgres@10000000@on}{2}
\@namedef{tb@f@mariadb@10000000@off}{11}
\@namedef{tb@f@mariadb@10000000@on}{1}
\@namedef{tb@f@mongo@10000000@off}{11}
\@namedef{tb@f@mongo@10000000@on}{1}
\@namedef{tb@f@valkey@10000000@off}{7}
\@namedef{tb@f@valkey@10000000@on}{0}
\makeatother
\newcommand{\fullscans}[3]{\tbget{f@#1@#2@#3}}
\newcommand{\PlanScenarioCount}{13}

% generated by tbench report
\makeatletter
\@namedef{tb@sc@postgres@base@1}{4914.9950}
\@namedef{tb@sc@postgres@base@100}{8487.9317}
\@namedef{tb@sc@postgres@role@1}{5221.3983}
\@namedef{tb@sc@postgres@role@100}{7127.6750}
\@namedef{tb@sc@postgres@tls@1}{4539.4200}
\@namedef{tb@sc@postgres@tls@100}{6302.7650}
\@namedef{tb@sc@postgres@tls_role@1}{3776.1883}
\@namedef{tb@sc@postgres@tls_role@100}{7243.5483}
\@namedef{tb@sc@mariadb@base@1}{4381.2467}
\@namedef{tb@sc@mariadb@base@100}{4209.6567}
\@namedef{tb@sc@mariadb@role@1}{4439.5400}
\@namedef{tb@sc@mariadb@role@100}{5502.2133}
\@namedef{tb@sc@mariadb@tls@1}{3807.6667}
\@namedef{tb@sc@mariadb@tls@100}{4812.5883}
\@namedef{tb@sc@mariadb@tls_role@1}{4007.6683}
\@namedef{tb@sc@mariadb@tls_role@100}{4872.5533}
\@namedef{tb@sc@mongo@base@1}{718.4050}
\@namedef{tb@sc@mongo@base@100}{3127.4667}
\@namedef{tb@sc@mongo@role@1}{668.6217}
\@namedef{tb@sc@mongo@role@100}{2930.1250}
\@namedef{tb@sc@mongo@tls@1}{664.5650}
\@namedef{tb@sc@mongo@tls@100}{3153.3917}
\@namedef{tb@sc@mongo@tls_role@1}{656.0850}
\@namedef{tb@sc@mongo@tls_role@100}{2970.6667}
\@namedef{tb@sc@valkey@base@1}{17180.0500}
\@namedef{tb@sc@valkey@base@100}{36390.8200}
\@namedef{tb@sc@valkey@role@1}{16987.3367}
\@namedef{tb@sc@valkey@role@100}{33504.8500}
\@namedef{tb@sc@valkey@tls@1}{16249.5667}
\@namedef{tb@sc@valkey@tls@100}{19236.9367}
\@namedef{tb@sc@valkey@tls_role@1}{16946.2517}
\@namedef{tb@sc@valkey@tls_role@100}{18292.9567}
\@namedef{tb@sp@postgres@role@C}{1.0584}
\@namedef{tb@sp@postgres@role@R}{-34.8305}
\@namedef{tb@sp@postgres@role@U}{-27.4629}
\@namedef{tb@sp@postgres@role@D}{2.4578}
\@namedef{tb@sp@postgres@role@J}{-0.1960}
\@namedef{tb@sp@postgres@role@all}{-14.6989}
\@namedef{tb@sp@postgres@tls@C}{4.8214}
\@namedef{tb@sp@postgres@tls@R}{-23.8611}
\@namedef{tb@sp@postgres@tls@U}{-20.1724}
\@namedef{tb@sp@postgres@tls@D}{6.0524}
\@namedef{tb@sp@postgres@tls@J}{-2.8802}
\@namedef{tb@sp@postgres@tls@all}{-8.6409}
\@namedef{tb@sp@postgres@tls_role@C}{8.2396}
\@namedef{tb@sp@postgres@tls_role@R}{-21.8393}
\@namedef{tb@sp@postgres@tls_role@U}{-25.7224}
\@namedef{tb@sp@postgres@tls_role@D}{16.4665}
\@namedef{tb@sp@postgres@tls_role@J}{11.7918}
\@namedef{tb@sp@postgres@tls_role@all}{-5.5389}
\@namedef{tb@sp@mariadb@role@C}{-24.2090}
\@namedef{tb@sp@mariadb@role@R}{18.5775}
\@namedef{tb@sp@mariadb@role@U}{27.4489}
\@namedef{tb@sp@mariadb@role@D}{-8.0775}
\@namedef{tb@sp@mariadb@role@J}{49.0605}
\@namedef{tb@sp@mariadb@role@all}{5.7390}
\@namedef{tb@sp@mariadb@tls@C}{-13.0028}
\@namedef{tb@sp@mariadb@tls@R}{15.2403}
\@namedef{tb@sp@mariadb@tls@U}{8.9656}
\@namedef{tb@sp@mariadb@tls@D}{-7.1496}
\@namedef{tb@sp@mariadb@tls@J}{16.7473}
\@namedef{tb@sp@mariadb@tls@all}{2.0577}
\@namedef{tb@sp@mariadb@tls_role@C}{-23.2687}
\@namedef{tb@sp@mariadb@tls_role@R}{1.2071}
\@namedef{tb@sp@mariadb@tls_role@U}{-5.4440}
\@namedef{tb@sp@mariadb@tls_role@D}{-1.3829}
\@namedef{tb@sp@mariadb@tls_role@J}{33.9070}
\@namedef{tb@sp@mariadb@tls_role@all}{-3.8522}
\@namedef{tb@sp@mongo@role@C}{8.2929}
\@namedef{tb@sp@mongo@role@R}{15.2936}
\@namedef{tb@sp@mongo@role@U}{15.0717}
\@namedef{tb@sp@mongo@role@D}{-7.7066}
\@namedef{tb@sp@mongo@role@J}{21.9950}
\@namedef{tb@sp@mongo@role@all}{8.8493}
\@namedef{tb@sp@mongo@tls@C}{-6.0874}
\@namedef{tb@sp@mongo@tls@R}{9.9633}
\@namedef{tb@sp@mongo@tls@U}{1.5044}
\@namedef{tb@sp@mongo@tls@D}{-10.1766}
\@namedef{tb@sp@mongo@tls@J}{32.6424}
\@namedef{tb@sp@mongo@tls@all}{1.8168}
\@namedef{tb@sp@mongo@tls_role@C}{6.0988}
\@namedef{tb@sp@mongo@tls_role@R}{22.4236}
\@namedef{tb@sp@mongo@tls_role@U}{4.3769}
\@namedef{tb@sp@mongo@tls_role@D}{18.9430}
\@namedef{tb@sp@mongo@tls_role@J}{23.1675}
\@namedef{tb@sp@mongo@tls_role@all}{13.8074}
\@namedef{tb@sp@valkey@role@C}{-3.9790}
\@namedef{tb@sp@valkey@role@R}{-9.3428}
\@namedef{tb@sp@valkey@role@U}{7.8227}
\@namedef{tb@sp@valkey@role@D}{-1.4746}
\@namedef{tb@sp@valkey@role@J}{8.7341}
\@namedef{tb@sp@valkey@role@all}{-0.8047}
\@namedef{tb@sp@valkey@tls@C}{15.7305}
\@namedef{tb@sp@valkey@tls@R}{7.5345}
\@namedef{tb@sp@valkey@tls@U}{3.4351}
\@namedef{tb@sp@valkey@tls@D}{-2.2768}
\@namedef{tb@sp@valkey@tls@J}{1.0193}
\@namedef{tb@sp@valkey@tls@all}{5.3504}
\@namedef{tb@sp@valkey@tls_role@C}{5.2201}
\@namedef{tb@sp@valkey@tls_role@R}{16.4528}
\@namedef{tb@sp@valkey@tls_role@U}{16.3471}
\@namedef{tb@sp@valkey@tls_role@D}{6.8535}
\@namedef{tb@sp@valkey@tls_role@J}{9.3830}
\@namedef{tb@sp@valkey@tls_role@all}{10.9043}
\@namedef{tb@sn@postgres}{11.2522}
\@namedef{tb@sn@mariadb}{-3.1134}
\@namedef{tb@sn@mongo}{18.3584}
\@namedef{tb@sn@valkey}{-9.1572}
\makeatother
\newcommand{\secpct}[3]{\ifcsname tb@sp@#1@#2@#3\endcsname\num[round-mode=places,round-precision=1,print-implicit-plus]{\csname tb@sp@#1@#2@#3\endcsname}\,\%\else --\fi}
\newcommand{\seccval}[3]{\tbnum{sc@#1@#2@#3}}
\newcommand{\secnoise}[1]{\ifcsname tb@sn@#1\endcsname\num[round-mode=places,round-precision=1,print-implicit-plus]{\csname tb@sn@#1\endcsname}\,\%\else --\fi}


% Wstawia wygenerowany plik albo informację o braku danych.
\newcommand{\tbinput}[1]{\IfFileExists{#1}{\input{#1}}{\emph{brak danych: \texttt{\detokenize{#1}}}}}

% \fourlines{plik}{ylabel}{xmode}{ymode}{xlabel}
\newcommand{\fourlines}[5]{%
  \begin{axis}[tb, xmode=#3, ymode=#4, ylabel={#2}, xlabel={#5}]
    \addplot[pg] table[x index=0, y=postgres] {#1}; \addlegendentry{PostgreSQL}
    \addplot[maria] table[x index=0, y=mariadb] {#1}; \addlegendentry{MariaDB}
    \addplot[mongo] table[x index=0, y=mongo] {#1}; \addlegendentry{MongoDB}
    \addplot[valkey] table[x index=0, y=valkey] {#1}; \addlegendentry{Valkey}
  \end{axis}}

\begin{document}

\begin{titlepage}
  \centering
  {\large Politechnika Krakowska im.~Tadeusza Kościuszki\par}
  {\large Wydział Informatyki i Telekomunikacji\par}
  \vspace{3cm}
  {\Large Zaawansowane technologie baz danych\par}
  \vspace{1.2cm}
  {\huge\bfseries Porównanie wydajności relacyjnych\\[4pt] i nierelacyjnych SZBD\par}
  \vspace{0.6cm}
  {\Large System rezerwacji biletów: PostgreSQL, MariaDB, MongoDB i Valkey\par}
  \vfill
  {\large Wykonali: Łukasz Ryczko, Łukasz Pacholec\par}
  \vspace{0.4cm}
  {\large Studia II stopnia, niestacjonarne, semestr 1\par}
  \vspace{1.5cm}
  {\large Kraków, listopad 2026\par}
\end{titlepage}

\tableofcontents
\newpage

\section{Temat i cel}

Temat projektu to porównanie wydajności relacyjnych i nierelacyjnych systemów
zarządzania bazą danych na przykładzie systemu rezerwacji biletów na
wydarzenia. Badamy cztery silniki: PostgreSQL~17, MariaDB~11.8, MongoDB~8.0
oraz Valkey~8.1.

Celem jest zmierzenie, jak te silniki wykonują operacje CRUD na zbiorach
danych o rozmiarze od pół miliona do dziesięciu milionów biletów. Drugim celem
jest sprawdzenie, jak radzą sobie z problemem typowym dla tej domeny: wielu
klientów jednocześnie próbuje zarezerwować te same miejsca. Domena łączy
transakcje, blokady, dane półstrukturalne (metadane wydarzeń i biletów w~JSON)
i duży ruch równoległy, więc każdy z tych elementów wynika z~zadania, a~nie
został dopisany na potrzeby testu.

\begin{table}[H]
\centering
\caption{Badane SZBD i ich rola w projekcie.}
\begin{tabular}{llp{7.4cm}}
\toprule
SZBD & model danych & rola \\
\midrule
PostgreSQL 17 & relacyjny + JSONB & 12 tabel, transakcje, blokady wierszy przy rezerwacji miejsca \\
MariaDB 11.8 & relacyjny (InnoDB) + JSON & ten sam schemat, bezpośrednie porównanie dwóch silników relacyjnych \\
MongoDB 8.0 & dokumentowy & wydarzenie z osadzonym cennikiem, rezerwacja z osadzonymi biletami i płatnością \\
Valkey 8.1 & klucz--wartość w pamięci & zbiory wolnych miejsc, blokady miejsc z TTL, liczniki sprzedaży, rezerwacje jako JSON \\
\bottomrule
\end{tabular}
\end{table}

Hipotezy badawcze sformułowane we wstępnej koncepcji:

\begin{description}
  \item[H1] Wraz ze wzrostem liczby równoległych klientów rezerwujących te same
  miejsca przewaga wydajnościowa Valkey i MongoDB nad bazami relacyjnymi rośnie.
  \item[H2] (zapasowa) Różnice wydajności między silnikami rosną wraz ze
  wzrostem liczby rekordów.
\end{description}

W stosunku do wstępnej koncepcji zmieniliśmy technologię aplikacji testowej.
Zamiast skryptu w Pythonie powstała aplikacja w języku Go podzielona na
mikroserwisy: osobny serwis gRPC dla każdej bazy, serwis wykonujący plan
testów i klient wiersza poleceń. Każdą operację śledzimy przez
OpenTelemetry, a identyfikator śladu (\emph{trace id}) zapisujemy przy każdym
wyniku. Kod znajduje się w~pięciu publicznych repozytoriach: repozytorium spinającym \texttt{ticketbench} i~czterech repozytoriach komponentów~\cite{repo}.

\section{Charakterystyka wybranych SZBD}
\label{sec:szbd}

\subsection{PostgreSQL 17}
\label{sec:opis}
%<*opis-pg>

PostgreSQL to relacyjny SZBD rozwijany od 1996 roku na licencji PostgreSQL
(odmiana BSD/MIT). Każde połączenie obsługuje osobny proces serwera.
Współbieżność zapewnia MVCC: zmiana wiersza tworzy jego nową wersję, czytelnicy
nie blokują piszących, a stare wersje usuwa proces \texttt{VACUUM}. Trwałość
gwarantuje dziennik zapisu z~wyprzedzeniem (WAL). Domyślny poziom izolacji to
READ COMMITTED, dostępne są też REPEATABLE READ i~SERIALIZABLE (SSI)~\cite{pg}.
Typ \texttt{JSONB} przechowuje dokumenty JSON w~postaci binarnej z~operatorami
(\texttt{->>}, \texttt{@>}, \texttt{||}) i~indeksami GIN. Silnik rozszerza się
przez rozszerzenia (PostGIS, TimescaleDB, pgvector, pg\_partman) i~obce
tabele (FDW).

%</opis-pg>
\subsection{MariaDB 11.8}
%<*opis-maria>

MariaDB powstała w~2009 roku jako rozwidlenie MySQL i~jest rozwijana przez
MariaDB Foundation na licencji GPLv2. Architektura opiera się na wątkach
w~jednym procesie i~wymiennych silnikach składowania. Domyślny silnik InnoDB
ma MVCC, dziennik redo i~undo oraz klucz główny jako indeks klastrowany: wiersze
leżą w~liściach B-drzewa klucza głównego, a~indeksy drugorzędne wskazują na
klucz główny. Domyślny poziom izolacji to REPEATABLE READ z~blokadami
następnych kluczy (\emph{next-key locks})~\cite{maria}. Typ \texttt{JSON} jest
aliasem \texttt{LONGTEXT} z~kontrolą poprawności, a~indeksować można go przez
kolumny wirtualne. Inne silniki to m.in. ColumnStore do analityki, Aria
i~Spider do shardingu.

%</opis-maria>
\subsection{MongoDB 8.0}
%<*opis-mongo>

MongoDB to dokumentowy SZBD przechowujący dokumenty BSON w~kolekcjach, bez
narzuconego schematu (opcjonalnie z~walidacją JSON Schema). Silnik WiredTiger
stosuje MVCC na poziomie dokumentu, kompresję i~dziennik. Operacja na jednym
dokumencie jest zawsze atomowa. Od wersji 4.0 dostępne są transakcje
wielodokumentowe z~izolacją migawkową, ale wymagają replica setu lub
klastra~\cite{mongo}. Zapytania obsługują filtry, projekcje i~potok agregacji
(\texttt{\$match}, \texttt{\$group}, \texttt{\$lookup}). Indeksy są B-drzewami
na polach, w~tym zagnieżdżonych i~tablicowych (indeksy wielokluczowe).
Licencja serwera to SSPL, która ogranicza oferowanie MongoDB jako usługi
chmurowej przez strony trzecie.

%</opis-mongo>
\subsection{Valkey 8.1}
%<*opis-valkey>

Valkey to magazyn klucz--wartość w~pamięci operacyjnej. Powstał w~2024 roku
jako rozwidlenie Redis 7.2.4 pod opieką Linux Foundation, po zmianie licencji
Redisa, i~jest na licencji BSD. Wartości mają typy: napisy, hasze, listy,
zbiory, zbiory posortowane, strumienie. Polecenia wykonuje jeden wątek
główny, dlatego każde polecenie i~każdy skrypt Lua są atomowe, a~wątki I/O
obsługują tylko sieć~\cite{valkey}. Nie ma języka zapytań, planera ani
indeksów drugorzędnych. Dostęp po atrybucie wymaga struktur utrzymywanych
przez aplikację. Moduły (valkey-json, valkey-bloom, valkey-search) dodają
m.in. typ JSON i~wyszukiwanie.

%</opis-valkey>
\subsection{Zalety i wady}
\label{sec:zalety}

\begin{table}[H]
\centering
\small
\caption{Udogodnienia i ograniczenia.}
\label{tab:zalety}
%<*zalety>
\begin{tabular}{>{\raggedright\arraybackslash}p{1.9cm}>{\raggedright\arraybackslash}p{6.4cm}>{\raggedright\arraybackslash}p{6.4cm}}
\toprule
SZBD & zalety i udogodnienia & wady i ograniczenia \\
\midrule
PostgreSQL & zgodność ze standardem SQL, bogaty planer (złączenia haszowe, równoległe skany, indeksy częściowe i~na wyrażeniach), JSONB z~indeksami GIN, rozszerzenia, SSI & proces na połączenie (przy setkach sesji potrzebny pooler, np. PgBouncer), koszt \texttt{VACUUM} przy częstych aktualizacjach, brak wbudowanego shardingu \\
MariaDB & wysoka wydajność prostych zapytań po kluczu, indeks klastrowany, replikacja i~Galera Cluster w~pakiecie, zgodność z~ekosystemem MySQL & JSON tylko jako tekst, słabszy planer złożonych zapytań (brak złączeń haszowych dla dużych tabel w~domyślnej konfiguracji), blokady przedziałów przy REPEATABLE READ \\
MongoDB & elastyczny schemat, dokumenty zagnieżdżone zamiast złączeń, wbudowany sharding i~automatyczne przełączanie w~replica secie, potok agregacji & transakcje wielodokumentowe droższe niż w~bazach relacyjnych, duplikacja danych przy denormalizacji, limit 16\,MB na dokument, licencja SSPL \\
Valkey & opóźnienia rzędu mikrosekund, atomowe operacje na strukturach danych, TTL na kluczach, skrypty Lua, Pub/Sub, licencja BSD & dane muszą mieścić się w~RAM, brak zapytań po atrybutach bez ręcznie utrzymywanych indeksów, trwałość tylko przez RDB/AOF (możliwa utrata ostatnich zapisów), brak transakcji z~wycofaniem \\
\bottomrule
\end{tabular}
%</zalety>
\end{table}

\subsection{Awaryjność, bezpieczeństwo, migracje, integracje i skalowalność}
\label{sec:ops}

\begin{table}[H]
\centering
\footnotesize
\caption{Porównanie cech operacyjnych (część teoretyczna).}
\label{tab:ops}
%<*ops>
\begin{tabular}{>{\raggedright\arraybackslash}p{1.75cm}>{\raggedright\arraybackslash}p{3.45cm}>{\raggedright\arraybackslash}p{3.45cm}>{\raggedright\arraybackslash}p{3.45cm}>{\raggedright\arraybackslash}p{3.45cm}}
\toprule
cecha & PostgreSQL & MariaDB & MongoDB & Valkey \\
\midrule
awaryjność i~odtwarzanie & WAL, odtwarzanie do punktu w~czasie (PITR) z~kopii bazowej i~archiwum WAL; replikacja strumieniowa (synchroniczna lub asynchroniczna), przełączanie przez Patroni & dziennik redo InnoDB, binlog; replikacja asynchroniczna lub półsynchroniczna, Galera (synchroniczna, wiele węzłów zapisujących); kopie Mariabackup & replica set (zwykle 3 węzły), automatyczny wybór nowego węzła głównego w~kilka sekund, \emph{write concern} \texttt{majority} chroni przed utratą potwierdzonych zapisów & RDB (migawki) i/lub AOF (dziennik poleceń, \texttt{fsync} co sekundę lub przy każdym zapisie); replikacja asynchroniczna, Sentinel lub Cluster do przełączania \\
bezpieczeństwo & role i~\texttt{GRANT} do poziomu kolumny, Row-Level Security, SCRAM-SHA-256, TLS, \texttt{pgcrypto}; brak wbudowanego szyfrowania danych na dysku w~wersji społecznościowej & role i~\texttt{GRANT}, wtyczki uwierzytelniania (ed25519, PAM, GSSAPI), TLS, szyfrowanie przestrzeni tabel i~dzienników w~spoczynku (\emph{data-at-rest encryption}) & RBAC, SCRAM, x.509, LDAP; TLS; szyfrowanie w~spoczynku w~wersji Enterprise i~Atlas; szyfrowanie pól po stronie klienta (CSFLE) i~Queryable Encryption & ACL (użytkownicy, kategorie poleceń, wzorce kluczy), TLS; brak szyfrowania w~spoczynku, domyślnie bez hasła, więc konfiguracja jest niezbędna \\
migracje & \texttt{pg\_dump}/\texttt{pg\_restore}, \texttt{pg\_upgrade}, replikacja logiczna do aktualizacji bez przestoju; narzędzia schematu: Flyway, Liquibase & \texttt{mariadb-dump}, Mariabackup, prosta migracja z~MySQL (zgodny protokół i~składnia), Flyway, Liquibase & \texttt{mongodump}/\texttt{mongorestore}, \texttt{mongoimport}, Relational Migrator (z~baz SQL), migracje schematu w~kodzie aplikacji & kopiowanie pliku RDB, \texttt{MIGRATE}, replikacja z~Redisa 7.2 (zgodny protokół), więc migracja z~Redisa jest bezpośrednia \\
integracje & sterowniki dla wszystkich języków, FDW, replikacja logiczna i~Debezium (CDC), PostGIS, pgvector & sterowniki MySQL, Debezium, MaxScale (proxy, podział odczytów), CONNECT (obce źródła) & change streams, konektor Kafka, Atlas Search, sterowniki z~mapowaniem obiektów & klienty dla wszystkich języków, Pub/Sub, strumienie, zgodność z~bibliotekami Redisa \\
skalowalność & pionowa; odczyty przez repliki; partycjonowanie deklaratywne; sharding przez rozszerzenia (Citus) & pionowa; odczyty przez repliki; Galera dla wysokiej dostępności; Spider do shardingu & pozioma: wbudowany sharding (klucz haszowy lub zakresowy), balanser, repliki do odczytów & pozioma: Valkey Cluster (16\,384 sloty haszowe), repliki; ograniczenie: pamięć RAM węzłów \\
\bottomrule
\end{tabular}
%</ops>
\end{table}
%<*ops-uwaga>

W~naszym środowisku testowym wszystkie bazy działają jako pojedyncze węzły
(MongoDB jako jednowęzłowy replica set), bez TLS i~bez ról ograniczających
uprawnienia. Wyniki opisują więc wydajność silnika, a~nie konfiguracji
produkcyjnej z~replikacją i~szyfrowaniem.

%</ops-uwaga>
\subsection{Obszary zastosowań biznesowych}
\label{sec:zastosowania}
%<*zastosowania>

\begin{description}
  \item[PostgreSQL] systemy transakcyjne wymagające spójności i~złożonych
  zapytań: bankowość i~płatności, ERP, systemy rezerwacyjne, SaaS
  wielodostępne (RLS), systemy geoinformacyjne (PostGIS), hurtownie
  średniej wielkości.
  \item[MariaDB] aplikacje webowe i~e-commerce (stos LAMP, WordPress,
  Magento), systemy zastępujące MySQL, klastry wysokiej dostępności z~Galerą,
  analityka kolumnowa w~ColumnStore.
  \item[MongoDB] katalogi produktów i~treści o~zmiennej strukturze, profile
  użytkowników, backendy aplikacji mobilnych, IoT i~dane czasowe, systemy,
  w~których dokument odpowiada agregatowi domenowemu (zamówienie z~pozycjami).
  \item[Valkey] pamięć podręczna przed bazą główną, sesje, limity zapytań
  (\emph{rate limiting}), rankingi (zbiory posortowane), kolejki i~Pub/Sub,
  czasowe blokady zasobów z~TTL. W~systemie biletowym to dokładnie
  przytrzymanie miejsca na czas płatności.
\end{description}

W~praktyce system rezerwacji biletów łączy kilka silników: dane transakcyjne
(rezerwacje, płatności) w~bazie relacyjnej, katalog wydarzeń w~bazie
dokumentowej, a~dostępność miejsc, blokady i~kolejkę wejścia w~Valkey.
Wyniki testów w~kolejnych rozdziałach pokazują, gdzie każdy z~nich ma przewagę.
%</zastosowania>

\section{Model danych}

\subsection{Schemat relacyjny}

PostgreSQL i MariaDB używają tego samego schematu z~12 tabel: \texttt{kategorie},
\texttt{organizatorzy}, \texttt{obiekty}, \texttt{sektory}, \texttt{miejsca},
\texttt{wydarzenia}, \texttt{cenniki}, \texttt{klienci}, \texttt{kody\_rabatowe},
\texttt{rezerwacje}, \texttt{bilety} i \texttt{platnosci}.

Bilet oznacza jedno miejsce oferowane na jedno wydarzenie. Jest pulą
sprzedażową: przy imporcie powstaje bilet dla każdego miejsca obiektu, a jego
status przechodzi z~\texttt{wolny} do \texttt{zablokowany},
\texttt{zarezerwowany} albo \texttt{sprzedany}. Dzięki temu rezerwacja miejsca
jest warunkową aktualizacją jednego wiersza
(\texttt{UPDATE bilety SET status = 'zarezerwowany' ... WHERE id = ? AND status = 'wolny'}).
Ta sama operacja istnieje w~każdym z~czterech modeli, co pozwala porównać
silniki na identycznym problemie.

\begin{figure}[H]
\centering
\begin{tikzpicture}[
  tbl/.style={draw, rounded corners=2pt, fill=gray!8, font=\footnotesize\ttfamily, minimum width=2.3cm, minimum height=0.6cm},
  rel/.style={-{Stealth[length=2mm]}, gray!70!black}]
  \node[tbl] (obj) {obiekty};
  \node[tbl, right=1.1cm of obj] (sek) {sektory};
  \node[tbl, right=1.1cm of sek] (mie) {miejsca};
  \node[tbl, below=0.9cm of obj] (wyd) {wydarzenia};
  \node[tbl, left=0.8cm of wyd] (kat) {kategorie};
  \node[tbl, below=0.9cm of kat] (org) {organizatorzy};
  \node[tbl, below=0.9cm of wyd] (cen) {cenniki};
  \node[tbl, below=0.9cm of mie] (bil) {bilety};
  \node[tbl, below=0.9cm of bil] (rez) {rezerwacje};
  \node[tbl, right=0.8cm of rez] (pla) {platnosci};
  \node[tbl, below=0.9cm of rez] (kli) {klienci};
  \node[tbl, left=0.8cm of kli] (kod) {kody\_rabatowe};
  \draw[rel] (sek) -- (obj); \draw[rel] (mie) -- (sek);
  \draw[rel] (wyd) -- (obj); \draw[rel] (wyd) -- (kat); \draw[rel] (wyd) -- (org);
  \draw[rel] (cen) -- (wyd); \draw[rel] (cen.east) -| (sek.south);
  \draw[rel] (bil) -- (mie); \draw[rel] (bil.west) -- ++(-1.2,0) |- (wyd.east);
  \draw[rel] (bil) -- (rez); \draw[rel] (pla) -- (rez);
  \draw[rel] (rez) -- (kli); \draw[rel] (rez.west) -| (kod.north);
\end{tikzpicture}
\caption{Schemat relacyjny (strzałka: kolumna odwołująca się do klucza głównego).}
\end{figure}

Klucze obce zostały celowo pominięte. InnoDB automatycznie tworzy indeks dla
każdego klucza obcego, więc w~wariancie ,,bez indeksów'' MariaDB miałaby
kilka indeksów, których PostgreSQL by nie miał. Spójność referencyjną
zapewnia aplikacja, a~sprawdzają ją testy (rozdz.~\ref{sec:testy}).

Wariant ,,z~indeksami'' dodaje w~obu bazach te same 15 indeksów drugorzędnych.
Wśród nich są \texttt{bilety(wydarzenie\_id, status)},
\texttt{bilety(rezerwacja\_id)}, \texttt{rezerwacje(klient\_id, utworzono)}
i \texttt{rezerwacje(status, wygasa\_o)}. Na danych JSON indeksujemy pole
\texttt{metadane->>'gatunek'}: w~PostgreSQL przez indeks na wyrażeniu
i~indeks GIN \texttt{jsonb\_path\_ops}, w~MariaDB przez kolumnę wirtualną
\texttt{gatunek} z~indeksem B-drzewa.

\subsection{Model dokumentowy (MongoDB)}

\begin{itemize}
  \item \texttt{wydarzenia}: wydarzenie z~osadzonym cennikiem (tablica
  \texttt{cennik}) oraz zdenormalizowanymi polami \texttt{miasto},
  \texttt{obiekt\_nazwa} i \texttt{kategoria\_nazwa}, żeby wyszukiwanie po
  mieście nie wymagało \texttt{\$lookup};
  \item \texttt{obiekty}: obiekt z~osadzonymi sektorami;
  \item \texttt{rezerwacje}: rezerwacja z~osadzonymi migawkami biletów
  (\texttt{bilety}: id, wydarzenie, sektor, rząd, numer, cena) i~osadzoną
  płatnością (\texttt{platnosc});
  \item \texttt{bilety}: pula sprzedażowa jako osobna kolekcja, ponieważ to ją
  modyfikują równolegli klienci; osadzenie 2--5 tys. biletów w~dokumencie
  wydarzenia zamieniłoby rezerwację różnych miejsc w~konflikt o~jeden dokument;
  \item \texttt{klienci}, \texttt{organizatorzy}, \texttt{kategorie},
  \texttt{kody\_rabatowe}.
\end{itemize}

MongoDB działa jako jednowęzłowy \emph{replica set} \texttt{rs0}, bo
transakcje wielodokumentowe wymagają replikacji. Operacje zmieniające kilka
dokumentów (np. rezerwacja z~biletami, usunięcie konta klienta) wykonujemy
w~transakcjach, tak jak w~bazach relacyjnych.

\subsection{Model klucz--wartość (Valkey)}

\begin{center}
\small
\begin{tabular}{llp{8.2cm}}
\toprule
klucz & typ & zawartość \\
\midrule
\texttt{w:\{id\}} & hash & wydarzenie: nazwa, miasto, data, organizator, metadane (JSON), cennik (JSON) \\
\texttt{wb:\{wydarzenie\}} & hash & id biletu $\to$ JSON biletu (miejsce, sektor, cena, status, rezerwacja) \\
\texttt{wf:\{wydarzenie\}} & set & identyfikatory wolnych biletów \\
\texttt{r:\{id\}} & string & rezerwacja jako JSON z~biletami i~płatnością \\
\texttt{k:\{id\}} & hash & klient, preferencje jako JSON \\
\texttt{blok:\{bilet\}} & string z~TTL & czasowa blokada miejsca \\
\texttt{sprzedaz}, \texttt{przychod} & zset, hash & liczniki sprzedaży i~przychodu per wydarzenie \\
\texttt{ix:*} & set, zset & struktury indeksujące, tylko w~wariancie z~indeksami \\
\bottomrule
\end{tabular}
\end{center}

Valkey nie ma planera zapytań ani indeksów drugorzędnych. W~wariancie
,,z~indeksami'' aplikacja utrzymuje dodatkowe zbiory posortowane, np.
\texttt{ix:k:\{klient\}} (rezerwacje klienta według daty),
\texttt{ix:miasto:\{miasto\}} (wydarzenia według daty) czy \texttt{ix:oczek}
(rezerwacje według terminu wygaśnięcia). Bez nich zapytanie po atrybucie
wymaga \texttt{SCAN} po całej przestrzeni kluczy i~dekodowania każdego
dokumentu. Operacje zmieniające kilka kluczy są atomowe dzięki skryptom Lua.
Liczniki sprzedaży są częścią modelu bazowego, bo w~tej domenie są naturalną
strukturą, a~nie indeksem.

\subsection{Dane i rozmiary zbiorów}

Dane tworzy deterministyczny generator. Każda encja wynika z~trójki (ziarno,
rodzaj, identyfikator), więc każdy serwis generuje dane u~siebie i~wszystkie
cztery bazy dostają identyczne wiersze. Po imporcie serwis zwraca sumę
kontrolną FNV-64 wygenerowanych danych. Raport porównuje ją między bazami.

\begin{table}[H]
\centering
\caption{Rozmiary zbiorów (ziarno 42).}
\begin{tabular}{rrrrrr}
\toprule
bilety & wydarzenia & obiekty & miejsca & klienci & rezerwacje \\
\midrule
\num{500000} & 244 & 8 & \num{16684} & \num{62500} & \num{82807} \\
\num{1000000} & 476 & 16 & \num{33922} & \num{125000} & \num{170396} \\
\num{10000000} & \num{4156} & 166 & \num{396795} & \num{1250000} & \num{1674389} \\
\bottomrule
\end{tabular}
\end{table}

Około 60\% biletów jest sprzedanych albo zarezerwowanych, rezerwacja obejmuje
od 1 do 6 sąsiednich miejsc, 85\% rezerwacji jest opłaconych, a~10\% ma kod
rabatowy. Metadane wydarzeń (gatunek, wykonawcy, limit wiekowy, tagi),
biletów, klientów, rezerwacji i~płatności są zapisane jako JSON. Cztery
wydarzenia w~każdym zbiorze nie mają sprzedaży (,,piaskownica''). Korzystają
z~nich scenariusze rezerwujące miejsca i~test współbieżności, więc nigdy nie
kolidują z~danymi ani ze sobą nawzajem.

\section{Aplikacja testowa}
\label{sec:aplikacja}

\subsection{Wymagania}
\label{sec:wymagania}
%<*wymagania>

Wymagania funkcjonalne:
\begin{enumerate}[label=F\arabic*.]
  \item Generowanie tego samego syntetycznego zbioru danych o~zadanym rozmiarze
  (liczba biletów) i~ziarnie dla każdej bazy oraz potwierdzenie identyczności
  danych sumą kontrolną.
  \item Import hurtowy danych do PostgreSQL, MariaDB, MongoDB i~Valkey.
  \item Wykonanie każdego z~27 scenariuszy na każdej bazie z~zadaną liczbą
  powtórzeń, w~wariancie bez indeksów drugorzędnych i~z~nimi.
  \item Zapis planów zapytań dla każdego scenariusza i~wariantu indeksów.
  \item Test współbieżności: $N$ klientów rezerwujących te same miejsca, wiele
  strategii blokowania, kontrola braku podwójnej sprzedaży.
  \item Zapis surowych wyników do CSV i~wygenerowanie z~nich tabel, wykresów
  oraz oceny hipotez bez ręcznego przepisywania liczb.
  \item Uruchomienie całości jednym poleceniem.
\end{enumerate}

Wymagania niefunkcjonalne:
\begin{enumerate}[label=N\arabic*.]
  \item Powtarzalność: deterministyczne dane i~parametry, wersje obrazów
  i~konfiguracja baz zapisane w~repozytorium.
  \item Rzetelność pomiaru: czas mierzony w~serwisie bazy, bez narzutu sieci
  między CLI a~serwisem; rozgrzewka przed pomiarem; fikstury poza pomiarem.
  \item Obserwowalność: każdy wynik powiązany z~identyfikatorem śladu
  OpenTelemetry, metryki w~Prometheusie.
  \item Przenośność: uruchomienie na macOS i~Linuksie (amd64, arm64) w~Dockerze.
  \item Jakość: testy automatyczne na prawdziwych bazach w~CI przy każdej zmianie.
\end{enumerate}

%</wymagania>
\subsection{Technologie i narzędzia}
\label{sec:technologie}

\begin{table}[H]
\centering
\small
\caption{Wykorzystane technologie.}
%<*technologie>
\begin{tabular}{>{\raggedright\arraybackslash}p{3.6cm}>{\raggedright\arraybackslash}p{11cm}}
\toprule
obszar & technologie \\
\midrule
język i~komunikacja & Go 1.26, gRPC i~Protocol Buffers (buf), cobra (CLI) \\
sterowniki baz & pgx v5 (PostgreSQL), go-sql-driver/mysql (MariaDB), mongo-driver v2, valkey-go \\
obserwowalność & OpenTelemetry SDK, otelgrpc, otelpgx, otelsql, otelmongo, valkeyotel, OpenTelemetry Collector, Jaeger, Prometheus, Grafana \\
środowisko & Docker, Docker Compose, obrazy distroless, GHCR \\
testy & testing (Go), testcontainers-go, golangci-lint, ansible-lint, shellcheck \\
CI/CD i~wdrożenie & GitHub Actions, GoReleaser, Ansible, skrypt \texttt{ticketbench.sh} \\
raport & LaTeX (XeLaTeX, tectonic), pgfplots, siunitx; tabele generowane przez \texttt{tbench report} \\
\bottomrule
\end{tabular}
%</technologie>
\end{table}

\subsection{Architektura}
\label{sec:architektura}

Aplikacja \emph{ticketbench} jest napisana w~Go (1.26) i~podzielona na
mikroserwisy komunikujące się przez gRPC (rys.~\ref{fig:arch}). Kontrakty
(\texttt{StoreService}, \texttt{RunnerService}), model danych, generator,
definicje scenariuszy i~konfiguracja OpenTelemetry tworzą wspólny moduł
\texttt{ticketbench-contracts}.

\begin{figure}[H]
\centering
\begin{tikzpicture}[
  box/.style={draw, rounded corners=3pt, align=center, font=\footnotesize, minimum height=0.85cm, inner sep=4pt},
  svc/.style={box, fill=blue!6},
  db/.style={box, fill=gray!10},
  obs/.style={box, fill=orange!10},
  arr/.style={-{Stealth[length=2mm]}},
  otlp/.style={-{Stealth[length=2mm]}, dashed, gray}]
  \node[box, fill=green!8] (cli) {\texttt{tbench} CLI\\root span};
  \node[svc, right=1.3cm of cli] (run) {bench-runner\\plan, wyniki};
  \node[svc, right=1.6cm of run, yshift=1.9cm] (spg) {store-postgres};
  \node[svc, below=0.35cm of spg] (smy) {store-mariadb};
  \node[svc, below=0.35cm of smy] (smo) {store-mongo};
  \node[svc, below=0.35cm of smo] (sva) {store-valkey};
  \node[db, right=1.1cm of spg] (pg) {PostgreSQL 17};
  \node[db, right=1.1cm of smy] (my) {MariaDB 11.8};
  \node[db, right=1.1cm of smo] (mo) {MongoDB 8 (rs0)};
  \node[db, right=1.1cm of sva] (va) {Valkey 8.1};
  \node[obs, below=1.3cm of run] (col) {OTel Collector};
  \node[obs, left=0.7cm of col] (jae) {Jaeger};
  \node[obs, below=0.5cm of col, xshift=-1.2cm] (pro) {Prometheus};
  \node[obs, right=0.5cm of pro] (gra) {Grafana};
  \draw[arr] (cli) -- node[above, font=\scriptsize]{gRPC} node[below, font=\scriptsize]{traceparent} (run);
  \foreach \s in {spg,smy,smo,sva} \draw[arr] (run.east) -- (\s.west);
  \draw[arr] (spg) -- (pg); \draw[arr] (smy) -- (my); \draw[arr] (smo) -- (mo); \draw[arr] (sva) -- (va);
  \draw[otlp] (run) -- (col); \draw[otlp] (sva.south) |- (col.east);
  \draw[otlp] (col) -- (jae); \draw[otlp] (col) -- (pro); \draw[otlp] (pro) -- (gra);
\end{tikzpicture}
\caption{Architektura: linie ciągłe to wywołania gRPC i~połączenia z~bazą, przerywane to eksport OTLP.}
\label{fig:arch}
\end{figure}

%<*dzialanie>
\begin{description}
  \item[store-*] Jeden serwis na bazę, wspólny interfejs \texttt{Engine} w~Go.
  Serwis generuje dane z~ziarna, ładuje je hurtowo (PostgreSQL: \texttt{COPY}
  w~formacie binarnym, MariaDB: wielowierszowe \texttt{INSERT} po 1000 wierszy,
  MongoDB: \texttt{insertMany} po 5000 dokumentów, Valkey: potoki poleceń),
  tworzy i~usuwa indeksy, wykonuje scenariusze, zwraca plany zapytań
  i~przeprowadza test współbieżności.
  \item[bench-runner] Przyjmuje plan (listę kroków), wykonuje go po kolei
  i~strumieniuje wiersze wyników do klienta.
  \item[tbench] Klient CLI: buduje plan z~pliku \texttt{matrix.yaml}, zapisuje
  wyniki do CSV na bieżąco, generuje tabele i~dane wykresów do tego
  sprawozdania (\texttt{tbench report}).
\end{description}
%</dzialanie>

\subsection{Śledzenie i trace id}
\label{sec:sledzenie}
%<*sledzenie>

CLI rozpoczyna span główny dla każdego segmentu (jeden rozmiar zbioru na
jednej bazie). Kontekst śladu trafia do runnera jako nagłówek W3C
\texttt{traceparent} w~metadanych gRPC~\cite{w3c}. Runner tworzy spany kroków
i~wywołań \texttt{Execute}, przekazuje kontekst dalej do serwisu bazy, a~ten
do instrumentacji sterownika (otelpgx, otelsql, otelmongo, valkeyotel). Każde
polecenie SQL, operacja MongoDB i~polecenie Valkey jest więc spanem w~tym
samym śladzie. Logi (\texttt{slog}, JSON) zawierają pola \texttt{trace\_id}
i~\texttt{span\_id}. Identyfikator śladu jest zapisany w~każdym wierszu
plików CSV z~wynikami, więc z~dowolnego pomiaru można przejść do jego śladu
w~Jaegerze. Test jednostkowy runnera sprawdza, że nagłówek
\texttt{traceparent} z~identyfikatorem śladu klienta dociera do serwisu bazy.

Przy imporcie danych serwis oznacza kontekst jako niepróbkowany
(\texttt{telemetry.Quiet}): identyfikator śladu zostaje, ale tysiące spanów
pojedynczych partii \texttt{INSERT} nie są zapisywane. W~teście
współbieżności śledzimy 1\% prób rezerwacji. Metryki (histogramy czasu
scenariuszy i~prób rezerwacji, licznik konfliktów) trafiają przez kolektor do
Prometheusa i~na pulpit Grafany.

%</sledzenie>
\subsection{Automatyzacja}
\label{sec:testy}

%<*automatyzacja>
\begin{itemize}
  \item \textbf{Docker Compose} uruchamia 4 bazy, 4 serwisy, runner, kolektor
  OTel, Jaeger, Prometheus i~Grafanę. Obrazy serwisów są oparte na
  \texttt{distroless/static} (ok.~30\,MB) i~mają kontrolę zdrowia przez
  gRPC Health.
  \item \textbf{Jeden program} (\texttt{ticketbench.sh}) pobiera wszystkie
  serwisy z~GitHuba (repozytoria, obrazy, binarkę \texttt{tbench}), uruchamia
  środowisko i~całą macierz testów, generuje tabele i~wykresy, a~na koniec
  kompiluje to sprawozdanie i~prezentację. Instrukcja: runbook
  \texttt{docs/runbook.html}.
  \item \textbf{Testy automatyczne}: testy jednostkowe generatora (determinizm,
  spójność rezerwacji i~biletów), scenariuszy, generatora raportu i~CLI,
  a~także test zgodności (\texttt{enginetest}) uruchamiany na prawdziwych
  bazach przez testcontainers-go. Test zgodności porównuje wyniki zapytań
  z~wyrocznią liczoną bezpośrednio z~generatora, wykonuje wszystkie
  scenariusze w~obu wariantach indeksów, plany zapytań i~test współbieżności
  każdą strategią.
  \item \textbf{GitHub Actions}: lint (golangci-lint, buf, ansible-lint),
  testy z~detektorem wyścigów, test zgodności jako macierz czterech baz,
  benchmarki Go, binarki dla linux/amd64, linux/arm64 i~darwin/arm64,
  wieloarchitekturowe obrazy w~GHCR, wydania przez GoReleaser oraz test
  end-to-end, w~którym Ansible stawia środowisko na maszynie CI i~uruchamia
  macierz \texttt{smoke}. Sprawozdanie również kompiluje się w~CI.
  \item \textbf{Wariant bezpieczeństwa}: nakładka
  \texttt{docker-compose.security.yml} (certyfikaty z~własnego CA, konta
  aplikacyjne) i~polecenie \texttt{./ticketbench.sh security}.
  \item \textbf{Ansible}: role \texttt{docker} (instalacja silnika na
  Debianie/Ubuntu), \texttt{ticketbench} (wdrożenie i~start stosu),
  \texttt{bench} (uruchomienie macierzy i~pobranie wyników)
  i~\texttt{teardown}.
\end{itemize}
%</automatyzacja>

\section{Metodyka}

\subsection{Scenariusze}

Zgodnie z~koncepcją mamy 24 scenariusze CRUD, po 6 na każdą operację.
Dodaliśmy do nich 3 scenariusze na danych JSON (tab.~\ref{tab:scen}).
Parametry każdego wykonania (identyfikatory, miasta, przedziały dat) wynikają
wyłącznie z~(zbioru, scenariusza, numeru iteracji), więc wszystkie bazy
wykonują dokładnie te same operacje na tych samych wierszach.

\begin{table}[H]
\centering
\small
\caption{Scenariusze. ,,F'' oznacza scenariusz z~fiksturą: dane, na których
działa operacja, są wstawiane przed pomiarem i~nie wliczają się do czasu.}
\label{tab:scen}
\begin{tabular}{lp{4.6cm}p{8.6cm}}
\toprule
ID & nazwa & co dokładnie jest mierzone \\
\midrule
C1 & Pojedyncza rezerwacja & transakcja: nowa rezerwacja + zajęcie 1 wolnego biletu \\
C2 & Rezerwacja z~kilkoma biletami & jak C1, 4 bilety \\
C3 (F) & Masowy import biletów & \num{10000} nowych biletów (pula sprzedażowa nowych wydarzeń) \\
C4 & Rejestracja klienta & wstawienie klienta z~preferencjami JSON \\
C5 & Wydarzenie z~cennikiem & wydarzenie i~cena dla każdego sektora obiektu \\
C6 (F) & Zapis płatności & płatność + zmiana statusu rezerwacji na opłaconą \\
\midrule
R1 & Rezerwacja po identyfikatorze & rezerwacja z~biletami, miejscami i~sektorami \\
R2 & Wolne miejsca na wydarzenie & wolne bilety wydarzenia z~sektorem, rzędem i~numerem, posortowane \\
R3 & Wydarzenia po dacie i~mieście & wydarzenia w~mieście w~30-dniowym oknie \\
R4 & Historia rezerwacji klienta & 50 ostatnich rezerwacji klienta z~liczbą biletów \\
R5 & Sprzedaż per wydarzenie & liczba i~suma cen sprzedanych biletów dla wydarzeń organizatora \\
R6 & 10 najpopularniejszych wydarzeń & agregacja po wszystkich sprzedanych biletach \\
\midrule
U1 & Zmiana statusu rezerwacji & rezerwacja $\to$ opłacona \\
U2 & Zajęcie miejsca (blokada) & warunkowe \texttt{wolny}$\to$\texttt{zablokowany} z~terminem blokady \\
U3 & Masowa zmiana cen w~sektorze & cena $\times 1{,}1$ dla wolnych biletów sektora + cennik \\
U4 & Zmiana danych klienta & e-mail, telefon, scalenie pól w~preferencjach JSON \\
U5 & Zmiana terminu wydarzenia & przesunięcie daty o~tydzień \\
U6 (F) & Wygaszenie nieopłaconych rezerwacji & \num{1000} przeterminowanych rezerwacji $\to$ wygasłe, zwolnienie ich biletów \\
\midrule
D1 & Anulowanie rezerwacji & usunięcie rezerwacji i~płatności, zwolnienie biletów \\
D2 & Rezerwacja wraz z~biletami & usunięcie rezerwacji, płatności i~jej biletów \\
D3 (F) & Masowe usunięcie wygasłych & usunięcie \num{1000} wygasłych rezerwacji \\
D4 & Usunięcie konta klienta & klient, jego rezerwacje i~płatności, zwolnienie biletów \\
D5 (F) & Usunięcie odwołanego wydarzenia & wydarzenie, jego bilety (całe obiekty) i~cennik \\
D6 (F) & Usunięcie danych starszych niż data & wydarzenia sprzed 2 lat z~biletami, cennikami, rezerwacjami i~płatnościami \\
\midrule
J1 & Filtr po polu JSON & wydarzenia o~danym gatunku z~danym tagiem w~metadanych \\
J2 & Agregacja po polu JSON & liczba i~suma rezerwacji z~7 dni pogrupowana po \texttt{metadane.kanal} \\
J3 & Modyfikacja tablicy w~JSON & dopisanie tagu do \texttt{metadane.tagi} \\
\bottomrule
\end{tabular}
\end{table}

\subsection{Przebieg pomiaru}

Dla każdego rozmiaru zbioru i~każdej bazy plan wygląda tak:
\begin{enumerate}
  \item reset schematu, import danych, pomiar rozmiaru;
  \item wszystkie scenariusze bez indeksów drugorzędnych: 1 przebieg rozgrzewkowy
  i~5 mierzonych; plany zapytań;
  \item utworzenie indeksów (czas tworzenia jest osobnym wynikiem), pomiar rozmiaru;
  \item wszystkie scenariusze z~indeksami, plany zapytań;
  \item test współbieżności dla każdej strategii rezerwacji;
  \item reset (zwolnienie miejsca).
\end{enumerate}

Czas jest mierzony w~serwisie bazy, wokół samej operacji na bazie
(\texttt{time.Now()} przed i~po wywołaniu \texttt{Engine.Run}). Fikstura
(np. \num{1000} rezerwacji do usunięcia w~D3) jest wstawiana przed
rozpoczęciem pomiaru. Runner mierzy niezależnie czas całego wywołania gRPC,
dzięki czemu narzut sieci i~serializacji jest znany
(rozdz.~\ref{sec:grpc}), ale nie wpływa na porównanie baz. Wynik scenariusza
to średnia z~5 powtórzeń. W~plikach z~danymi są też mediana, odchylenie
standardowe, p95, minimum i~maksimum. Jedno wykonanie ma limit 60~s. Po
przekroczeniu limitu dalsze powtórzenia są pomijane, a~w~tabeli pojawia się
,,t/o''.

Kolejne powtórzenia działają na różnych wierszach (inny numer iteracji),
a~przebieg z~indeksami używa innego zakresu iteracji niż przebieg bez
indeksów. Scenariusze rezerwujące miejsca nie trafiają więc na bilety zajęte
we wcześniejszym przebiegu. Plany zapytań powstają przez
\texttt{EXPLAIN (ANALYZE, BUFFERS)} w~PostgreSQL, \texttt{ANALYZE FORMAT=JSON}
w~MariaDB (oba w~transakcji wycofywanej po pomiarze) oraz
\texttt{explain} z~\texttt{executionStats} w~MongoDB. Dla Valkey zapisujemy
listę wykonywanych poleceń.

\subsection{Test współbieżności}

Test odpowiada na pytanie z~hipotezy H1. W~serwisie bazy startuje
$N \in \{1, 10, 100, 1000\}$ gorutyn, każda z~własną sesją bazy. Wszystkie
próbują zarezerwować te same 500 biletów jednego wydarzenia, każda w~innej
losowej kolejności. Runda kończy się, gdy wszystkie bilety są zajęte. Mierzymy
czas rundy, przepustowość (udane rezerwacje na sekundę), liczbę konfliktów
(próba trafiła na zajęte miejsce albo transakcja została przerwana)
i~percentyle czasu pojedynczej próby. Po każdej rundzie serwis sprawdza
w~bazie, że istnieje dokładnie 500 rezerwacji i~500 zajętych biletów, czyli
że żadne miejsce nie zostało sprzedane dwa razy. Każdą konfigurację powtarzamy
3 razy.

Strategie rezerwacji:
\begin{itemize}
  \item PostgreSQL i~MariaDB: \texttt{cond} (warunkowy \texttt{UPDATE ... WHERE status='wolny'}),
  \texttt{forupdate} (\texttt{SELECT ... FOR UPDATE}, sprawdzenie, \texttt{UPDATE}),
  \texttt{nowait} (jak wyżej, ale zablokowany wiersz od razu oznacza konflikt),
  \texttt{serializable} (izolacja SERIALIZABLE i~ponawianie po błędzie serializacji lub zakleszczeniu);
  \item MongoDB: \texttt{atomic} (warunkowy \texttt{updateOne} na dokumencie biletu, potem
  \texttt{insertOne} rezerwacji) i~\texttt{txn} (obie operacje w~transakcji, \texttt{WriteConflict} to konflikt);
  \item Valkey: \texttt{lua} (sprawdzenie i~zapis w~jednym skrypcie),
  \texttt{setnx} (klucz blokady \texttt{SET NX PX}, czyli ,,blokada miejsca z~TTL''
  z~koncepcji, potem \texttt{MULTI}) i~\texttt{watch} (optymistyczne
  \texttt{WATCH}/\texttt{MULTI}/\texttt{EXEC} na zbiorze wolnych miejsc).
\end{itemize}
Pierwsza strategia z~listy jest domyślna i~to ona jest używana do weryfikacji H1.

\subsection{Środowisko}

\begin{itemize}
  \item Komputer: Apple M5 (10 rdzeni), 24\,GB RAM, macOS 27, Docker Desktop
  (maszyna wirtualna Linux 6.12, 10 vCPU, 8\,GB RAM).
  \item Każda z~trzech baz dyskowych dostaje 1\,GB pamięci podręcznej
  (\texttt{shared\_buffers}, \texttt{innodb\_buffer\_pool\_size},
  \texttt{wiredTigerCacheSizeGB}). Valkey trzyma cały zbiór w~RAM, bez RDB i~AOF.
  \item Trwałość zostaje domyślna: \texttt{fsync} i~\texttt{synchronous\_commit}
  w~PostgreSQL, \texttt{innodb\_flush\_log\_at\_trx\_commit=1} w~MariaDB,
  \emph{write concern} \texttt{majority} z~dziennikiem w~MongoDB. Valkey
  z~definicji nie zapisuje danych na dysk, co trzeba brać pod uwagę przy
  porównaniu zapisów.
  \item Poziomy izolacji są domyślne dla silnika: READ COMMITTED w~PostgreSQL,
  REPEATABLE READ w~MariaDB, \emph{snapshot} w~transakcjach MongoDB.
  \item Pule połączeń do 1100 sesji, \texttt{max\_connections=1200}.
  \item Przy zbiorach 1 i~10 mln biletów, na początku każdego segmentu,
  restartowane są kontenery pozostałych trzech baz. Mierzona baza ma wtedy do
  dyspozycji całą pamięć maszyny wirtualnej, a~pamięć podręczna innych silników
  jej nie zabiera.
\end{itemize}

\section{Wyniki operacji CRUD}

Macierz obejmuje \SizesList{} biletów i~dała \ScenarioRuns{} pomiarów
scenariuszy oraz \ContentionRuns{} rund testu współbieżności. Wszystkie czasy
w~tym rozdziale są czasami po stronie serwisu bazy, w~milisekundach, uśrednionymi
z~5 powtórzeń.

\subsection{Import danych i rozmiar}

\begin{table}[H]
\centering
\caption{Czas importu i~tworzenia indeksów drugorzędnych [s].}
\label{tab:load}
\small
\tbinput{generated/tables/load.tex}
\end{table}

\begin{table}[H]
\centering
\caption{Zajętość danych [MB]. Dla baz dyskowych jest to rozmiar plików bazy,
dla Valkey zużycie pamięci (\texttt{used\_memory}).}
\label{tab:storage}
\small
\tbinput{generated/tables/storage.tex}
\end{table}

Import 10 mln biletów najszybciej wykonuje Valkey
(\lval{valkey}{10000000}{load}\,s), potem PostgreSQL
(\lval{postgres}{10000000}{load}\,s, binarny \texttt{COPY}), MariaDB
(\lval{mariadb}{10000000}{load}\,s, wielowierszowe \texttt{INSERT}) i~MongoDB
(\lval{mongo}{10000000}{load}\,s). MongoDB płaci za \emph{write concern}
\texttt{majority} z~dziennikiem i~za budowę dokumentów BSON z~osadzonymi
tablicami. Czas importu rośnie z~rozmiarem w~przybliżeniu liniowo we
wszystkich bazach.

Tworzenie 15 indeksów na 10 mln biletów trwa w~PostgreSQL
\lval{postgres}{10000000}{indexes_on}\,s, w~MariaDB
\lval{mariadb}{10000000}{indexes_on}\,s, a~w~MongoDB
\lval{mongo}{10000000}{indexes_on}\,s. W~Valkey ,,indeksy'' buduje aplikacja:
musi przeczytać przez \texttt{SCAN} wszystkie 1,67 mln rezerwacji i~zapisać
zbiory posortowane, co zajmuje \lval{valkey}{10000000}{indexes_on}\,s.

Najmniej miejsca zajmuje MongoDB dzięki kompresji WiredTiger (Snappy):
\sval{mongo}{10000000}{off}\,MB bez indeksów wobec
\sval{postgres}{10000000}{off}\,MB w~PostgreSQL i~\sval{mariadb}{10000000}{off}\,MB
w~MariaDB. Valkey przy 10 mln biletów zajmuje \sval{valkey}{10000000}{off}\,MB
pamięci RAM. Cały zbiór musi się w~niej zmieścić, co jest praktyczną granicą
tego modelu. Wartości ,,z~indeksami'' obejmują też wiersze dopisane przez
scenariusze z~pierwszego przebiegu (np. C3 importuje \num{10000} biletów
w~każdym powtórzeniu).


\subsection{Czasy scenariuszy}

Tabele~\ref{tab:s500}--\ref{tab:s10m} zawierają średni czas każdego
scenariusza. Pogrubiona jest najlepsza wartość w~danym wariancie indeksów.
,,t/o'' oznacza przekroczenie limitu 60~s.

\newcommand{\scentable}[3]{%
\begin{table}[htbp]
\centering
\caption{Średni czas scenariuszy [ms], #2 biletów.}
\label{#3}
\resizebox{\linewidth}{!}{\tbinput{generated/tables/scenarios_#1.tex}}
\end{table}}
\scentable{500000}{500 tys.}{tab:s500}
\scentable{1000000}{1 mln}{tab:s1m}
\scentable{10000000}{10 mln}{tab:s10m}

\newcommand{\scenplot}[2]{%
  \nextgroupplot[title={#1 #2}]
  \addplot[pg] table[x=tickets, y=postgres] {generated/plots/scen_#1_#2.dat};
  \addplot[maria] table[x=tickets, y=mariadb] {generated/plots/scen_#1_#2.dat};
  \addplot[mongo] table[x=tickets, y=mongo] {generated/plots/scen_#1_#2.dat};
  \addplot[valkey] table[x=tickets, y=valkey] {generated/plots/scen_#1_#2.dat};}

\begin{figure}[p]
\centering
\begin{tikzpicture}
\begin{groupplot}[group style={group size=3 by 4, horizontal sep=1.3cm, vertical sep=1.5cm},
  width=0.36\linewidth, height=4.4cm, xmode=log, ymode=log, grid=major, grid style={gray!25},
  tick label style={font=\scriptsize}, title style={font=\footnotesize}, unbounded coords=jump, log ticks with fixed point, every axis plot/.append style={mark options={scale=0.7}},
  xtick={500000,1000000,10000000}, xticklabels={0{,}5,1,10}]
  \scenplot{R1}{off}
  \scenplot{R2}{off}
  \scenplot{R4}{off}
  \scenplot{R1}{on}
  \scenplot{R2}{on}
  \scenplot{R4}{on}
  \scenplot{R6}{off}
  \scenplot{U6}{off}
  \scenplot{D4}{off}
  \scenplot{R6}{on}
  \scenplot{U6}{on}
  \scenplot{D4}{on}
\end{groupplot}
\end{tikzpicture}

\vspace{2mm}
\footnotesize
\textcolor{cpg}{$\bullet$}~PostgreSQL\quad
\textcolor{cmaria}{$\blacksquare$}~MariaDB\quad
\textcolor{cmongo}{$\blacktriangle$}~MongoDB\quad
\textcolor{cvalkey}{$\blacklozenge$}~Valkey
\caption{Czas wybranych scenariuszy [ms] w~funkcji liczby biletów w~mln (obie osie
logarytmiczne). ,,off'': bez indeksów drugorzędnych, ,,on'': z~indeksami.}
\label{fig:scen}
\end{figure}

\paragraph{Operacje punktowe.} Wstawienie i~zmiana pojedynczego rekordu
po kluczu (C1, C2, C4--C6, U1, U2, U4, U5, J3) we wszystkich bazach trwają
poniżej 2\,ms i~nie zależą od rozmiaru zbioru. Pojedyncza rezerwacja (C1) na
10 mln biletów trwa \val{valkey}{10000000}{off}{C1}\,ms w~Valkey,
\val{postgres}{10000000}{off}{C1}\,ms w~PostgreSQL,
\val{mariadb}{10000000}{off}{C1}\,ms w~MariaDB i~\val{mongo}{10000000}{off}{C1}\,ms
w~MongoDB. MongoDB jest tu najwolniejszy, bo każdy zapis czeka na
potwierdzenie z~dziennika (\texttt{majority}), a~C1 dodatkowo jest
transakcją wielodokumentową. Valkey nie zapisuje niczego na dysk, więc jego
przewaga w~zapisach wynika częściowo z~innego poziomu trwałości.

\paragraph{Zapytania bez indeksów rosną liniowo.} W~bazach relacyjnych
i~w~MongoDB każde zapytanie po kolumnie innej niż klucz główny wymaga pełnego
skanu. Wolne miejsca na wydarzenie (R2) w~PostgreSQL rosną z~\val{postgres}{500000}{off}{R2}\,ms
(500 tys.) do \val{postgres}{10000000}{off}{R2}\,ms (10 mln). W~MariaDB rosną
z~\val{mariadb}{500000}{off}{R2} do \val{mariadb}{10000000}{off}{R2}\,ms,
a~w~MongoDB z~\val{mongo}{500000}{off}{R2} do \val{mongo}{10000000}{off}{R2}\,ms.
W~Valkey ten sam scenariusz trwa stale około \val{valkey}{10000000}{off}{R2}\,ms,
bo zbiór wolnych miejsc wydarzenia jest podstawową strukturą modelu, a~nie
indeksem.

\paragraph{MariaDB bez indeksów jest najwolniejsza z~relacyjnych.} Przy 10 mln
biletów MariaDB wykonuje R1 w~\val{mariadb}{10000000}{off}{R1}\,ms, a~PostgreSQL
w~\val{postgres}{10000000}{off}{R1}\,ms. Złączenie rezerwacji z~biletami bez
indeksu MariaDB realizuje pętlą zagnieżdżoną z~buforem bloków (\texttt{ALL} na
tabeli \texttt{bilety}, ok. 10 mln przeczytanych wierszy, tab.~\ref{tab:plans-rel}).
PostgreSQL używa złączenia haszowego i~równoległego skanu. Ta sama różnica
widać w~R5, U3, U6 i~D1--D6.

\paragraph{Osadzanie zamiast złączeń.} Odczyt rezerwacji z~biletami (R1)
w~MongoDB trwa \val{mongo}{10000000}{off}{R1}\,ms, a~w~Valkey
\val{valkey}{10000000}{off}{R1}\,ms, także bez indeksów drugorzędnych.
Bilety są osadzone w~dokumencie albo w~JSON-ie rezerwacji, więc wystarcza
wyszukiwanie po kluczu głównym. Ten sam efekt widać przy anulowaniu
rezerwacji (D1): MongoDB zwalnia bilety po identyfikatorach z~osadzonej
migawki, przez indeks \texttt{\_id}, w~\val{mongo}{10000000}{off}{D1}\,ms.
PostgreSQL bez indeksu na \texttt{bilety.rezerwacja\_id} potrzebuje do tego
\val{postgres}{10000000}{off}{D1}\,ms, a~MariaDB \val{mariadb}{10000000}{off}{D1}\,ms.
To przykład kompromisu normalizacji i~denormalizacji: odczyt agregatu jest
tani, ale dane biletów są przechowywane w~dwóch miejscach.

\paragraph{Agregacje.} Ranking 10 najpopularniejszych wydarzeń (R6) wymaga
w~bazach relacyjnych i~w~MongoDB przejścia przez wszystkie sprzedane bilety.
PostgreSQL z~indeksem robi to w~\val{postgres}{10000000}{on}{R6}\,ms (Index
Only Scan), MariaDB w~\val{mariadb}{10000000}{on}{R6}\,ms, MongoDB
w~\val{mongo}{10000000}{on}{R6}\,ms. Valkey odczytuje gotowy licznik ze zbioru
posortowanego w~\val{valkey}{10000000}{on}{R6}\,ms. Utrzymuje go przy każdej
sprzedaży, czyli przenosi koszt agregacji na zapis.

\paragraph{Valkey bez struktur indeksujących.} Zapytania po atrybucie
(historia klienta R4, wygaszanie U6, usuwanie D3, D4, D6, agregacja J2)
wymagają w~Valkey przejścia \texttt{SCAN} przez wszystkie klucze rezerwacji
i~zdekodowania każdego JSON-a. Przy 10 mln biletów to około
\val{valkey}{10000000}{off}{R4}\,ms na operację, najgorszy wynik w~całym badaniu.
Po dodaniu zbiorów \texttt{ix:*} ta sama historia klienta trwa
\val{valkey}{10000000}{on}{R4}\,ms. W~modelu klucz--wartość brak ,,indeksu''
nie oznacza wolniejszego planu, tylko brak jakiejkolwiek ścieżki dostępu.

\paragraph{Zmienność.} Mediany, odchylenia standardowe i~p95 są w~pliku
\texttt{generated/data/scenarios\_summary.csv}. Najwyższy współczynnik
zmienności mają operacje punktowe trwające ułamek milisekundy. Pojedyncze
powtórzenie trafiające na zimną pamięć podręczną albo punkt kontrolny WAL
wydłuża się wtedy kilkukrotnie. Dla operacji trwających powyżej 10\,ms
rozrzut jest mały i~nie zmienia kolejności baz.


\subsection{Wpływ indeksów}

\begin{figure}[H]
\centering
\begin{tikzpicture}
\begin{axis}[tb, ybar, bar width=1.6pt, ymode=log, log origin=infty, height=6.5cm,
  ylabel={przyspieszenie (bez / z~indeksami)}, xtick={0,...,26},
  xticklabels={C1,C2,C3,C4,C5,C6,R1,R2,R3,R4,R5,R6,U1,U2,U3,U4,U5,U6,D1,D2,D3,D4,D5,D6,J1,J2,J3},
  x tick label style={font=\scriptsize, rotate=90}, enlarge x limits=0.02,
  extra y ticks={1}, extra y tick style={grid=major, grid style={black, thick}}]
  \addplot[fill=cpg, draw=none] table[x=i, y=postgres] {generated/plots/speedup_10000000.dat}; \addlegendentry{PostgreSQL}
  \addplot[fill=cmaria, draw=none] table[x=i, y=mariadb] {generated/plots/speedup_10000000.dat}; \addlegendentry{MariaDB}
  \addplot[fill=cmongo, draw=none] table[x=i, y=mongo] {generated/plots/speedup_10000000.dat}; \addlegendentry{MongoDB}
  \addplot[fill=cvalkey, draw=none] table[x=i, y=valkey] {generated/plots/speedup_10000000.dat}; \addlegendentry{Valkey}
\end{axis}
\end{tikzpicture}
\caption{Ile razy scenariusz jest szybszy z~indeksami, 10 mln biletów (skala
logarytmiczna, wartość poniżej 1 oznacza spowolnienie). Brak słupka: brak
pomiaru bez indeksów (np. przekroczony limit czasu).}
\label{fig:speedup}
\end{figure}

Indeksy przyspieszają przede wszystkim scenariusze filtrujące po kolumnie
obcej. Przy 10 mln biletów historia klienta (R4) w~PostgreSQL skraca się
ze \val{postgres}{10000000}{off}{R4} do \val{postgres}{10000000}{on}{R4}\,ms,
w~MariaDB z~\val{mariadb}{10000000}{off}{R4} do \val{mariadb}{10000000}{on}{R4}\,ms,
a~w~MongoDB z~\val{mongo}{10000000}{off}{R4} do \val{mongo}{10000000}{on}{R4}\,ms.
Wygaszanie rezerwacji (U6) w~MariaDB spada z~\val{mariadb}{10000000}{off}{U6}
do \val{mariadb}{10000000}{on}{U6}\,ms.

Koszt indeksów widać przy wstawianiu. Import \num{10000} biletów (C3)
w~PostgreSQL wydłuża się z~\val{postgres}{10000000}{off}{C3} do
\val{postgres}{10000000}{on}{C3}\,ms, w~MariaDB z~\val{mariadb}{10000000}{off}{C3}
do \val{mariadb}{10000000}{on}{C3}\,ms, a~w~MongoDB z~\val{mongo}{10000000}{off}{C3}
do \val{mongo}{10000000}{on}{C3}\,ms. Każdy nowy wiersz trzeba wstawić także
do dwóch indeksów tabeli \texttt{bilety}. Pojedyncze zapisy po kluczu
(C1, U1, U2) zmieniają się w~granicach zmienności pomiaru. Indeksy
zwiększają też rozmiar bazy (tab.~\ref{tab:storage}) i~wymagają czasu na
budowę (tab.~\ref{tab:load}).

\subsection{Analiza planów zapytań}

Tabele~\ref{tab:plans-rel} i~\ref{tab:plans-nosql} zestawiają automatycznie
odczytane z~planów ścieżki dostępu (dla 13 scenariuszy odczytujących dane),
liczbę przeczytanych wierszy lub dokumentów i~czas wykonania raportowany przez
silnik. Komórki z~pełnym skanem są wyróżnione. Bez indeksów pełny skan
występuje w~\fullscans{postgres}{10000000}{off} z~\PlanScenarioCount{} planów
PostgreSQL, \fullscans{mariadb}{10000000}{off} MariaDB i~\fullscans{mongo}{10000000}{off}
MongoDB. Po dodaniu indeksów zostaje ich odpowiednio \fullscans{postgres}{10000000}{on},
\fullscans{mariadb}{10000000}{on} i~\fullscans{mongo}{10000000}{on}.

\begin{table}[htbp]
\centering
\scriptsize
\caption{Plany zapytań, bazy relacyjne, 10 mln biletów: ścieżki dostępu,
przeczytane wiersze, czas silnika. Tło: pełny skan tabeli.}
\label{tab:plans-rel}
\resizebox{\linewidth}{!}{\tbinput{generated/tables/plans_postgres_mariadb_10000000.tex}}
\end{table}

\begin{table}[htbp]
\centering
\scriptsize
\caption{Plany zapytań, MongoDB (\texttt{explain} z~\texttt{executionStats})
i~polecenia Valkey, 10 mln biletów.}
\label{tab:plans-nosql}
\resizebox{\linewidth}{!}{\tbinput{generated/tables/plans_mongo_valkey_10000000.tex}}
\end{table}

Wnioski z~planów:
\begin{itemize}
  \item W~R4 PostgreSQL po dodaniu indeksów zastępuje dwa skany sekwencyjne
  dwoma skanami indeksów (\texttt{rezerwacje\_klient},
  \texttt{bilety\_rezerwacja}). Czas silnika spada z~setek milisekund do
  dziesiątych części milisekundy (rys.~\ref{lst:r4}). MariaDB
  przechodzi z~\texttt{ALL} na \texttt{ref}: ponad 11 mln przeczytanych wierszy
  zmienia się w~kilka.
  \item R6 pokazuje, że indeks nie zawsze pomaga. PostgreSQL wykonuje
  Index Only Scan na \texttt{bilety(wydarzenie\_id, status)}, bez czytania
  tabeli. MariaDB czyta cały indeks (\texttt{index}). MongoDB nadal robi
  \texttt{COLLSCAN}, bo \texttt{\$match} po samym polu \texttt{status} nie
  pasuje do prefiksu indeksu złożonego \texttt{(wydarzenie\_id, status)}.
  Kolejność pól w~indeksie złożonym ma znaczenie.
  \item W~R2 z~indeksem PostgreSQL zostawia pełny skan małej tabeli
  \texttt{sektory} (kilkaset wierszy). Planer uznaje go za tańszy od
  indeksu, co jest poprawną decyzją kosztową, a~nie brakiem indeksu.
  \item W~J1 PostgreSQL używa indeksu GIN na \texttt{metadane} (Bitmap Index
  Scan), MongoDB indeksu na zagnieżdżonych polach \texttt{metadane.gatunek}
  i~\texttt{metadane.tagi}. Dane JSON są indeksowane na równi z~kolumnami.
  \item Valkey nie ma planów. Tabela pokazuje polecenia: bez struktur
  \texttt{ix:*} zostaje \texttt{SCAN} po przestrzeni kluczy, z~nimi
  \texttt{ZRANGE} po zbiorze posortowanym.
\end{itemize}

\begin{figure}[H]
{\VerbatimInput[firstline=9,frame=single,breaklines=true,breakanywhere=true,fontsize=\tiny,label={bez indeksów}]{data/results/plans/10000000/postgres/indexes-off/R4.txt}}
{\VerbatimInput[firstline=9,frame=single,breaklines=true,breakanywhere=true,fontsize=\tiny,label={z indeksami}]{data/results/plans/10000000/postgres/indexes-on/R4.txt}}
\caption{Plan R4 w~PostgreSQL bez indeksów i~z~indeksami, 10 mln biletów.
Surowe pliki z~\texttt{data/results/plans}, wstawione bez zmian.}
\label{lst:r4}
\end{figure}


\subsection{Narzut gRPC i śledzenia}
\label{sec:grpc}

\begin{table}[H]
\centering
\caption{Różnica między czasem wywołania mierzonym w~runnerze a~czasem
w~serwisie bazy, uśredniona po wszystkich scenariuszach i~rozmiarach.}
\label{tab:grpc}
\small
\tbinput{generated/tables/grpc.tex}
\end{table}

\begin{table}[H]
\centering
\caption{Narzut OpenTelemetry: mediana czasu wywołania (runner) dla tych
samych scenariuszy z~próbkowanym i~niepróbkowanym spanem głównym, 500 tys.
biletów, 30 powtórzeń.}
\label{tab:otel}
\small
\tbinput{generated/tables/overhead.tex}
\end{table}

Hop gRPC między runnerem a~serwisem bazy dodaje w~medianie dziesiąte części
milisekundy (tab.~\ref{tab:grpc}). Przy najszybszych operacjach, takich jak
odczyt rezerwacji po kluczu, to od jednej piątej do prawie połowy całego czasu
wywołania. Gdyby czas mierzył klient, różnice między bazami w~operacjach
punktowych byłyby rozmyte przez transport. Dlatego porównujemy czas mierzony
w~serwisie bazy. Średni narzut jest wyższy od mediany, bo obejmuje też
przesłanie dużych odpowiedzi (np. tysięcy wolnych miejsc w~R2).

Pomiar narzutu OpenTelemetry (tab.~\ref{tab:otel}) porównuje te same
scenariusze ze śladem próbkowanym i~niepróbkowanym, w~którym identyfikator
śladu nadal jest propagowany. Różnice mają oba znaki i~mieszczą się w~rozrzucie
pomiaru operacji trwających ułamek milisekundy. Pełne śledzenie z~eksportem
spanów nie zmienia więc wyników w~sposób mierzalny w~tym teście. Pozostaje
pomijalne także dlatego, że import hurtowy i~test współbieżności zapisują
spany tylko dla części operacji.

\subsection{Dane półstrukturalne (JSON)}
\label{sec:json}
%<*json>

Metadane wydarzeń, biletów, klientów, rezerwacji i~płatności są zapisane jako
JSON w~każdej bazie, choć na różne sposoby. PostgreSQL używa \texttt{JSONB}
(binarnie, z~operatorami i~indeksem GIN). MariaDB używa \texttt{JSON}
(tekst z~kontrolą poprawności, kolumna wirtualna do indeksu). W~MongoDB
metadane są natywnymi poddokumentami BSON. Valkey zapisuje JSON jako napis
w~polu hasza, a~częściową aktualizację wykonuje skrypt Lua z~biblioteką
\texttt{cjson}.

Filtr po polach JSON (J1) z~indeksem trwa przy 10 mln biletów
\val{postgres}{10000000}{on}{J1}\,ms w~PostgreSQL, \val{mariadb}{10000000}{on}{J1}\,ms
w~MariaDB, \val{mongo}{10000000}{on}{J1}\,ms w~MongoDB i~\val{valkey}{10000000}{on}{J1}\,ms
w~Valkey. Agregacja po polu \texttt{metadane.kanal} rezerwacji z~tygodnia (J2)
pokazuje koszt wyłuskiwania pola z~każdego dokumentu: PostgreSQL
\val{postgres}{10000000}{on}{J2}\,ms, MariaDB \val{mariadb}{10000000}{on}{J2}\,ms,
MongoDB \val{mongo}{10000000}{on}{J2}\,ms, Valkey \val{valkey}{10000000}{on}{J2}\,ms
(dekodowanie JSON-a po stronie klienta). Dopisanie elementu do tablicy w~JSON
(J3) jest operacją punktową we wszystkich bazach. PostgreSQL i~MariaDB
przepisują wtedy całą wartość kolumny, MongoDB wykonuje \texttt{\$push} na
miejscu, a~Valkey dekoduje i~koduje JSON w~skrypcie Lua.
%</json>


\section{Współbieżna rezerwacja tych samych miejsc}
\label{sec:wspolbieznosc}

\newcommand{\contplot}[2]{%
\begin{subfigure}{0.49\linewidth}
\begin{tikzpicture}
\begin{axis}[tb, height=5.4cm, xmode=log, ymode=log, xlabel={liczba klientów}, ylabel={#2},
  xtick={1,10,100,1000}, legend to name=contlegend]
  \addplot[pg] table[x=clients, y=postgres] {#1}; \addlegendentry{PostgreSQL}
  \addplot[maria] table[x=clients, y=mariadb] {#1}; \addlegendentry{MariaDB}
  \addplot[mongo] table[x=clients, y=mongo] {#1}; \addlegendentry{MongoDB}
  \addplot[valkey] table[x=clients, y=valkey] {#1}; \addlegendentry{Valkey}
\end{axis}
\end{tikzpicture}
\end{subfigure}}

\begin{figure}[H]
\centering
\contplot{generated/plots/contention_500000.dat}{rezerwacje/s, 0,5 mln}\hfill
\contplot{generated/plots/contention_10000000.dat}{rezerwacje/s, 10 mln}

\contplot{generated/plots/contention_p95_500000.dat}{p95 próby [ms], 0,5 mln}\hfill
\contplot{generated/plots/contention_p95_10000000.dat}{p95 próby [ms], 10 mln}

\ref{contlegend}
\caption{Przepustowość udanych rezerwacji i~95.~percentyl czasu pojedynczej
próby w~funkcji liczby równoległych klientów, domyślne strategie
(\texttt{cond}, \texttt{cond}, \texttt{atomic}, \texttt{lua}). Obie osie
logarytmiczne.}
\label{fig:cont}
\end{figure}

\begin{table}[H]
\centering
\caption{Przepustowość [rezerwacje/s] wszystkich strategii, 1 mln biletów,
500 spornych miejsc, średnia z~3 rund. $^{*}$ strategia domyślna, $^{!}$ runda
niespójna (nie wystąpiła).}
\label{tab:strat}
\small
\tbinput{generated/tables/strategies_1000000.tex}
\end{table}

\begin{table}[H]
\centering
\caption{Liczba konfliktów na rundę i~p95 czasu próby, strategie domyślne,
1 mln biletów.}
\label{tab:conf}
\small
\tbinput{generated/tables/contention_1000000.tex}
\end{table}

We wszystkich \ContentionRuns{} rundach testu (4 bazy, 13 strategii, 4 poziomy
współbieżności, 3 rozmiary, 3 powtórzenia) liczba rezerwacji w~bazie zgadzała
się z~liczbą zajętych miejsc. Liczba niespójnych rund wynosi
\ContentionInconsistent. Żadna strategia nie dopuściła do podwójnej sprzedaży,
więc wyniki różnią się wydajnością, a~nie poprawnością.

\paragraph{Valkey.} Skrypt Lua przy 10 mln biletów obsługuje
\cval{valkey}{10000000}{lua}{1}\,rezerwacji/s z~jednym klientem
i~\cval{valkey}{10000000}{lua}{1000}\,rez./s z~tysiącem klientów. Przepustowość
rośnie z~liczbą klientów, bo jednowątkowy serwer wykonuje skrypty jeden po
drugim bez blokad, a~sterownik valkey-go łączy równoległe wywołania w~potoki.
95.~percentyl czasu próby przy 1000 klientach to \cpn{valkey}{10000000}{lua}{1000}\,ms.
Strategia \texttt{watch} (optymistyczna, na całym zbiorze wolnych miejsc)
załamuje się: \ckn{valkey}{10000000}{watch}{1000} przerwanych transakcji na
rundę i~\cval{valkey}{10000000}{watch}{1000}\,rez./s. Ten sam silnik daje więc
wyniki różniące się o~kilka rzędów wielkości w~zależności od sposobu blokowania.

\paragraph{PostgreSQL.} Warunkowy \texttt{UPDATE} (\texttt{cond}) osiąga
maksimum przy 100 klientach (\cval{postgres}{10000000}{cond}{100}\,rez./s),
a~przy 1000 klientach spada do \cval{postgres}{10000000}{cond}{1000}\,rez./s
z~p95 równym \cpn{postgres}{10000000}{cond}{1000}\,ms. Tysiąc procesów
serwera na 10 rdzeniach czeka w~kolejce na blokady tych samych 500 wierszy.
\texttt{SELECT ... FOR UPDATE} przy 1000 klientach daje więcej
(\cval{postgres}{10000000}{forupdate}{1000}\,rez./s), choć przy 100 klientach
kolejność jest odwrotna. Przy dużej współbieżności sposób oczekiwania na
blokadę wiersza ma więc znaczenie.
\texttt{serializable} traci najwięcej na ponowieniach
(\ckn{postgres}{10000000}{serializable}{1000} konfliktów na rundę).

\paragraph{MariaDB.} Do 100 klientów MariaDB dorównuje PostgreSQL, a~z~\texttt{nowait}
osiąga nawet \cval{mariadb}{10000000}{nowait}{100}\,rez./s. Przy 1000 klientach
strategie czekające na blokadę spadają do kilkudziesięciu rezerwacji na
sekundę (\texttt{cond}: \cval{mariadb}{10000000}{cond}{1000}\,rez./s,
p95 \cpn{mariadb}{10000000}{cond}{1000}\,ms). Rundy trwają wtedy dziesiątki
sekund. Obserwujemy kolejki tysiąca wątków czekających na blokady rekordów
InnoDB w~trybie REPEATABLE READ. \texttt{nowait} nie czeka, tylko od razu
zgłasza konflikt, i~zachowuje \cval{mariadb}{10000000}{nowait}{1000}\,rez./s.
Przy wysokiej współbieżności lepiej szybko odrzucić próbę niż kolejkować.

\paragraph{MongoDB.} Z~jednym klientem MongoDB jest najwolniejszy
(\cval{mongo}{10000000}{atomic}{1}\,rez./s), bo każdy zapis czeka na dziennik.
Przepustowość rośnie z~liczbą klientów, bo zapisy wielu klientów dzielą
jeden zapis dziennika (\emph{group commit}). Przy 1000 klientach osiąga
\cval{mongo}{10000000}{atomic}{1000}\,rez./s, więcej niż PostgreSQL
i~MariaDB w~strategiach domyślnych. Transakcje wielodokumentowe (\texttt{txn})
są wolniejsze od pojedynczej atomowej aktualizacji dokumentu.

\paragraph{Rozmiar zbioru.} Wykresy dla 500 tys. i~10 mln biletów mają ten
sam kształt. Test dotyczy 500 gorących wierszy, więc decyduje o~nim sposób
blokowania, a~nie rozmiar tabeli.


\section{Analiza bezpieczeństwa: TLS i ograniczone role}
\label{sec:bezpieczenstwo}

%<*sec-metoda>
Rozdział mierzy wpływ dwóch mechanizmów bezpieczeństwa na operacje CRUD
i~na rezerwację miejsc przy współbieżności. Pierwszy to szyfrowanie połączenia
między aplikacją a~bazą (TLS), drugi to praca aplikacji na koncie
z~minimalnymi uprawnieniami zamiast na koncie administratora. Pomiar
wykonaliśmy na osobnym stosie (nakładka \texttt{docker-compose.security.yml}),
na zbiorze 1 mln biletów, z~indeksami.

\begin{table}[H]
\centering
\small
\caption{Konfiguracja bezpieczeństwa w~eksperymencie.}
\label{tab:sec-konfig}
\begin{tabular}{>{\raggedright\arraybackslash}p{1.9cm}>{\raggedright\arraybackslash}p{6.9cm}>{\raggedright\arraybackslash}p{5.9cm}}
\toprule
SZBD & konto aplikacyjne (profil \texttt{role}) & TLS (profil \texttt{tls}) \\
\midrule
PostgreSQL & rola \texttt{tb\_app}: \texttt{SELECT}, \texttt{INSERT}, \texttt{UPDATE}, \texttt{DELETE} przez \texttt{ALTER DEFAULT PRIVILEGES}; brak DDL i~własności tabel & \texttt{ssl=on}, klient \texttt{sslmode=verify-full} z~własnym CA \\
MariaDB & użytkownik \texttt{tb\_app}: te same cztery uprawnienia na \texttt{ticketbench.*} & \texttt{--ssl-cert}, klient weryfikuje CA (\texttt{tls=tb}) \\
MongoDB & uwierzytelnianie (SCRAM, plik klucza replica setu), rola z~akcjami \texttt{find}, \texttt{insert}, \texttt{update}, \texttt{remove}; bez \texttt{readWrite}, które pozwala usuwać kolekcje i~indeksy & \texttt{tlsMode=allowTLS}, klient \texttt{tls=true} z~CA \\
Valkey & użytkownik ACL \texttt{tbapp}: tylko prefiksy kluczy modelu (\texttt{w:*}, \texttt{r:*}, \ldots), bez kategorii \texttt{@admin} i~\texttt{@dangerous} (\texttt{FLUSHDB}, \texttt{CONFIG}, \texttt{KEYS}) & osobny port TLS 6380 \\
\bottomrule
\end{tabular}
\end{table}

Dane ładuje i~indeksy tworzy administrator, tak jak przy wdrożeniu
produkcyjnym. Następnie te same scenariusze i~test współbieżności (1 i~100
klientów, 500 spornych miejsc) przechodzą przez cztery profile połączenia:
\texttt{base} (administrator, bez TLS), \texttt{role}, \texttt{tls}
i~\texttt{tls\_role}. Przed pomiarem wykonujemy nieliczony przebieg
rozgrzewkowy wszystkich scenariuszy. Profile mierzymy w~kolejności
zrównoważonej (\texttt{base}, \texttt{role}, \texttt{tls}, \texttt{tls\_role},
a~potem odwrotnie), więc każdy profil ma dwa przebiegi po 5 powtórzeń
w~symetrycznych pozycjach i~liniowy dryf (nagrzewanie pamięci podręcznej,
przyrost danych) się znosi. Porównujemy mediany czasu po stronie serwisu,
który mierzy wywołanie sterownika, więc szyfrowanie i~deszyfrowanie TLS są
wliczone. Narzut grupy operacji to średnia geometryczna ilorazów
$t_{\mathrm{profil}}/t_{\mathrm{base}}$ po scenariuszach tej grupy.

Przed pomiarem sprawdziliśmy, że ograniczenia działają. Konto aplikacyjne
dostaje odmowę przy \texttt{DROP TABLE} w~PostgreSQL i~MariaDB, przy usunięciu
kolekcji w~MongoDB oraz przy \texttt{FLUSHDB} i~dostępie do klucza spoza
prefiksów w~Valkey (\texttt{NOPERM}). Widoki \texttt{pg\_stat\_ssl}, licznik
\texttt{Ssl\_accepts} i~lista klientów Valkey potwierdzają, że profile
\texttt{tls} i~\texttt{tls\_role} faktycznie szyfrują połączenia.
%</sec-metoda>

\begin{table}[H]
\centering
\small
\caption{Zmiana mediany czasu operacji względem profilu \texttt{base}
[\%], 1 mln biletów, z~indeksami. Wartość dodatnia oznacza spowolnienie.}
\label{tab:sec}
\tbinput{generated/tables/security.tex}
\end{table}

\begin{table}[H]
\centering
\small
\caption{Przepustowość rezerwacji [rez./s] w~teście współbieżności przy
1 i~100 klientach, strategie domyślne.}
\label{tab:sec-cont}
\resizebox{\linewidth}{!}{\tbinput{generated/tables/security_contention.tex}}
\end{table}

%<*sec-wyniki>
\paragraph{Punkt odniesienia.} Profil \texttt{base} mierzymy dwa razy:
na początku i~na końcu sekwencji. Różnica między tymi dwoma przebiegami tej
samej konfiguracji (kolumna ,,base A/B'' w~tab.~\ref{tab:sec}) wynosi
\secnoise{postgres} w~PostgreSQL, \secnoise{mariadb} w~MariaDB,
\secnoise{mongo} w~MongoDB i~\secnoise{valkey} w~Valkey. Dla operacji
trwających ułamek milisekundy to naturalna zmienność środowiska. Narzut
profilu, który się w~niej mieści, jest w~tym teście niemierzalny.

\paragraph{Ograniczone role nie kosztują.} Zmiana łącznego czasu po
przejściu na konto aplikacyjne wynosi \secpct{postgres}{role}{all}
w~PostgreSQL, \secpct{mariadb}{role}{all} w~MariaDB,
\secpct{mongo}{role}{all} w~MongoDB i~\secpct{valkey}{role}{all} w~Valkey.
Znaki są różne i~wszystkie wartości mieszczą się w~dryfie ,,base A/B''.
Sprawdzenie uprawnień to porównanie z~listą uprawnień roli. W~bazach
relacyjnych odbywa się raz przy rozpoczęciu wykonania instrukcji,
w~MongoDB raz na operację z~pamięci podręcznej sesji, a~w~Valkey raz na
polecenie (dopasowanie prefiksu klucza). Wobec kosztu dostępu do danych jest
to niezauważalne. Zasada najmniejszych uprawnień nie ma więc w~tym systemie
ceny wydajnościowej.

\paragraph{TLS przy pojedynczych operacjach.} Szyfrowanie połączenia zmienia
łączny czas o~\secpct{postgres}{tls}{all} (PostgreSQL),
\secpct{mariadb}{tls}{all} (MariaDB), \secpct{mongo}{tls}{all} (MongoDB)
i~\secpct{valkey}{tls}{all} (Valkey), również w~granicach dryfu. Pule
utrzymują połączenia otwarte, więc uzgadnianie sesji TLS (najdroższa część)
odbywa się raz na połączenie, a~nie na operację. Szyfrowanie symetryczne
kilkuset bajtów zapytania i~odpowiedzi kosztuje mikrosekundy. Profil
\texttt{tls\_role}, łączący oba mechanizmy, daje podobny obraz.

\paragraph{TLS przy wysokiej współbieżności.} Wyraźny efekt widać w~teście
rezerwacji (tab.~\ref{tab:sec-cont}). Valkey przy 100 klientach obsługuje
bez TLS \seccval{valkey}{base}{100}\,rez./s, a~z~TLS
\seccval{valkey}{tls}{100}\,rez./s. Valkey w~domyślnej konfiguracji (jeden
wątek I/O) sam szyfruje i~deszyfruje ruch wszystkich połączeń w~wątku, który
wykonuje polecenia, więc przy dziesiątkach tysięcy krótkich poleceń na
sekundę koszt TLS konkuruje bezpośrednio z~obsługą danych. Przy jednym
kliencie różnica jest mała (\seccval{valkey}{base}{1} wobec
\seccval{valkey}{tls}{1}\,rez./s), bo serwer nie jest wtedy wąskim gardłem.
W~bazach dyskowych zmiany przepustowości przy 100 klientach nie układają się
w~stały wzór: każda transakcja czeka na zapis dziennika, a~nie na
szyfrowanie.

\paragraph{Wniosek.} Konto o~minimalnych uprawnieniach i~TLS można włączyć
w~PostgreSQL, MariaDB i~MongoDB bez mierzalnej straty wydajności operacji
CRUD. W~Valkey TLS obniża przepustowość przy dużym ruchu. Trzeba to
uwzględnić przy planowaniu (wątki I/O, więcej instancji albo szyfrowanie na
poziomie sieci, np. w~prywatnej sieci klastra).
%</sec-wyniki>


\section{Weryfikacja hipotez}
\label{sec:hipotezy}

\subsection{H1: przewaga Valkey i MongoDB rośnie z liczbą klientów}

Dla każdego rozmiaru i~każdej pary (baza nierelacyjna, baza relacyjna)
liczymy przewagę jako iloraz przepustowości przy strategiach domyślnych, dla
1, 10, 100 i~1000 klientów. Przewaga ,,rośnie'', jeśli współczynnik korelacji
rang Spearmana między liczbą klientów a~przewagą jest dodatni i~przewaga przy
1000 klientów jest większa niż przy jednym.

\begin{table}[H]
\centering
\caption{H1: przewaga przepustowości (iloraz) przy 1 i~przy 1000 klientach.}
\label{tab:h1}
\small
\tbinput{generated/tables/h1.tex}
\end{table}

Przewaga rośnie w~\HOneGrows{} z~\HOneCases{} przypadków. Wynik
automatycznej oceny: hipoteza H1 \textbf{\HOneVerdict}.

Iloraz przepustowości rośnie z~liczbą klientów dla wszystkich par i~rozmiarów.
Współczynnik Spearmana wynosi wszędzie 1 (zależność ściśle monotoniczna). Dla
10 mln biletów przewaga Valkey nad PostgreSQL rośnie
z~\cratio{valkey}{postgres}{10000000}{lua}{cond}{1}$\times$ przy jednym kliencie
do \cratio{valkey}{postgres}{10000000}{lua}{cond}{1000}$\times$ przy tysiącu.
Nad MariaDB rośnie z~\cratio{valkey}{mariadb}{10000000}{lua}{cond}{1}$\times$
do \cratio{valkey}{mariadb}{10000000}{lua}{cond}{1000}$\times$.

Wynik dla MongoDB trzeba czytać ostrożniej. Iloraz MongoDB do PostgreSQL rośnie
z~\cratio{mongo}{postgres}{10000000}{atomic}{cond}{1} do
\cratio{mongo}{postgres}{10000000}{atomic}{cond}{1000}: przy małej
współbieżności MongoDB jest kilka razy wolniejszy, a~przy 1000 klientach
dorównuje PostgreSQL. Względna pozycja MongoDB rośnie, ale realną przewagę
ma on tylko nad MariaDB. Hipotezę H1 uznajemy za
\textbf{potwierdzoną dla Valkey}. Dla MongoDB potwierdza się teza
o~rosnącym stosunku wydajności, a~nie o~przewadze w~każdym punkcie.

Dwa czynniki wzmacniają wynik Valkey i~trzeba je brać pod uwagę: brak zapisu
na dysk oraz łączenie poleceń w~potoki przez sterownik. Strategia \texttt{watch}
pokazuje też, że ten sam silnik z~nieodpowiednią metodą synchronizacji
przegrywa ze wszystkimi bazami relacyjnymi. H1 jest więc prawdziwa dla
właściwie zaprojektowanej operacji atomowej, a~nie dla silnika
niezależnie od sposobu użycia.


\subsection{H2: różnice między silnikami rosną z liczbą rekordów}

Dla każdego scenariusza i~wariantu indeksów liczymy iloraz czasu
najwolniejszej i~najszybszej bazy. Różnica ,,rośnie'', jeśli ten iloraz przy
10 mln biletów jest większy niż przy 500 tys.

\begin{table}[H]
\centering
\caption{H2: iloraz czasu najwolniejszej i~najszybszej bazy.}
\label{tab:h2}
\footnotesize
\IfFileExists{generated/tables/h2.tex}{\begin{minipage}{\linewidth}\centering% generated by tbench report
\begin{tabular}{llrrrr}
\toprule
Scen. & indeksy & 500 tys. & 1 mln & 10 mln & rośnie? \\
\midrule
C1 & bez & \num{4.89} & \num{4.42} & \num{9.17} & tak \\
C1 & z & \num{5.57} & \num{13.0} & \num{12.5} & tak \\
C2 & bez & \num{7.38} & \num{3.13} & \num{7.09} & nie \\
C2 & z & \num{6.17} & \num{6.70} & \num{6.84} & tak \\
C3 & bez & \num{5.55} & \num{4.59} & \num{6.01} & tak \\
C3 & z & \num{8.86} & \num{6.81} & \num{9.00} & tak \\
C4 & bez & \num{2.98} & \num{3.83} & \num{4.17} & tak \\
C4 & z & \num{4.58} & \num{4.16} & \num{4.43} & nie \\
C5 & bez & \num{3.06} & \num{1.57} & \num{1.58} & nie \\
C5 & z & \num{3.70} & \num{2.37} & \num{3.43} & nie \\
C6 & bez & \num{8.92} & \num{5.43} & \num{8.02} & nie \\
C6 & z & \num{13.3} & \num{9.17} & \num{12.0} & nie \\
R1 & bez & \num{447.8} & \num{776.1} & \num{24351} & tak \\
R1 & z & \num{1.45} & \num{1.79} & \num{12.0} & tak \\
R2 & bez & \num{81.5} & \num{160.0} & \num{1954} & tak \\
R2 & z & \num{1.60} & \num{2.20} & \num{1.82} & tak \\
R3 & bez & \num{3.56} & \num{4.51} & \num{15.6} & tak \\
R3 & z & \num{2.52} & \num{2.20} & \num{2.60} & tak \\
R4 & bez & \num{137.6} & \num{73.4} & \num{83.4} & nie \\
R4 & z & \num{1.92} & \num{2.08} & \num{1.60} & nie \\
R5 & bez & \num{511.8} & \num{551.6} & \num{1036} & tak \\
R5 & z & \num{80.2} & \num{80.9} & \num{258.0} & tak \\
R6 & bez & \num{826.0} & \num{1434} & \num{15256} & tak \\
R6 & z & \num{923.9} & \num{1784} & \num{17619} & tak \\
U1 & bez & \num{7.37} & \num{10.7} & \num{4.82} & nie \\
U1 & z & \num{5.37} & \num{5.25} & \num{12.4} & tak \\
U2 & bez & \num{7.21} & \num{13.1} & \num{11.2} & tak \\
U2 & z & \num{8.02} & \num{6.20} & \num{5.33} & nie \\
U3 & bez & \num{67.3} & \num{119.6} & \num{1069} & tak \\
U3 & z & \num{3.19} & \num{2.68} & \num{1.77} & nie \\
U4 & bez & \num{5.17} & \num{4.60} & \num{5.16} & nie \\
U4 & z & \num{3.45} & \num{3.45} & \num{4.65} & tak \\
U5 & bez & \num{3.31} & \num{3.98} & \num{5.49} & tak \\
U5 & z & \num{3.48} & \num{7.40} & \num{5.18} & tak \\
U6 & bez & \num{16.7} & \num{20.1} & \num{20.4} & tak \\
U6 & z & \num{3.04} & \num{2.46} & \num{2.86} & nie \\
D1 & bez & \num{858.5} & \num{239.6} & \num{16378} & tak \\
D1 & z & \num{7.82} & \num{7.09} & \num{9.76} & tak \\
D2 & bez & \num{771.7} & \num{416.6} & \num{16771} & tak \\
D2 & z & \num{8.11} & \num{6.46} & \num{9.42} & tak \\
D3 & bez & \num{240.8} & \num{264.0} & \num{358.0} & tak \\
D3 & z & \num{23.0} & \num{17.8} & \num{17.6} & nie \\
D4 & bez & \num{23.3} & \num{18.8} & \num{20.7} & nie \\
D4 & z & \num{5.05} & \num{6.53} & \num{60.3} & tak \\
D5 & bez & \num{1170} & \num{2152} & \num{26468} & tak \\
D5 & z & \num{94.4} & \num{107.8} & \num{110.8} & tak \\
D6 & bez & \num{20.2} & \num{22.3} & \num{17.0} & nie \\
D6 & z & \num{11.0} & \num{11.4} & \num{9.61} & nie \\
J1 & bez & \num{3.76} & \num{4.24} & \num{8.11} & tak \\
J1 & z & \num{5.51} & \num{3.02} & \num{7.00} & tak \\
J2 & bez & \num{235.6} & \num{275.3} & \num{417.1} & tak \\
J2 & z & \num{16.1} & \num{23.2} & \num{28.3} & tak \\
J3 & bez & \num{5.40} & \num{7.78} & \num{8.19} & tak \\
J3 & z & \num{11.0} & \num{10.4} & \num{9.98} & nie \\
\bottomrule
\end{tabular}
\end{minipage}}{brak danych}
\end{table}

Różnica rośnie w~\HTwoGrows{} z~\HTwoCases{} przypadków. Wynik automatycznej
oceny: hipoteza H2 \textbf{\HTwoVerdict}.

Rozrzut między najszybszą a~najwolniejszą bazą rośnie z~rozmiarem w~większości
przypadków, ale nie we wszystkich. Rośnie tam, gdzie co najmniej jedna baza
skanuje dane w~czasie $O(n)$, a~inna ma ścieżkę stałą lub logarytmiczną.
Przykłady to scenariusze bez indeksów R1, R2, R5, U3, D1, D2, D5 i~J2 oraz
R6, gdzie Valkey odczytuje licznik, a~pozostałe bazy agregują wszystkie bilety
(tab.~\ref{tab:h2}). Przy R4, D4 i~D6 bez indeksów rozrzut nie rośnie, bo
wszystkie bazy, także najszybsza, skanują tam pełne dane. Czas najszybszej
rośnie więc co najmniej w~tym samym tempie co czas najwolniejszej (Valkey ze
skanem \texttt{SCAN}).
Na przykład w~R1 bez indeksów MariaDB jest
\ratio{mariadb}{valkey}{500000}{off}{R1}$\times$ wolniejsza od Valkey przy
500 tys. biletów i~\ratio{mariadb}{valkey}{10000000}{off}{R1}$\times$ przy
10 mln.

Rozrzut nie rośnie dla operacji punktowych po kluczu (C1, C4, U1, U2, J3),
których czas w~każdej bazie jest stały. Nie rośnie też dla zapytań z~indeksem
o~małej selektywności, gdzie wszystkie bazy czytają kilka wierszy. H2 jest
\textbf{potwierdzona częściowo}. Różnice rosną z~liczbą rekordów wtedy, gdy
bazy mają różną złożoność ścieżki dostępu, czyli przy braku indeksu, przy
różnicy modeli danych (osadzenie a~złączenie) lub przy agregacji
utrzymywanej w~locie. Dla dobrze zaindeksowanych operacji punktowych rozmiar
zbioru nie ma znaczenia.


\section{Wnioski}
\label{sec:wnioski}

%<*wnioski>
\begin{enumerate}
  \item \textbf{Nie ma jednej najszybszej bazy.} Valkey wygrywa w~operacjach
  po kluczu, w~licznikach i~w~rezerwacji miejsc przy dużej współbieżności.
  Bez przygotowanych przez aplikację struktur dostępu jest jednak najwolniejszy
  ze wszystkich: historia klienta (R4) bez indeksów trwa w~nim
  \val{valkey}{10000000}{off}{R4}\,ms, a~w~PostgreSQL \val{postgres}{10000000}{off}{R4}\,ms. PostgreSQL jest najbardziej
  przewidywalny: bez indeksów najszybszy z~baz dyskowych, z~indeksami
  w~czołówce wszystkich scenariuszy.
  \item \textbf{Indeksy decydują bardziej niż silnik.} Ten sam scenariusz
  w~tej samej bazie przyspiesza po dodaniu indeksu nawet o~kilka rzędów
  wielkości (rys.~\ref{fig:speedup}). Po zaindeksowaniu różnice między
  silnikami w~zapytaniach selektywnych są zwykle mniejsze niż rząd wielkości.
  Duże pozostają w~agregacjach po całej tabeli (R6, J2). Koszt to wolniejsze wstawianie hurtowe
  i~dodatkowe miejsce.
  \item \textbf{Model danych zmienia złożoność operacji.} Osadzenie biletów
  w~dokumencie rezerwacji daje w~MongoDB odczyt i~anulowanie rezerwacji
  w~czasie stałym bez indeksów drugorzędnych. Baza relacyjna potrzebuje do
  tego indeksu na kluczu obcym.
  \item \textbf{Współbieżność to kwestia strategii blokowania.} Przy 1000
  klientów strategia (warunkowy UPDATE, \texttt{FOR UPDATE NOWAIT}, skrypt
  Lua, WATCH) zmienia przepustowość tego samego silnika nawet o~trzy
  rzędy wielkości (Valkey: \texttt{lua} a~\texttt{watch}). Szybkie odrzucanie konfliktów (\texttt{NOWAIT}, operacja
  atomowa) skaluje się lepiej niż kolejkowanie na blokadach.
  \item \textbf{Bezpieczeństwo bez kosztu, z~jednym wyjątkiem.} Konta
  o~minimalnych uprawnieniach i~TLS nie zmieniają czasu operacji CRUD ponad
  zmienność pomiaru. Wyjątkiem jest Valkey przy 100 klientach, gdzie TLS
  obniża przepustowość z~\seccval{valkey}{base}{100} do
  \seccval{valkey}{tls}{100}\,rez./s (rozdz.~\ref{sec:bezpieczenstwo}).
  \item \textbf{H1 potwierdzona dla Valkey, częściowo dla MongoDB, H2
  potwierdzona częściowo} (rozdz.~\ref{sec:hipotezy}).
  \item \textbf{Praktyczna rekomendacja dla systemu biletowego:}
  rezerwacje i~płatności w~PostgreSQL (spójność, złożone zapytania), blokady
  miejsc z~TTL i~liczniki dostępności w~Valkey (operacja atomowa w~skrypcie
  Lua). Katalog wydarzeń może być w~MongoDB, jeśli jego struktura często się
  zmienia.
\end{enumerate}
%</wnioski>

\subsection{Zagrożenia dla trafności wyników}
\label{sec:zagrozenia}
%<*zagrozenia>

\begin{itemize}
  \item Wszystkie bazy działały na jednym komputerze w~maszynie wirtualnej
  Docker Desktop (8\,GB RAM, Apple M5). Wyniki pokazują relacje między
  silnikami, a~nie przepustowość produkcyjną. Generator obciążenia działał na
  tej samej maszynie co bazy.
  \item Poziomy trwałości nie są równe: Valkey bez RDB i~AOF nie czeka na dysk,
  MongoDB czeka na dziennik (\texttt{majority}), bazy relacyjne na
  \texttt{fsync} WAL i~redo. Zapisy w~Valkey są przez to tańsze z~założenia.
  \item Każda baza działała jako pojedynczy węzeł, a~główna macierz bez TLS.
  Wpływ TLS i~ról zmierzyliśmy osobno (rozdz.~\ref{sec:bezpieczenstwo}),
  wpływu replikacji synchronicznej nie badaliśmy.
  \item W~każdym wariancie mierzyliśmy 5 powtórzeń po jednym przebiegu
  rozgrzewkowym. Dla operacji trwających poniżej 1\,ms pojedyncze odchylenia
  (zimna pamięć podręczna, punkt kontrolny) przesuwają średnią, dlatego
  w~danych są też mediany i~p95.
  \item Indeksy w~Valkey są projektem autorów (zbiory \texttt{ix:*}), a~nie
  funkcją silnika. Inny projekt struktur dałby inne wyniki.
  \item Parametry konfiguracyjne baz zostały ustawione tak samo (1\,GB
  pamięci podręcznej), ale nie były strojone pod obciążenie. Strojenie
  (np. \texttt{innodb\_thread\_concurrency}, pooler połączeń dla PostgreSQL)
  zmieniłoby wyniki przy 1000 klientach.
\end{itemize}
%</zagrozenia>

\section{Harmonogram i podział pracy}

Projekt realizowaliśmy w~zespole dwuosobowym według harmonogramu ze wstępnej
koncepcji. Na zajęciach przedstawialiśmy postęp kolejnych etapów.

\begin{table}[H]
\centering
\small
\caption{Etapy projektu i~odpowiedzialność.}
\begin{tabular}{c>{\raggedright\arraybackslash}p{8.2cm}>{\raggedright\arraybackslash}p{5cm}}
\toprule
etap & zakres & odpowiedzialni \\
\midrule
1 & temat, model danych (12 tabel, model dokumentowy i~klucz--wartość), środowisko Docker z~czterema bazami & Łukasz Ryczko, Łukasz Pacholec \\
2 & deterministyczny generator danych, import do wszystkich baz, kontrakty gRPC & Łukasz Ryczko \\
3 & implementacja 24 scenariuszy CRUD i~3 JSON w~czterech serwisach, automatyzacja pomiarów (runner, CLI) & Łukasz Pacholec \\
4 & indeksy, plany zapytań, test współbieżności i~strategie blokowania, OpenTelemetry, CI, Ansible & Łukasz Ryczko, Łukasz Pacholec \\
5 & pomiary, analiza wyników, weryfikacja hipotez, sprawozdanie i~prezentacja & Łukasz Ryczko, Łukasz Pacholec \\
\bottomrule
\end{tabular}
\end{table}

\section{Zgodność z wymaganiami regulaminu}
\label{sec:zgodnosc}

\begin{table}[H]
\centering
\footnotesize
\caption{Wymagania z~§3 regulaminu zajęć i~miejsce ich realizacji.}
\begin{tabular}{>{\raggedright\arraybackslash}p{0.8cm}>{\raggedright\arraybackslash}p{7.2cm}>{\raggedright\arraybackslash}p{7.2cm}}
\toprule
ocena & wymaganie & realizacja \\
\midrule
3.0 & cel i~zakres pracy & rozdz.~1 \\
3.0 & opis SZBD, zalety i~wady & rozdz.~\ref{sec:szbd}, tab.~\ref{tab:zalety} \\
3.0 & awaryjność, bezpieczeństwo, migracje, integracje, skalowalność & tab.~\ref{tab:ops} \\
3.0 & obszary zastosowań biznesowych & rozdz.~\ref{sec:szbd} \\
3.0 & opis zbioru danych, min. 5 tabel & rozdz.~3: 12 tabel, model dokumentowy i~klucz--wartość \\
3.0 & aplikacja: wymagania, technologie, działanie & rozdz.~4 \\
3.0 & CRUD dla małego, średniego i~dużego zbioru & 0,5 mln, 1 mln, 10 mln biletów, tab.~\ref{tab:s500}--\ref{tab:s10m} \\
3.0 & min. 4 SZBD (2 relacyjne, 2 nierelacyjne) & PostgreSQL, MariaDB, MongoDB, Valkey \\
3.0 & min. 12 scenariuszy, min. 3 na operację CRUD & 24 CRUD (6 na operację) + 3 JSON \\
3.0 & średnia z~min. 3 prób & 5 prób + rozgrzewka, średnia, mediana, odchylenie, p95 \\
3.0 & opis i~wizualizacja wyników, sprawozdanie i~prezentacja & rozdz.~5--8, wykresy, \texttt{prezentacja.pdf} \\
\midrule
4.0 & min. dwa modele danych & trzy: relacyjny, dokumentowy, klucz--wartość \\
4.0 & indeksy & 15 indeksów drugorzędnych w~każdej bazie, zbiory \texttt{ix:*} w~Valkey \\
4.0 & min. 10 tabel & 12 tabel \\
4.0 & analiza planów zapytań & EXPLAIN ANALYZE, ANALYZE FORMAT=JSON, explain executionStats; tab.~\ref{tab:plans-rel}, \ref{tab:plans-nosql}, rys.~\ref{lst:r4} \\
4.0 & porównanie przed i~po indeksach & każdy scenariusz w~obu wariantach, rys.~\ref{fig:speedup} \\
4.0 & 24 scenariusze, min. 6 na operację, min. 3 próby & 24 + 3, po 5 prób \\
4.0 & rozszerzona analiza wyników i~wnioski & rozdz.~5--8, zagrożenia trafności \\
4.0 & zbiory 0,5 mln, 1 mln, 10 mln & tak \\
\midrule
5.0 & testy skalowalności (równoległe zapytania, wielu użytkowników) & rozdz.~6: 1--1000 klientów, 13 strategii, kontrola spójności \\
5.0 & dane półstrukturalne & JSONB, JSON, BSON, JSON w~Valkey; scenariusze J1--J3, rozdz.~\ref{sec:json} \\
5.0 & automatyzacja testów i~generowania wyników & \texttt{ticketbench.sh}, Docker Compose, CI, Ansible, \texttt{tbench report} \\
5.0 & sformułowanie i~weryfikacja hipotezy & H1 (rozdz.~\ref{sec:hipotezy}), dodatkowo H2 \\
\midrule
§3 & harmonogram i~podział zadań & rozdz.~9 \\
\bottomrule
\end{tabular}
\end{table}

Z~czterech elementów zaawansowanych na ocenę 5.0 (wymagane dwa) projekt
realizuje trzy: testy skalowalności, dane półstrukturalne i~automatyzację.
Analiza bezpieczeństwa jest opisana teoretycznie (tab.~\ref{tab:ops}), bez
pomiaru jej wpływu na operacje CRUD.


\appendix
\section{Repozytoria i uruchomienie}

\begin{table}[H]
\centering
\small
\begin{tabular}{lp{10.5cm}}
\toprule
repozytorium & zawartość \\
\midrule
\texttt{ryczek02/ticketbench} & repozytorium spinające: README, \texttt{ticketbench.sh}, runbook (\texttt{docs/runbook.html}), to sprawozdanie z~surowymi wynikami, submoduły pozostałych repozytoriów \\
\texttt{ticketbench-contracts} & pliki \texttt{.proto}, wygenerowany kod, model, generator danych, scenariusze, telemetria \\
\texttt{ticketbench-services} & \texttt{bench-runner}, \texttt{store-postgres}, \texttt{store-mariadb}, \texttt{store-mongo}, \texttt{store-valkey}, testy zgodności \\
\texttt{ticketbench-cli} & \texttt{tbench}: plany, zapis wyników, generator tabel i~wykresów \\
\texttt{ticketbench-infra} & Docker Compose, OTel/Jaeger/Prometheus/Grafana, macierze, Ansible, \texttt{bench.sh} \\
\bottomrule
\end{tabular}
\end{table}

Odtworzenie wszystkich wyników:
\begin{Verbatim}[frame=single,fontsize=\small]
git clone --recursive https://github.com/ryczek02/ticketbench
cd ticketbench
./ticketbench.sh all                # narzędzia, pobranie z GitHuba, stos, test dymny
./ticketbench.sh bench              # pełna macierz
./ticketbench.sh report             # tabele, wykresy, PDF
\end{Verbatim}

Wybrane polecenia CLI:
\begin{Verbatim}[frame=single,fontsize=\small]
tbench ping
tbench run --stores postgres,valkey --tickets 1M --scenarios R4,U6 --reps 10
tbench contend --clients 1,10,100,1000 --strategies all
tbench report results/<run> --out ../report/generated
\end{Verbatim}

Wszystkie liczby w~tabelach i~na wykresach pochodzą z~plików w~katalogu
\texttt{generated/}, które \texttt{tbench report} tworzy z~surowych wyników
w~\texttt{data/results}. CI repozytorium \texttt{ticketbench} generuje te pliki
ponownie i~sprawdza, czy są identyczne z~zatwierdzonymi. Szczegółowa instrukcja
uruchomienia jest w~runbooku \texttt{docs/runbook.html}.


\begin{thebibliography}{9}
\bibitem{pg} PostgreSQL Global Development Group, \emph{PostgreSQL 17 Documentation}, rozdz. 13 (Concurrency Control) i 14 (Performance Tips), \url{https://www.postgresql.org/docs/17/}.
\bibitem{maria} MariaDB Foundation, \emph{MariaDB Server Documentation: InnoDB, Optimization and Indexes, JSON Functions}, \url{https://mariadb.com/kb/}.
\bibitem{mongo} MongoDB Inc., \emph{MongoDB 8.0 Manual: Data Modeling, Transactions, Explain Results}, \url{https://www.mongodb.com/docs/manual/}.
\bibitem{valkey} Valkey contributors, \emph{Valkey Documentation: Commands, Scripting with Lua, Transactions}, \url{https://valkey.io/docs/}.
\bibitem{otel} OpenTelemetry Authors, \emph{OpenTelemetry Specification: Trace Context Propagation}, \url{https://opentelemetry.io/docs/specs/}.
\bibitem{w3c} W3C, \emph{Trace Context, W3C Recommendation}, 2021, \url{https://www.w3.org/TR/trace-context/}.
\bibitem{repo} Ł.~Ryczko, Ł.~Pacholec, \emph{ticketbench}, kod źródłowy, \url{https://github.com/ryczek02/ticketbench}.
\end{thebibliography}

\end{document}
